% !TeX program = lualatex
\documentclass[11pt,a4paper]{article}
\usepackage{szzmaterial}
\usepackage{tikz}
\usetikzlibrary{positioning,arrows.meta,calc,fit,backgrounds,shapes.geometric,shapes.misc,matrix,decorations.pathreplacing,patterns,intersections,angles,quotes}
\usepackage{pgfplots}
\pgfplotsset{compat=1.18}
\usepgfplotslibrary{fillbetween}
\setlist[enumerate]{nosep,leftmargin=2.1em}
\setlist[itemize]{nosep,leftmargin=1.8em}

\DeclareMathOperator*{\argmax}{argmax}
\DeclareMathOperator*{\argmin}{argmin}
\DeclareMathOperator{\Var}{Var}
\DeclareMathOperator{\Cov}{Cov}
\DeclareMathOperator{\sign}{sign}
\DeclareMathOperator{\diag}{diag}
\DeclareMathOperator{\tr}{tr}

% --- Box pro algoritmy (pseudokód). Lokální doplněk preambule okruhu (jako v S1). ---
\newtcolorbox{algo}[1][]{enhanced,breakable,colback=teal!4,colframe=teal!55!black,
  boxrule=0.6pt,arc=2pt,left=6pt,right=8pt,top=3pt,bottom=4pt,
  fonttitle=\bfseries,coltitle=white,colbacktitle=teal!55!black,
  title={\faCogs\enspace Algoritmus \textemdash\ #1}}
% Komentář v pseudokódu
\newcommand{\cmt}[1]{\hfill{\footnotesize\itshape\textcolor{black!55}{$\triangleright$ #1}}}

% --- Společná barevná konvence v obrázcích S3 ---
% Třída / shluk 1 = modrá, 2 = červená, 3 = zelená. Centroid = hvězda. Podpůrný vektor = zakroužkovaný.
\tikzset{
  osa/.style={->,>=Stealth,black!55},
  cls1/.style={circle,draw=blue!60!black,fill=blue!35,thick,inner sep=0pt,minimum size=4.5pt},
  cls2/.style={rectangle,draw=red!65!black,fill=red!45,thick,inner sep=0pt,minimum size=4.5pt},
  cls3/.style={diamond,draw=green!45!black,fill=green!45,thick,inner sep=0pt,minimum size=5pt},
  centroid/.style={star,star points=5,draw=black,thick,fill=yellow!85!orange,inner sep=0pt,minimum size=9pt},
  podpora/.style={circle,draw=black,thick,inner sep=0pt,minimum size=8pt},
  hranice/.style={very thick,black!75},
  margin/.style={dashed,black!55,thick},
  treenode/.style={draw,thick,rounded corners=2pt,fill=blue!8,font=\small,minimum height=6mm,inner sep=3pt},
  treeleaf/.style={draw,thick,fill=green!12,font=\small,minimum height=6mm,inner sep=3pt},
  vetev/.style={->,>=Stealth,thick},
  popis/.style={font=\scriptsize,fill=white,inner sep=1pt},
}

\szztitul{Strojové učení}{S3}{NAIL121 (Seminář dobývání znalostí), NPFL054 (Úvod do SU v R), NPFL129 (Úvod do SU v Pythonu), NPGR035 (SU v počítačovém vidění)}

\begin{document}
\maketitle
\tableofcontents
\bigskip

% ============================================================
\section*{Úvod}
\addcontentsline{toc}{section}{Úvod}

Tento materiál pokrývá okruh \textbf{S3 -- Strojové učení}, druhý ze dvou okruhů specializace UI-SU. Výklad stojí především na kurzu \textbf{NPFL129 \emph{Úvod do strojového učení v~Pythonu}} (Milan Straka, výuka Jindřich Libovický; ÚFAL MFF, zdrojový repozitář \url{https://github.com/ufal/npfl129}). Téma \emph{metody podpůrných vektorů} (SVM) čerpá z~ročníku, kdy bylo v~osnově NPFL129 (jádrové funkce a~SVM, přednášky 6--7), \emph{statistické testy} (t-test, $\chi^2$) navazují na okruh \textbf{M5} (Pravděpodobnost a~statistika). Strojové učení je probíráno i~v~úvodním kurzu UI (okruh \textbf{S1}, sekce \emph{Strojové učení}) na úrovni „základní principy``; \emph{tento} okruh jde do hloubky a~na S1 jen navazuje (kde stačí přehled, odkazujeme tam).

\medskip
\noindent\textbf{Role důkazů.} Požadavky okruhu jmenují pojmy, modely, algoritmy a~metody evaluace, ne formální důkazy. Těžiště je v~definici pojmu, v~popisu a~aplikaci algoritmu na konkrétních datech (matice záměn a~metriky, stavba rozhodovacího stromu podle informačního zisku, iterace $k$-means, krok gradientního sestupu, odhad koeficientů regrese, statistický test) a~ve zdůvodnění. Výzvy „odvoďte`` a~„ukažte`` se v~zařazených otázkách týkají \emph{výpočtu}: derivace ztrátové funkce logistické regrese (sekce~\ref{sec:logisticka}) a~explicitního řešení nejmenších čtverců (sekce~\ref{sec:linreg}). Věty a~vlastnosti (např.\ konvergence $k$-means, optimum nejmenších čtverců) proto uvádíme se zněním, intuicí a~krátkým \emph{odvozením} tam, kde pomáhá pochopení.

Do výkladu je zařazeno všech $35$ reálných otázek SZZ k~okruhu z~let 2022--2026 jako řešené příklady; jejich originální zadání (a~grafy oříznuté z~PDF) jsou i~v~samostatném dokumentu \emph{Příklady otázek}. Použité zdroje jsou uvedeny na konci.

\medskip
\noindent\textbf{Značení.} \emph{Vstup} (vektor příznaků) je $\mathbf{x}\in\mathbb{R}^D$, jednotlivé příznaky $x_1,\dots,x_D$. \emph{Trénovací množina} má $N$ příkladů; \emph{matice dat} $\mathbf{X}\in\mathbb{R}^{N\times D}$ má příklad $\mathbf{x}_i$ na $i$-tém řádku. \emph{Cíl} (target) je $t$ -- reálné číslo u~regrese ($t\in\mathbb{R}$), třída u~klasifikace ($t\in\{0,1,\dots,K-1\}$, případně $t\in\{0,1\}$ binárně); vektor cílů $\mathbf{t}\in\mathbb{R}^N$. \emph{Predikci} modelu značíme $y(\mathbf{x})$ nebo $\hat{t}$. \emph{Váhy} lineárního modelu jsou $\mathbf{w}\in\mathbb{R}^D$, \emph{bias} $b$ (ve statistickém značení regrese koeficienty $\beta_0,\beta_1,\dots$); skalární součin $\mathbf{w}^{\!\top}\mathbf{x}=\sum_i w_i x_i$. \emph{Rychlost učení} (krok) je $\alpha$, \emph{regularizační síla} $\lambda$ (nebo $C$ u~SVM). $\sigma(z)=1/(1+e^{-z})$ je \emph{logistická} (sigmoidní) funkce. Logaritmus $\log$ bez základu v~entropii je $\log_2$ (entropie v~bitech). $\mathbb{E}$ je střední hodnota, $\Var$ rozptyl, $\|\cdot\|$ eukleidovská ($L^2$) norma.

% ============================================================
\section{Učení s učitelem a evaluace modelů}\label{sec:evaluace}

Tato sekce zavádí společný rámec celého okruhu: co je \emph{učení s~učitelem}, jak rozlišujeme \emph{klasifikaci} a~\emph{regresi}, a~především \textbf{jak model ohodnotíme} -- testovacími daty, křížovou validací a~principem maximum likelihood estimation, a~jakými \emph{metrikami} měříme jeho kvalitu. Konkrétní modely (regrese, stromy, SVM, $\dots$) přijdou v~dalších sekcích; evaluace je společná všem.

\begin{poznamka}[Vztah k~S1]
Okruh S1 (sekce \emph{Strojové učení}) podává učení s~učitelem na úrovni „základní principy``: druhy učení, klasifikace vs.\ regrese, Ockhamova břitva. Zde to nezdvojujeme -- jen krátce zrekapitulujeme a~rovnou jdeme do hloubky, kterou S1 výslovně přenechává okruhu S3 (evaluační metriky, křížová validace).
\end{poznamka}

\subsection{Úlohy učení s učitelem}

\begin{definice}[Učení s~učitelem, klasifikace, regrese]
Při \textbf{učení s~učitelem} máme \emph{trénovací množinu} $N$ příkladů $(\mathbf{x}_i,t_i)$, kde $\mathbf{x}_i\in\mathbb{R}^D$ je vstup (vektor příznaků) a~$t_i$ je \emph{cíl} (požadovaný výstup, label). Hledáme funkci (model) $y$, která dobře předpovídá cíl i~pro \emph{nové, neviděné} vstupy. Podle typu cíle:
\begin{itemize}
  \item \textbf{regrese} -- cíl je reálné číslo, $t\in\mathbb{R}$ (cena, teplota, výnos);
  \item \textbf{klasifikace} -- cíl je jedna z~$K$ tříd, $t\in\{0,1,\dots,K-1\}$ (spam/ham, druh kytky). Model může vracet jen třídu, nebo celé \emph{rozdělení pravděpodobností} tříd.
\end{itemize}
\end{definice}

\begin{intuice}
Příklady chápeme jako nezávislé vzorky z~neznámého \emph{rozdělení generujícího data}. Cílem není zapamatovat trénovací množinu (to je pouhá \emph{optimalizace}), ale \textbf{generalizovat} -- dobře předpovídat data ze stejného rozdělení, která model neviděl. Optimalizace na trénovacích datech je jen prostředek; tím, co nás zajímá, je chyba na nezávislých datech. Celá tato sekce je o~tom, jak takovou chybu \emph{poctivě změřit}.
\end{intuice}

\subsection{Ohodnocení modelu: testovací data, křížová validace, maximum likelihood}

\subsubsection*{Trénovací, validační a~testovací množina}

Chybu na neviděných datech nelze měřit na trénovací množině (model je na ni „naladěn`` a~chybu podhodnotí). Data proto dělíme:
\begin{itemize}
  \item \textbf{trénovací} -- na nich se učí parametry modelu (váhy);
  \item \textbf{validační} (development) -- na nich ladíme \emph{hyperparametry} (např.\ stupeň polynomu, sílu regularizace, $k$ u~$k$-NN) a~rozhodujeme o~\emph{včasném zastavení}; viz sekce~\ref{sec:preuceni};
  \item \textbf{testovací} -- použijí se \emph{jen jednou na konci} k~odhadu generalizační chyby. Kdo ladí podle testovacích dat, klame sám sebe.
\end{itemize}

\begin{definice}[$k$-násobná křížová validace]
Dat je vždy málo na to, vyčlenit velkou validační množinu. \textbf{$k$-násobná křížová validace} ($k$-fold cross-validation) data rozdělí na $k$ stejných dílů (\emph{foldů}). Postupně bere každý díl jako validační a~zbylých $k-1$ dílů jako trénovací; natrénuje a~změří chybu. Výsledný odhad je \emph{průměr} $k$ takto získaných chyb. Využije tak na validaci postupně \emph{všechna} data, ale model nikdy netestuje na datech, na nichž se učil.
\end{definice}

\begin{center}
\begin{tikzpicture}[scale=1.0]
  \def\W{1.0} \def\H{0.42} \def\G{0.13}
  \foreach \r [count=\ri] in {0,1,2,3,4}{
    \pgfmathsetmacro\y{0-\r*(\H+\G)}
    \pgfmathsetmacro\yc{\y+\H/2}
    \node[anchor=east,font=\scriptsize] at (-0.15,\yc) {běh~\ri};
    \foreach \c in {0,1,2,3,4}{
      \pgfmathsetmacro\x{\c*\W}
      \pgfmathsetmacro\xc{\x+\W/2}
      \ifnum\c=\r
        \draw[fill=yellow!55,draw=black!60] (\x,\y) rectangle ++(\W,\H);
        \node[font=\scriptsize] at (\xc,\yc) {valid.};
      \else
        \draw[fill=blue!12,draw=black!60] (\x,\y) rectangle ++(\W,\H);
        \node[font=\scriptsize] at (\xc,\yc) {trén.};
      \fi
    }
  }
\end{tikzpicture}
\obrpopis{$5$-násobná křížová validace: data jsou v~$5$ dílech, v~každém z~$5$ běhů slouží jiný díl jako validační (žlutě) a~zbylé čtyři k~trénování (modře). Odhad chyby je průměrem chyb z~pěti běhů. Krajní případ $k=N$ (každý příklad sám fold) se nazývá \emph{leave-one-out}.}
\end{center}

\subsubsection*{Maximum likelihood estimation}

Mnoho modelů (lineární i~logistická regrese, naivní Bayes) se učí \emph{pravděpodobnostně}: model definuje pravděpodobnost dat při daných parametrech a~my hledáme parametry, které pozorovaná data činí \emph{nejvíc pravděpodobnými}.

\begin{definice}[Maximum likelihood estimation (MLE)]
Mějme model s~parametry $\mathbf{w}$, který datům přiřazuje pravděpodobnost (věrohodnost) $p(\text{data}\mid\mathbf{w})$. Při nezávislých příkladech se faktorizuje, $p(\text{data}\mid\mathbf{w})=\prod_{i=1}^N p(t_i\mid \mathbf{x}_i;\mathbf{w})$. \textbf{Maximum likelihood estimation} (maximálně věrohodný odhad) parametrů je
\[ \mathbf{w}_{\mathrm{MLE}} = \argmax_{\mathbf{w}} \prod_{i=1}^N p(t_i\mid \mathbf{x}_i;\mathbf{w}) = \argmin_{\mathbf{w}} \sum_{i=1}^N \big(-\log p(t_i\mid \mathbf{x}_i;\mathbf{w})\big). \]
\end{definice}

\begin{intuice}
Logaritmus mění součin na součet a~je rostoucí, takže maximum nezmění; navíc se lépe počítá a~derivuje. Veličina $-\log p(t_i\mid\mathbf{x}_i;\mathbf{w})$ je \emph{negative log-likelihood} (NLL, záporná logaritmická věrohodnost) a~slouží přímo jako \textbf{ztrátová funkce}: minimalizovat NLL $=$ maximalizovat věrohodnost. Pro normálně rozdělený šum dává MLE přesně \emph{nejmenší čtverce} (sekce~\ref{sec:linreg}), pro klasifikaci \emph{křížovou entropii} (sekce~\ref{sec:logisticka}). Maximum likelihood tak propojuje „intuitivní`` ztrátové funkce s~pravděpodobnostním modelem.
\end{intuice}

\subsection{Evaluační metriky pro klasifikaci}

Základem evaluace binárního klasifikátoru je \textbf{matice záměn} (confusion matrix): tabulka, kolik příkladů které skutečné třídy bylo predikováno do které třídy.

\begin{center}
\begin{tikzpicture}[scale=1.0]
  \def\W{2.5} \def\H{1.0}
  % cells
  \fill[green!16] (0,0) rectangle (\W,\H);          % TN
  \fill[red!12]  (\W,0) rectangle (2*\W,\H);         % FP (skut. neg, pred pos)
  \fill[red!12]  (0,\H) rectangle (\W,2*\H);         % FN
  \fill[green!16] (\W,\H) rectangle (2*\W,2*\H);     % TP
  \draw (0,0) grid[step=\W] (2*\W,2*\H); \draw (0,0) rectangle (2*\W,2*\H);
  % cell labels (velká zkratka + malý popisek na druhém řádku, vejde se do buňky)
  \node[align=center] at (1.5*\W,1.5*\H) {\large\textbf{TP}\\[1pt]{\footnotesize true positive}};
  \node[align=center] at (0.5*\W,1.5*\H) {\large\textbf{FN}\\[1pt]{\footnotesize false negative}};
  \node[align=center] at (1.5*\W,0.5*\H) {\large\textbf{FP}\\[1pt]{\footnotesize false positive}};
  \node[align=center] at (0.5*\W,0.5*\H) {\large\textbf{TN}\\[1pt]{\footnotesize true negative}};
  % column headers (predikce) -- bez prefixu "pred." (dimenzi značí tučný titulek "predikce")
  \node[font=\small] at (0.5*\W,2*\H+0.3) {\emph{neg}};
  \node[font=\small] at (1.5*\W,2*\H+0.3) {\emph{poz}};
  \node[font=\small\bfseries] at (\W,2*\H+0.75) {predikce};
  % row headers (skutečnost) -- bez prefixu "skut."
  \node[font=\small,rotate=90] at (-0.35,1.5*\H) {\emph{poz}};
  \node[font=\small,rotate=90] at (-0.35,0.5*\H) {\emph{neg}};
  \node[font=\small\bfseries,rotate=90] at (-0.9,\H) {skutečnost};
\end{tikzpicture}
\obrpopis{Matice záměn binárního klasifikátoru. Zeleně správné predikce (TP, TN), červeně chyby (FP $=$ falešný poplach, FN $=$ přehlédnutí). Z~těchto čtyř čísel se počítají všechny níže uvedené metriky. Pozitivní třída je obvykle ta „zajímavá`` (nemoc, podvod, anomálie).}
\end{center}

\begin{definice}[Metriky binární klasifikace]
Označme TP, TN, FP, FN počty příkladů v~matici záměn. Pak:
\begin{itemize}
  \item \textbf{accuracy} (úspěšnost): podíl správně klasifikovaných,
    \[ \mathrm{acc} = \frac{\mathrm{TP}+\mathrm{TN}}{\mathrm{TP}+\mathrm{TN}+\mathrm{FP}+\mathrm{FN}}. \]
  \item \textbf{precision} (přesnost): podíl skutečně pozitivních mezi predikovanými pozitivními,
    \[ \mathrm{precision} = \frac{\mathrm{TP}}{\mathrm{TP}+\mathrm{FP}}. \]
  \item \textbf{recall} (úplnost), též \emph{sensitivity} a~\emph{true positive rate} (TPR): podíl zachycených ze skutečně pozitivních,
    \[ \mathrm{recall} = \mathrm{TPR} = \frac{\mathrm{TP}}{\mathrm{TP}+\mathrm{FN}}. \]
  \item \textbf{specificity} (specificita), též \emph{true negative rate} (TNR): podíl zachycených ze skutečně negativních,
    \[ \mathrm{specificity} = \mathrm{TNR} = \frac{\mathrm{TN}}{\mathrm{TN}+\mathrm{FP}}. \]
  \item \textbf{F1} (F1 skóre): harmonický průměr precision a~recall,
    \[ F_1 = \frac{2}{\frac{1}{\mathrm{precision}}+\frac{1}{\mathrm{recall}}} = \frac{2\,\mathrm{TP}}{2\,\mathrm{TP}+\mathrm{FP}+\mathrm{FN}}. \]
\end{itemize}
\end{definice}

\begin{intuice}
\emph{Precision} a~\emph{recall} jsou v~napětí: agresivnější klasifikátor (více predikcí „pozitivní``) zvýší recall, ale sníží precision, a~naopak. F1 je vyvažuje jedním číslem; harmonický průměr (na rozdíl od aritmetického) je nízký, je-li \emph{kterákoli} z~obou nízká. Falešný poplach (FP) a~přehlédnutí (FN) mívají \emph{různě velké} náklady (přehlédnutý nádor vs.\ zbytečná biopsie), proto neexistuje jedna „správná`` metrika -- volíme podle úlohy.
\end{intuice}

\begin{poznamka}[Pozor na nevyvážené třídy]
U~silně nevyvážených dat (např.\ $1\,\%$ podvodů) je \emph{accuracy zavádějící}: triviální klasifikátor „vše negativní`` má accuracy $99\,\%$, a~přitom nezachytí jediný podvod. Tam dává smysl sledovat \emph{precision, recall, F1} nebo \emph{AUC}, ne accuracy.
\end{poznamka}

\subsubsection*{ROC křivka a~AUC}

Klasifikátor obvykle vrací \emph{skóre} (např.\ pravděpodobnost pozitivní třídy) a~třídu přiřadí podle \emph{prahu}. Posouváním prahu měníme kompromis mezi TPR a~FPR.

\begin{definice}[ROC křivka, AUC]
\textbf{ROC křivka} (receiver operating characteristic) vynáší pro všechny prahy bod $(\mathrm{FPR},\mathrm{TPR})$, kde $\mathrm{FPR}=\frac{\mathrm{FP}}{\mathrm{FP}+\mathrm{TN}}=1-\mathrm{specificity}$. Diagonála $\mathrm{TPR}=\mathrm{FPR}$ odpovídá náhodnému hádání, levý horní roh $(0,1)$ dokonalému klasifikátoru. \textbf{AUC} (area under curve) je plocha pod ROC křivkou; rovná se pravděpodobnosti, že náhodný pozitivní příklad dostane vyšší skóre než náhodný negativní. AUC${}=1$ je dokonalý, AUC${}=0{,}5$ náhodný klasifikátor.
\end{definice}

\begin{center}
\begin{tikzpicture}
\begin{axis}[width=7.4cm,height=6.2cm,xlabel={FPR $=1-$specificity},ylabel={TPR $=$ recall},
  xmin=0,xmax=1,ymin=0,ymax=1,xtick={0,0.5,1},ytick={0,0.5,1},
  axis lines=left,every axis plot/.append style={thick},clip=false]
  \addplot[name path=roc,blue!60!black,domain=0:1,samples=80]{x^0.32};
  \addplot[name path=diag,black!55,dashed,domain=0:1,samples=2]{x};
  \addplot[blue!12,opacity=0.5] fill between[of=roc and diag];
  \addplot[only marks,mark=*,mark size=1.6pt,red!70!black] coordinates {(0.2,0.597)};
  \node[font=\scriptsize,anchor=west] at (axis cs:0.235,0.5) {práh};
  \node[font=\scriptsize,blue!60!black] at (axis cs:0.55,0.93) {ROC};
  \node[font=\scriptsize,black!55,rotate=33] at (axis cs:0.68,0.55) {náhodné hádání};
  \node[font=\scriptsize] at (axis cs:0.4,0.58) {AUC};
\end{axis}
\end{tikzpicture}
\obrpopis{ROC křivka (modře): každý bod odpovídá jednomu prahu rozhodování. Čím blíž levému hornímu rohu $(0,1)$, tím lepší. Čárkovaná diagonála je náhodné hádání; plocha pod křivkou (AUC) měří kvalitu nezávisle na volbě prahu.}
\end{center}

\subsubsection*{Řešené příklady: matice záměn a~metriky}

\begin{priklad}[SZZ léto 2022 č.~11: Matice konfuze a~specificita]
\szzotazka{léto 2022}{11}{Evaluace binárního klasifikátoru}\par
\textbf{Zadání.} Testovací množina $1000$ příkladů, $50\,\%$ negativních. Klasifikátor má sensitivity $60\,\%$ a~accuracy $70\,\%$. (1)~Napište kompletní matici konfuze. (2)~Vypočtěte specificity.

\smallskip
\textbf{Řešení.} Je $P=500$ pozitivních a~$N=500$ negativních. Sensitivity (recall) $=\mathrm{TP}/P=0{,}6\Rightarrow\mathrm{TP}=300$, $\mathrm{FN}=200$. Accuracy $=(\mathrm{TP}+\mathrm{TN})/1000=0{,}7\Rightarrow\mathrm{TP}+\mathrm{TN}=700$, odtud $\mathrm{TN}=400$ a~$\mathrm{FP}=N-\mathrm{TN}=100$.
\begin{center}
\begin{tabular}{lcc}
\toprule
& predikováno pozitivní & predikováno negativní \\
\midrule
skutečně pozitivní & $\mathrm{TP}=300$ & $\mathrm{FN}=200$ \\
skutečně negativní & $\mathrm{FP}=100$ & $\mathrm{TN}=400$ \\
\bottomrule
\end{tabular}
\end{center}
Specificity $=\dfrac{\mathrm{TN}}{\mathrm{TN}+\mathrm{FP}}=\dfrac{400}{500}=\boxed{0{,}8}=80\,\%$.
\end{priklad}

\begin{priklad}[SZZ podzim 2024 č.~30: Evaluace z~accuracy a~specificity]
\szzotazka{podzim 2024}{30}{Evaluace binárního klasifikátoru}\par
\textbf{Zadání.} Testovací množina má $1000$ příkladů; klasifikátor $400$ označil jako pozitivní a~$600$ jako negativní, s~accuracy $70\,\%$ a~specificity $80\,\%$. Doplňte matici záměn a~spočtěte precision, recall a~F1.

\smallskip
\textbf{Řešení.} Z~predikcí: $\mathrm{TP}+\mathrm{FP}=400$, $\mathrm{TN}+\mathrm{FN}=600$. Z~accuracy $\mathrm{TP}+\mathrm{TN}=0{,}7\cdot 1000=700$. Specificity $\frac{\mathrm{TN}}{\mathrm{TN}+\mathrm{FP}}=0{,}8$. Soustavu vyřeší
\[ \mathrm{TN}=400,\quad \mathrm{FN}=200,\quad \mathrm{TP}=300,\quad \mathrm{FP}=100 \]
(kontrola: $\mathrm{TN}/(\mathrm{TN}+\mathrm{FP})=400/500=0{,}8$ \faCheck, $(\mathrm{TP}+\mathrm{TN})/1000=0{,}7$ \faCheck). Pak
\[ \text{precision}=\tfrac{300}{400}=\boxed{0{,}75},\quad \text{recall}=\tfrac{300}{500}=\boxed{0{,}6},\quad F_1=\tfrac{2\cdot 0{,}75\cdot 0{,}6}{0{,}75+0{,}6}=\boxed{0{,}\overline{6}}. \]
\end{priklad}

\begin{priklad}[SZZ jaro 2023 č.~7: Matice záměn ze sensitivity a~specificity]
\szzotazka{jaro 2023}{7}{Evaluace binárního klasifikátoru}\par
\textbf{Zadání.} Klasifikátor má sensitivity $0{,}8$ a~specificity $0{,}9$; testováno na $10\,000$ příkladech, z~nichž $500$ je pozitivních. (1)~Spočtěte matici záměn. (2)~Vstup je predikován jako pozitivní -- jaká je pravděpodobnost, že je skutečně pozitivní? (3)~Jaká je accuracy?

\smallskip
\textbf{Řešení.} (1)~Pozitivních je $500$, negativních $9500$. Sensitivity (recall) $=\mathrm{TP}/500=0{,}8\Rightarrow \mathrm{TP}=400$, $\mathrm{FN}=100$. Specificity $=\mathrm{TN}/9500=0{,}9\Rightarrow \mathrm{TN}=8550$, $\mathrm{FP}=950$.
\[ \mathrm{TP}=400,\ \mathrm{FN}=100,\ \mathrm{FP}=950,\ \mathrm{TN}=8550. \]
(2)~Hledaná pravděpodobnost je právě \emph{precision}: $\text{precision}=\frac{400}{400+950}=\frac{400}{1350}\approx\boxed{0{,}296}$. (Ač má klasifikátor vysokou sensitivity, kvůli vzácnosti pozitivní třídy je jen $\sim\!30\,\%$ predikovaných pozitivních skutečně pozitivních -- ukázka, jak nevyváženost sráží precision.) (3)~$\mathrm{acc}=\frac{400+8550}{10000}=\boxed{0{,}895}$.
\end{priklad}

\begin{priklad}[SZZ léto 2023 č.~6: ROC a~precision v~jednom bodě]
\szzotazka{léto 2023}{6}{Testování binárního klasifikátoru}\par
\textbf{Zadání.} Testovací sada má $100$ pozitivních a~$400$ negativních příkladů. (a)~ROC křivka prochází bodem $\mathrm{TPR}=\mathrm{FPR}=20\,\%$; spočtěte precision v~tomto bodě. (b)~Jiný klasifikátor má konstantní $\text{precision}=50\,\%$; znázorněte jeho ROC křivku.

\smallskip
\textbf{Řešení.} (a)~$\mathrm{TPR}=0{,}2\Rightarrow \mathrm{TP}=0{,}2\cdot 100=20$; $\mathrm{FPR}=0{,}2\Rightarrow \mathrm{FP}=0{,}2\cdot 400=80$. Precision $=\frac{20}{20+80}=\boxed{0{,}2}$.

(b)~Konstantní precision $\frac{\mathrm{TP}}{\mathrm{TP}+\mathrm{FP}}=\tfrac12$ znamená $\mathrm{FP}=\mathrm{TP}$ při každém prahu. Protože $\mathrm{TPR}=\frac{\mathrm{TP}}{100}$ a~$\mathrm{FPR}=\frac{\mathrm{FP}}{400}=\frac{\mathrm{TP}}{400}$, platí $\mathrm{FPR}=\tfrac{1}{4}\mathrm{TPR}$, tj.\ $\mathrm{TPR}=4\cdot\mathrm{FPR}$. ROC je tedy \emph{úsečka} z~$(0,0)$ do $(\mathrm{FPR},\mathrm{TPR})=(0{,}25,\,1)$ (při $\mathrm{TPR}=1$ je $\mathrm{FP}=\mathrm{TP}=100$, $\mathrm{FPR}=100/400=0{,}25$).
\begin{center}
\begin{tikzpicture}[scale=2.5]
  \draw[osa] (0,0) -- (1.08,0) node[right,font=\scriptsize]{FPR};
  \draw[osa] (0,0) -- (0,1.12) node[above,font=\scriptsize]{TPR};
  \foreach \t in {0.25,0.5,1}{\draw (\t,0.01)--(\t,-0.01) node[below,font=\tiny]{\t}; \draw (0.01,\t)--(-0.01,\t) node[left,font=\tiny]{\t};}
  \draw[black!45,dashed] (0,0)--(1,1);
  \draw[very thick,blue!60!black] (0,0)--(0.25,1);
  \fill[red!70!black] (0.25,1) circle (0.5pt) node[above right,font=\tiny]{$(0{,}25,\,1)$};
\end{tikzpicture}
\obrpopis{ROC křivka klasifikátoru s~konstantní precision $50\,\%$ na datech $100\!:\!400$: úsečka $\mathrm{TPR}=4\,\mathrm{FPR}$ od počátku do $(0{,}25,1)$. Leží nad diagonálou (lepší než náhoda).}
\end{center}
\end{priklad}

\begin{priklad}[SZZ podzim 2025 č.~28: Dva nástroje na detekci anomálií]
\szzotazka{podzim 2025}{28}{Evaluace binárního klasifikátoru}\par
\textbf{Zadání.} $10\,000$ instancí, z~nichž $1\,\%$ ($100$) jsou anomálie (pozitivní), $9\,900$ normálních. Nástroj~1: accuracy $99{,}5\,\%$, zachytí $90\,\%$ anomálií. Nástroj~2: zachytí $81\,\%$ anomálií, $10\,\%$ nahlášených je falešně pozitivních. Spočtěte matice záměn a~metriky a~rozhodněte, který nástroj zvolit.

\smallskip
\textbf{Řešení.}

\textbf{Nástroj 1.} Recall $0{,}9\Rightarrow \mathrm{TP}=90$, $\mathrm{FN}=10$. Správně klasifikováno $0{,}995\cdot 10000=9950$, z~toho $\mathrm{TN}=9950-90=9860$, takže $\mathrm{FP}=9900-9860=40$.
\[ \text{precision}=\tfrac{90}{130}\approx 0{,}69,\quad \text{recall}=0{,}9,\quad F_1\approx 0{,}78,\quad \mathrm{acc}=0{,}995. \]

\textbf{Nástroj 2.} Recall $0{,}81\Rightarrow \mathrm{TP}=81$, $\mathrm{FN}=19$. „$10\,\%$ nahlášených falešně pozitivních`` znamená $\text{precision}=0{,}9=\frac{\mathrm{TP}}{\mathrm{TP}+\mathrm{FP}}$, odkud $\mathrm{FP}=\tfrac{81}{9}=9$, $\mathrm{TN}=9900-9=9891$.
\[ \text{precision}=0{,}9,\quad \text{recall}=0{,}81,\quad F_1\approx 0{,}85,\quad \mathrm{acc}=0{,}9972. \]

\textbf{Závěr.} Nástroj~2 je lepší ve všech metrikách kromě recall ($0{,}81$ vs.\ $0{,}9$). Hlásí mnohem méně falešných poplachů (FP $9$ vs.\ $40$), takže neobtěžuje správce. Naopak nástroj~1 odhalí víc anomálií; jsou-li kritické, může se i~přes nižší precision vyplatit. Jednoznačná odpověď neexistuje, rozhoduje cena FP vs.\ FN.

\smallskip
\noindent\footnotesize V~původním PDF jsou u~nástroje~2 drobné nekonzistence (zaměněný vzorec precision, „$80$ ze $100$`` místo $81$ a~hodnota $\mathrm{FP}$ označená jako $\mathrm{FN}$); zde uvádíme spočítanou konzistentní variantu, viz též poznámka ve fragmentu v~\texttt{priklady/}.
\end{priklad}

\begin{priklad}[SZZ podzim 2023 č.~31: Metriky a~detekce podvodů]
\szzotazka{podzim 2023}{31}{Evaluační metriky}\par
\textbf{Zadání (zkráceně).} (1)~Definujte accuracy, recall, precision, F1. (2)~Detekce podvodných plateb, kde z~milionů jen stovky jsou podvody -- jaké metriky (ne)použít? (3)~Systém~A: recall $0{,}99$, precision $0{,}2$; systém~B: recall $0{,}9$, precision $0{,}95$. Co znamenají a~který zvolit?

\smallskip
\textbf{Řešení.} (1)~Viz definice výše. (2)~Data jsou silně nevyvážená, proto \emph{accuracy je nevhodná} (triviální „vše legitimní`` má accuracy $\approx 100\,\%$). Vhodné jsou \emph{precision a~recall} (případně F1 nebo AUC), které se soustředí na vzácnou pozitivní třídu. (3)~Systém~A odhalí $99\,\%$ podvodů, ale jen $20\,\%$ jeho hlášení jsou skutečné podvody (oproti B zhruba $5\times$ víc hlášení k~prošetření: $0{,}99/0{,}2\approx 4{,}95$ vs.\ $0{,}9/0{,}95\approx 0{,}95$ na jeden podvod). Systém~B odhalí $90\,\%$ podvodů a~$95\,\%$ hlášení sedí. Pokud nás trápí náklady na prošetřování a~obtěžování zákazníků, volíme \emph{B}; pokud je prioritní zachytit i~poslední podvod, volíme \emph{A}. Jednoznačná odpověď neexistuje, rozhoduje argumentace (cena FP vs.\ FN).
\end{priklad}

\subsection{Evaluační metriky pro regresi}

U~regrese měříme \emph{velikost chyby} predikce $\hat{t}_i=y(\mathbf{x}_i)$ vůči cíli $t_i$.

\begin{definice}[Metriky regrese]
\[ \mathrm{MSE}=\frac1N\sum_i (t_i-\hat t_i)^2,\quad \mathrm{RMSE}=\sqrt{\mathrm{MSE}},\quad \mathrm{MAE}=\frac1N\sum_i |t_i-\hat t_i|. \]
RMSE je ve stejných jednotkách jako cíl; oproti MAE \emph{víc penalizuje velké chyby} (kvadrát). Koeficient determinace $R^2=1-\frac{\sum_i (t_i-\hat t_i)^2}{\sum_i (t_i-\bar t)^2}$ udává podíl rozptylu cíle vysvětlený modelem ($R^2=1$ dokonale, $R^2=0$ jako predikce průměrem).
\end{definice}

\subsection{Naivní Bayesův klasifikátor}

\begin{poznamka}[Nad rámec požadavků]
Naivní Bayes není v~požadavcích jmenován, ale je to standardní pravděpodobnostní klasifikátor probíraný v~NPFL129 (vedle $k$-NN) a~objevuje se v~reálné SZZ otázce; uvádíme ho proto jako řešený příklad klasifikace.
\end{poznamka}

\begin{definice}[Naivní Bayesův klasifikátor]
Pro vstup s~příznaky $x_1,\dots,x_D$ přiřadíme třídu $c$ podle \emph{aposteriorní} pravděpodobnosti (Bayesova věta):
\[ P(c\mid x_1,\dots,x_D)=\frac{P(c)\,P(x_1,\dots,x_D\mid c)}{P(x_1,\dots,x_D)}\;\propto\; P(c)\,P(x_1,\dots,x_D\mid c). \]
\textbf{Naivní} předpoklad je \emph{podmíněná nezávislost} příznaků při dané třídě, takže sdruženou věrohodnost rozložíme na součin:
\[ \hat c=\argmax_{c}\; P(c)\prod_{j=1}^{D} P(x_j\mid c). \]
Pravděpodobnosti $P(c)$ (apriorní třídy) a~$P(x_j\mid c)$ odhadneme z~trénovacích dat relativními četnostmi (MLE). Jmenovatel $P(x_1,\dots,x_D)$ je stejný pro všechny třídy, takže pro výběr maxima stačí čitatel.
\end{definice}

\begin{intuice}
„Naivnost`` je v~předpokladu, že příznaky jsou při dané třídě \emph{nezávislé} -- to skoro nikdy přesně neplatí (např.\ déšť a~oblíbená přednáška spolu mohou korelovat), ale klasifikátor i~tak často dobře funguje a~potřebuje málo dat. Bez tohoto předpokladu bychom museli odhadovat celé sdružené rozdělení $P(x_1,\dots,x_D\mid c)$, což je exponenciálně mnoho parametrů.
\end{intuice}

\begin{priklad}[SZZ jaro 2023 č.~16: Naivní Bayes -- půjde student do školy?]
\szzotazka{jaro 2023}{16}{Naivní Bayesovský klasifikátor}\par
\textbf{Zadání.} Z~$13$ pozorování (atributy \emph{Vstanu}, \emph{Oblíbená přednáška?}, \emph{Počasí}; cíl \emph{Jdu do školy}) určete naivním Bayesem, zda student půjde do školy, jestliže \emph{vstane pozdě}, \emph{nemá oblíbenou přednášku} a~\emph{prší}.
\begin{center}\small
\begin{tabular}{ccccc}
\toprule
Den & Vstanu & Oblíbená přednáška? & Počasí & Jdu do školy\\
\midrule
1 & Brzo & Ne & Slunce & Ne\\
2 & Normálně & Ne & Déšť & Ne\\
3 & Brzo & Ne & Slunce & Ano\\
4 & Pozdě & Ne & Slunce & Ano\\
5 & Pozdě & Ano & Slunce & Ano\\
6 & Pozdě & Ano & Déšť & Ne\\
7 & Brzo & Ano & Déšť & Ano\\
8 & Normálně & Ne & Slunce & Ne\\
9 & Normálně & Ano & Slunce & Ano\\
10 & Pozdě & Ano & Slunce & Ano\\
11 & Normálně & Ano & Déšť & Ano\\
12 & Brzo & Ne & Déšť & Ano\\
13 & Brzo & Ano & Slunce & Ano\\
\bottomrule
\end{tabular}
\end{center}

\smallskip
\textbf{Řešení.} Z~$13$ dnů jde do školy $9\times$ (Ano), nejde $4\times$ (Ne): $P(\text{Ano})=\tfrac{9}{13}$, $P(\text{Ne})=\tfrac{4}{13}$. Spočteme podmíněné četnosti potřebných hodnot:
\[
\begin{aligned}
&P(\text{Pozdě}\mid\text{Ano})=\tfrac{3}{9},\ P(\text{obl}{=}\text{Ne}\mid\text{Ano})=\tfrac{3}{9},\ P(\text{Déšť}\mid\text{Ano})=\tfrac{3}{9};\\
&P(\text{Pozdě}\mid\text{Ne})=\tfrac{1}{4},\ P(\text{obl}{=}\text{Ne}\mid\text{Ne})=\tfrac{3}{4},\ P(\text{Déšť}\mid\text{Ne})=\tfrac{2}{4}.
\end{aligned}
\]
Skóre (čitatel) pro obě třídy:
\[ \text{Ano}:\ \tfrac{9}{13}\cdot\tfrac13\cdot\tfrac13\cdot\tfrac13=\tfrac{1}{39}\approx 0{,}0256,\qquad \text{Ne}:\ \tfrac{4}{13}\cdot\tfrac14\cdot\tfrac34\cdot\tfrac24=\tfrac{3}{104}\approx 0{,}0288. \]
Větší je skóre třídy \emph{Ne}, takže naivní Bayes predikuje, že student \textbf{do školy nepůjde} (po normalizaci $P(\text{Ne}\mid\text{data})\approx 53\,\%$, $P(\text{Ano}\mid\text{data})\approx 47\,\%$ -- těsné). Zajímavé je, že přestože „chodí do školy`` je apriorně mnohem častější, kombinace tří příznaků (pozdě, bez oblíbené přednášky, déšť) převáží ve prospěch „nejde``.
\end{priklad}

\subsection{Praktický postup: od dat k modelu}

Úloha nemusí mířit na jeden algoritmus, ale na \emph{celý postup}: rozpoznat typ úlohy, zvolit model, předzpracovat data, vybrat metriku. Následující reálná otázka SZZ takový postup procvičuje -- nejde o~samostatné téma požadavků, ale o~syntézu pojmů této sekce (a~modelů z~dalších sekcí).

\begin{priklad}[SZZ podzim 2025 č.~29: Predikce z~tabulárních dat]
\szzotazka{podzim 2025}{29}{Predikce pomocí tabulárních dat}\par
\textbf{Zadání (zkráceně).} Zemědělský dataset (pH, srážky, teplota, typ hnojiva, odrůda, velikost farmy, zavlažování, nadm.\ výška) má předpovídat \emph{výnos} (t/ha). (1)~Typ úlohy? (2)~Vhodný model a~proč? (3)~Předzpracování a~zacházení s~typy příznaků. (4)~Evaluační metriky?

\smallskip
\textbf{Řešení.} (1)~\emph{Regrese} s~učitelem -- cíl (výnos) je spojité číslo a~máme trénovací data s~cílovými hodnotami. (2)~Pro tabulární data s~mixem číselných i~kategorických příznaků jsou typicky nejlepší \emph{stromové ansámbly} -- \emph{random forest} nebo \emph{gradient boosting} (XGBoost/LightGBM): zvládnou nelinearity i~interakce, nevyžadují škálování (viz sekce~\ref{sec:kombinace}). Lineární regrese nebo MLP fungují s~pečlivějším předzpracováním. (3)~\emph{Kategorické} příznaky (typ hnojiva, odrůda, zavlažování) zakódovat \emph{one-hot}; \emph{číselné} (pH, srážky, $\dots$) zkontrolovat odlehlé hodnoty a~u~modelů citlivých na škálu \emph{standardizovat}; chybějící hodnoty doplnit (průměr/medián, resp.\ modus). (4)~\emph{MAE}, (R)MSE, případně $R^2$.
\end{priklad}

\section{Overfitting a regularizace}\label{sec:preuceni}

Sekce~\ref{sec:evaluace} ukázala, \emph{jak} chybu modelu poctivě měřit. Tato sekce vysvětluje, \emph{proč} se trénovací a~generalizační chyba rozcházejí: model může být na úlohu příliš slabý (\emph{underfitting}, podučení) nebo příliš silný a~zapamatovat si i~šum (\emph{overfitting}, přeučení). Zavedeme \emph{kapacitu modelu} jako škálu mezi těmito extrémy, ukážeme, jak ji ovládat \emph{regularizací} ($L^2$, $L^1$) a~\emph{early stopping}, a~na závěr rozebereme \emph{prokletí dimenzionality} -- proč je vysoká dimenze příznaků sama o~sobě zdrojem overfittingu.

\begin{poznamka}[Vztah k~S1]
Okruh S1 zmiňuje overfitting jen v~souvislosti s~\emph{Ockhamovou břitvou} (preferuj nejjednodušší hypotézu konzistentní s~daty). Zde to rozvádíme do hloubky, kterou S1 výslovně přenechává S3: kapacita modelu, generalizační chyba, regularizační techniky a~prokletí dimenzionality. Konkrétní aplikace $L^2$ na lineární regresi je v~sekci~\ref{sec:linreg}.
\end{poznamka}

\subsection{Generalizační chyba, underfitting a overfitting}

\begin{definice}[Trénovací a~generalizační chyba]
\textbf{Trénovací chyba} je chyba modelu na trénovací množině, na které se učil. \textbf{Generalizační chyba} (chyba zobecnění, též \emph{testovací chyba}) je \emph{očekávaná} chyba na nových, neviděných příkladech ze stejného rozdělení generujícího data; odhadujeme ji chybou na nezávislé testovací množině (sekce~\ref{sec:evaluace}). Cílem učení je malá generalizační chyba, ne malá trénovací chyba.
\end{definice}

\begin{definice}[Underfitting a~overfitting]
\textbf{Underfitting} (podučení) nastane, když model nedokáže zachytit ani zákonitost v~trénovacích datech: trénovací \emph{i} generalizační chyba jsou velké. \textbf{Overfitting} (přeučení) nastane, když se model „naučí`` i~náhodný šum trénovací množiny: trénovací chyba je malá, ale generalizační chyba velká (mezera mezi nimi je velká).
\end{definice}

\begin{intuice}
Klasická ukázka je prokládání zašuměných bodů polynomem (obrázek níže). \emph{Přímka} (stupeň 1) je na sinusovku příliš tuhá -- nastává underfitting. Polynom \emph{vysokého stupně} (stupeň 9, deset koeficientů na $11$ bodů) projde body téměř přesně (trénovací chyba skoro nulová), ale mezi body divoce kmitá: zapamatoval si šum, ne tvar (overfitting). \emph{Střední stupeň} (zde 3) zachytí tvar a~šum ignoruje -- nejlépe generalizuje. Overfittovaný model je jako student, který nabiflováním odpoví na cvičné otázky doslova, ale u~nové otázky selže.
\end{intuice}

\begin{center}
\pgfplotsset{
  fitstyl/.style={width=5.3cm,height=4.4cm,
    xmin=-0.05,xmax=1.05,ymin=-1.7,ymax=1.7,
    xtick=\empty,ytick=\empty,axis lines=left,
    every axis plot/.append style={thick},
    title style={font=\small}},
}
\begin{tikzpicture}
\begin{axis}[fitstyl,title={underfitting (st.\ 1)}]
  \addplot[black!35,dashed] table[x=x,y=y]{src/fig/s2_poly_true.dat};
  \addplot[blue!60!black] table[x=x,y=y]{src/fig/s2_poly_deg1.dat};
  \addplot[only marks,mark=*,mark size=1.2pt,red!70!black] table[x=x,y=t]{src/fig/s2_poly_data.dat};
\end{axis}
\end{tikzpicture}\hfill
\begin{tikzpicture}
\begin{axis}[fitstyl,title={dobrý fit (st.\ 3)}]
  \addplot[black!35,dashed] table[x=x,y=y]{src/fig/s2_poly_true.dat};
  \addplot[green!45!black] table[x=x,y=y]{src/fig/s2_poly_deg3.dat};
  \addplot[only marks,mark=*,mark size=1.2pt,red!70!black] table[x=x,y=t]{src/fig/s2_poly_data.dat};
\end{axis}
\end{tikzpicture}\hfill
\begin{tikzpicture}
\begin{axis}[fitstyl,title={overfitting (st.\ 9)}]
  \addplot[black!35,dashed] table[x=x,y=y]{src/fig/s2_poly_true.dat};
  \addplot[red!60!black] table[x=x,y=y]{src/fig/s2_poly_deg9.dat};
  \addplot[only marks,mark=*,mark size=1.2pt,red!70!black] table[x=x,y=t]{src/fig/s2_poly_data.dat};
\end{axis}
\end{tikzpicture}
\obrpopis{Polynomy stupně 1, 3 a~9 proložené $11$~zašuměnými body (červeně) z~funkce $\sin(2\pi x)$ (čárkovaně pravá křivka). Trénovací RMSE klesá $0{,}56\to0{,}16\to0{,}01$: stupeň 1 nestačí (underfitting), stupeň 9 projde body téměř přesně, ale mimo ně kmitá (overfitting), stupeň 3 vystihne tvar nejlépe. Data vygenerována s~fixním seedem (\texttt{src/fig/gen\_s2\_polyfit.py}).}
\end{center}

\subsection{Kapacita modelu}

\begin{definice}[Kapacita modelu]
\textbf{Kapacita} modelu je neformálně jeho schopnost vyjádřit \emph{složité} funkce -- jak bohatá je třída funkcí, mezi nimiž učení vybírá. Rozlišujeme \emph{reprezentační kapacitu} (jaké funkce model vůbec umí vyjádřit, např.\ polynomy do stupně $M$) a~\emph{efektivní kapacitu} (kterých z~nich učící algoritmus reálně dosáhne, případně omezenou regularizací). U~polynomiální regrese roste kapacita se stupněm $M$; obecně s~počtem volných parametrů.
\end{definice}

\begin{intuice}
Kapacita je „regulátor`` mezi underfittingem a~overfittingem. \emph{Malá} kapacita $\Rightarrow$ underfitting (model je příliš tuhý). \emph{Velká} kapacita $\Rightarrow$ riziko overfittingu (model je tak ohebný, že se přizpůsobí i~šumu). Optimum je uprostřed. Kapacitu lze měnit nejen volbou modelu (stupeň polynomu, hloubka stromu, počet sousedů $k$), ale i~\emph{množstvím dat}: víc trénovacích příkladů relativní kapacitu snižuje (model už nemá dost volnosti zapamatovat si každý bod), takže více dat samo o~sobě overfitting tlumí.
\end{intuice}

\begin{center}
\begin{tikzpicture}
\begin{axis}[width=10.0cm,height=6.4cm,
  xlabel={kapacita modelu (např.\ stupeň polynomu)},ylabel={chyba},
  xmin=0,xmax=10,ymin=0,ymax=3.2,
  xtick=\empty,ytick=\empty,axis lines=left,clip=false,
  every axis plot/.append style={very thick}]
  % trénovací chyba: monotónně klesá
  \addplot[blue!60!black,domain=0.5:9.6,samples=80]{2.6*exp(0-0.42*x)+0.06};
  % generalizační chyba: U-tvar (klesá, pak roste)
  \addplot[red!65!black,domain=0.5:9.6,samples=80]{2.6*exp(0-0.42*x)+0.06+0.045*(x-3.3)^2};
  % optimum kapacity = minimum červené křivky (numericky x = 4,87)
  \draw[black!55,dashed] (axis cs:4.87,0) -- (axis cs:4.87,2.05);
  \node[font=\scriptsize,anchor=south] at (axis cs:4.87,2.05) {optimum};
  % oblasti
  \node[font=\scriptsize,blue!50!black,align=center] at (axis cs:2.05,3.0) {underfitting\\(velká obě chyby)};
  \node[font=\scriptsize,red!55!black,align=center] at (axis cs:7.7,3.0) {overfitting\\(roste gen.\ chyba)};
  % popisky křivek
  \node[font=\scriptsize,blue!60!black,anchor=west] at (axis cs:6.3,0.33) {trénovací chyba};
  \node[font=\scriptsize,red!65!black,anchor=west] at (axis cs:5.9,1.75) {generalizační chyba};
\end{axis}
\end{tikzpicture}
\obrpopis{Schematická závislost chyby na kapacitě modelu. Trénovací chyba (modře) s~rostoucí kapacitou monotónně klesá. Generalizační chyba (červeně) má tvar \emph{U}: nejdřív klesá (mizí underfitting), pak roste (model overfittuje). Optimální kapacita je v~minimu červené křivky; vlevo od něj je oblast underfittingu, vpravo overfittingu. Mezera mezi křivkami je míra overfittingu.}
\end{center}

\subsection{Regularizace}

\begin{definice}[Regularizace]
\textbf{Regularizace} je jakákoli úprava učícího algoritmu, jejímž záměrem je \emph{snížit generalizační chybu} (ne nutně trénovací). Typicky omezuje efektivní kapacitu modelu, takže ten musí dát přednost „jednodušším`` řešením a~hůře se přizpůsobí šumu.
\end{definice}

\begin{intuice}
Omezení kapacity zvenčí (menší stupeň polynomu) je hrubé -- skáče po celých číslech a~rovnou škrtá celé příznaky. Regularizace dovoluje \emph{plynule} škálovat, jak moc model „smí`` přizpůsobit váhy datům: necháme bohatou třídu funkcí, ale přidáme do ztrátové funkce \emph{trest za složitost}. Nejčastěji trestáme \emph{velikost vah} -- malé váhy znamenají hladší, méně rozkmitanou funkci.
\end{intuice}

\subsubsection*{$L^2$ a~$L^1$ regularizace}

\begin{definice}[$L^2$ a~$L^1$ regularizace]
K~původní ztrátové funkci $E_0(\mathbf{w})$ (např.\ součtu čtverců) přidáme penalizační člen úměrný normě vah:
\[
E_{L^2}(\mathbf{w}) = E_0(\mathbf{w}) + \frac{\lambda}{2}\,\|\mathbf{w}\|_2^2
  = E_0(\mathbf{w}) + \frac{\lambda}{2}\sum_{j} w_j^2,
\qquad
E_{L^1}(\mathbf{w}) = E_0(\mathbf{w}) + \lambda\,\|\mathbf{w}\|_1
  = E_0(\mathbf{w}) + \lambda\sum_{j} |w_j|.
\]
\textbf{$L^2$ regularizace} (též \emph{weight decay}, hřebenová / ridge regrese) trestá kvadrát eukleidovské normy; \textbf{$L^1$ regularizace} (lasso) trestá součet absolutních hodnot. \emph{Regularizační síla} $\lambda \ge 0$ je hyperparametr: $\lambda=0$ je bez regularizace, velké $\lambda$ tlačí váhy k~nule. Penalizace se obvykle \emph{nevztahuje na bias} $b$ (přes bias nelze overfittovat).
\end{definice}

\begin{intuice}[$L^2$ vyhlazuje, $L^1$ řídne]
Obě penalizace táhnou váhy k~nule, ale jinak. \textbf{$L^2$} zmenšuje \emph{všechny} váhy úměrně jejich velikosti -- velké srazí hodně, malé téměř nechá; vede k~mnoha \emph{malým, ale nenulovým} vahám, tedy k~hladkému modelu (gradient lineární predikce vůči vstupu jsou přímo váhy, $\nabla_{\mathbf{x}}\,\mathbf{w}^{\!\top}\mathbf{x}=\mathbf{w}$, takže malé váhy $=$ malé změny výstupu). \textbf{$L^1$} naopak tlačí část vah \emph{přesně na nulu} -- provádí tak \emph{výběr příznaků}: výsledný model je \emph{řídký} (sparse), používá jen podmnožinu příznaků. Proč rozdíl, ukazuje geometrie níže.
\end{intuice}

\begin{center}
\begin{tikzpicture}[scale=1.1]
% Geometrie ověřena numericky: střed elips W=(2,0.45), poloosy (1.4,1.0), osově zarovnané.
% Rostoucí kontura se DOTKNE jednotkového kruhu v (0.957,0.289) [mimo osy] (měřítko 0.762)
% a kosočtverce přesně v ROHU (1,0) [w2=0] (měřítko 0.844). Vykreslujeme dotykovou konturu
% jako vnější (obepíná celou oblast a dotýká se jen v bodě řešení).
\foreach \px/\title/\dia in {0/{$L^2$: dotyk mimo osy}/0, 6.6/{$L^1$: dotyk v~rohu}/1}{
  \begin{scope}[shift={(\px,0)}]
    \node[font=\small] at (0.6,2.35) {\title};
    \draw[osa] (-1.6,0) -- (3.2,0) node[below right=-1pt,font=\scriptsize]{$w_1$};
    \draw[osa] (0,-1.6) -- (0,2.1) node[right=1pt,font=\scriptsize]{$w_2$};
    \coordinate (W) at (2.0,0.45);
    \ifnum\dia=0
      % L2: kruh + kontury (vnější má měřítko 0.762, dotyk v (0.957,0.289))
      \foreach \s in {0.30,0.52,0.762}{ \draw[red!55!black] (W) ellipse ({1.4*\s} and {1.0*\s}); }
      \draw[blue!60!black,very thick,fill=blue!10,fill opacity=0.5] (0,0) circle (1.0);
      \fill[red!65!black] (W) circle (1.3pt);
      \node[font=\scriptsize,red!55!black,anchor=south west] at (2.05,0.5) {$\mathbf{w}^\star$};
      \coordinate (S) at (0.957,0.289);
      \fill[black] (S) circle (1.5pt);
      \node[font=\scriptsize,anchor=south west] at (1.0,0.34) {$\hat{\mathbf{w}}$};
      \node[font=\scriptsize,anchor=north,blue!55!black] at (0.0,-1.05) {obě složky $\neq 0$};
    \else
      % L1: kosočtverec + kontury (vnější má měřítko 0.844, dotyk v rohu (1,0))
      \foreach \s in {0.33,0.58,0.844}{ \draw[red!55!black] (W) ellipse ({1.4*\s} and {1.0*\s}); }
      \draw[blue!60!black,very thick,fill=blue!10,fill opacity=0.5]
        (1.0,0) -- (0,1.0) -- (-1.0,0) -- (0,-1.0) -- cycle;
      \fill[red!65!black] (W) circle (1.3pt);
      \node[font=\scriptsize,red!55!black,anchor=south west] at (2.05,0.5) {$\mathbf{w}^\star$};
      \coordinate (S) at (1.0,0);
      \fill[black] (S) circle (1.7pt);
      \node[font=\scriptsize,anchor=south west] at (1.06,0.07) {$\hat{\mathbf{w}}$};
      \node[font=\scriptsize,anchor=north,green!45!black] at (0.0,-1.05) {$w_2=0$ (řídké)};
    \fi
  \end{scope}
}
\end{tikzpicture}
\obrpopis{Geometrie regularizace. Červené elipsy jsou \emph{kontury} (vrstevnice) původní ztráty se středem v~nepenalizovaném optimu $\mathbf{w}^\star_{\text{MLE}}$ (maximum likelihood estimation, odhad maximální věrohodnosti). Regularizace nutí řešení $\hat{\mathbf{w}}$ dovnitř modré omezující oblasti; optimum leží tam, kde se nejmenší kontura oblasti \emph{dotkne}. Vlevo: \emph{$L^2$} oblast je kruh (hladký), dotyk nastane v~obecném bodě s~oběma složkami nenulovými. Vpravo: \emph{$L^1$} oblast je kosočtverec s~\emph{rohy} na osách; elipsa se ho typicky dotkne v~rohu, kde $w_2=0$ -- proto $L^1$ vynuluje váhy a~dává řídké řešení.}
\end{center}

\begin{poznamka}
$\lambda$ a~stupeň polynomu $M$ jsou \emph{hyperparametry} -- učící algoritmus je sám nenastaví, určujeme je zvenčí. Konkrétní dopad $L^2$ na uzavřené řešení lineární regrese (přechod $\mathbf{X}^{\!\top}\mathbf{X}\to\mathbf{X}^{\!\top}\mathbf{X}+\lambda\mathbf{I}$, vždy invertibilní) a~na krok SGD rozebíráme u~modelu, kde se to počítá -- v~sekci~\ref{sec:linreg}.
\end{poznamka}

\subsubsection*{Early stopping}

U~iterativně trénovaných modelů (gradientní sestup) lze kapacitu omezit i~\emph{časem}: čím déle trénujeme, tím víc se model přizpůsobuje datům včetně šumu.

\begin{definice}[Early stopping]
\textbf{Early stopping} (včasné zastavení) sleduje během trénování chybu na \emph{validační} množině a~trénink ukončí (resp.\ použije uložený model) ve chvíli, kdy validační chyba přestane klesat a~začne růst, i~když trénovací chyba dál klesá. Počet kroků (epoch) je tak vlastně hyperparametr laděný na validačních datech.
\end{definice}

\begin{intuice}
Trénovací chyba u~dostatečně kapacitního modelu klesá monotónně až k~nule, ale validační chyba má opět tvar \emph{U}: nejdřív klesá (model se učí zákonitost), pak roste (začíná se učit šum). Bod obratu je optimální okamžik zastavení -- model v~něm má nejlepší generalizaci. Early stopping je oblíbené, protože je levné (nevyžaduje volbu $\lambda$ ani opakované trénování) a~funguje jako forma regularizace: omezuje, jak daleko se váhy od počátku stačí pohnout. (Příbuznou technikou je \emph{dropout} u~neuronových sítí, mimo rozsah tohoto okruhu.)
\end{intuice}

\begin{center}
\begin{tikzpicture}
\begin{axis}[width=9.6cm,height=5.6cm,
  xlabel={epocha trénování},ylabel={chyba},
  xmin=0,xmax=10,ymin=0,ymax=2.6,
  xtick=\empty,ytick=\empty,axis lines=left,clip=false,
  every axis plot/.append style={very thick}]
  % trénovací chyba: klesá k nule
  \addplot[blue!60!black,domain=0.3:9.7,samples=80]{2.2*exp(0-0.55*x)+0.04};
  % validační chyba: U-tvar
  \addplot[red!65!black,domain=0.3:9.7,samples=80]{2.2*exp(0-0.55*x)+0.04+0.06*(x-2.8)^2};
  % bod early stopping = minimum validační křivky (numericky x = 3,95, y = 0,37)
  \draw[black!55,dashed] (axis cs:3.95,0) -- (axis cs:3.95,1.35);
  \fill[black] (axis cs:3.95,0.37) circle (2.2pt);
  \node[font=\scriptsize,anchor=south west] at (axis cs:4.05,0.45) {early stopping};
  \node[font=\scriptsize,blue!60!black,anchor=west] at (axis cs:6.5,0.2) {trénovací};
  \node[font=\scriptsize,red!65!black,anchor=west] at (axis cs:6.7,1.55) {validační};
\end{axis}
\end{tikzpicture}
\obrpopis{Early stopping. Trénovací chyba (modře) klesá k~nule, validační chyba (červeně) nejdřív klesá, pak roste -- model začíná overfittovat. V~minimu validační chyby (černá tečka) trénink zastavíme; další epochy už generalizaci jen zhoršují.}
\end{center}

\subsection{Prokletí dimenzionality}

\begin{definice}[Prokletí dimenzionality]
\textbf{Prokletí dimenzionality} (curse of dimensionality) je souhrn jevů, kvůli nimž se s~rostoucím počtem příznaků $D$ (dimenzí vstupu) učení dramaticky ztěžuje: objem prostoru roste exponenciálně, takže fixní množina dat je ve vysoké dimenzi \emph{řídká} (vzdálené, prázdné okolí každého bodu), a~vzdálenosti mezi body se vzájemně \emph{připodobňují}. K~pokrytí prostoru se stejnou hustotou bodů je třeba \emph{exponenciálně} mnoho příkladů v~$D$.
\end{definice}

\begin{intuice}
Tři tváře téhož jevu. \emph{(1)~Hustota:} $10$ bodů rovnoměrně pokryje úsečku $[0,1]$, ale $[0,1]^{10}$ jich vyžaduje $10^{10}$, aby byla mřížka stejně hustá -- jinak jsou data ve vysoké dimenzi téměř všude prázdná. \emph{(2)~Koncentrace objemu:} skoro celý objem vysokodimenzionální krychle (či koule) leží těsně u~povrchu, takže „typický`` bod je u~kraje, ne uvnitř. \emph{(3)~Koncentrace vzdáleností:} poměr nejbližšího a~nejvzdálenějšího souseda se blíží jedné, takže pojem „nejbližší soused`` ztrácí smysl -- to přímo dopadá na $k$-NN (sekce~\ref{sec:knn}). Důsledek pro učení: víc příznaků zvyšuje kapacitu, ale i~potřebu dat; nadbytečné příznaky proto raději odstraňujeme (výběr příznaků, $L^1$, PCA).
\end{intuice}

\begin{center}
\begin{tikzpicture}
\begin{axis}[width=10.0cm,height=6.2cm,
  xlabel={dimenze $d$},ylabel={podíl objemu ve slupce},
  xmin=0,xmax=50,ymin=0,ymax=1.05,
  xtick={0,10,20,30,40,50},ytick={0,0.25,0.5,0.75,1},
  axis lines=left,clip=false,legend pos=south east,
  legend style={font=\scriptsize,draw=black!30},
  every axis plot/.append style={very thick}]
  \addplot[blue!60!black,domain=1:50,samples=60]{1-(1-0.05)^x};
  \addplot[green!50!black,domain=1:50,samples=60]{1-(1-0.1)^x};
  \addplot[red!65!black,domain=1:50,samples=60]{1-(1-0.2)^x};
  \legend{{$\varepsilon=0{,}05$},{$\varepsilon=0{,}10$},{$\varepsilon=0{,}20$}}
  \draw[black!45,dashed] (axis cs:20,0) -- (axis cs:20,0.878);
  \fill[green!40!black] (axis cs:20,0.878) circle (2pt);
  \node[font=\tiny,anchor=south east] at (axis cs:20,0.9) {$d{=}20,\ \varepsilon{=}0{,}1$};
\end{axis}
\end{tikzpicture}
\obrpopis{Koncentrace objemu: podíl objemu jednotkové krychle $[0,1]^d$, který leží ve vnější \emph{slupce} tloušťky $\varepsilon$ (mimo vnitřní krychli $[\varepsilon,1-\varepsilon]^d$), je $1-(1-2\varepsilon)^d$; zde kreslíme zjednodušeně podíl mimo zmenšenou krychli $[0,1-\varepsilon]^d$, tj.\ $1-(1-\varepsilon)^d$. Pro $\varepsilon=0{,}1$ a~$d=20$ vychází $1-0{,}9^{20}\approx0{,}878$, tj.\ skoro $88\,\%$ objemu je v~tenké slupce u~povrchu; pro $d=50$ je to $1-0{,}9^{50}\approx0{,}995$. Funkce je analytická, bez šumu.}
\end{center}

\begin{priklad}[Ilustrace: $L^1$ vynuluje váhu, $L^2$ ne]
\textbf{Zadání.} Mějme model se dvěma vahami, kde nepenalizované optimum ztráty je $\mathbf{w}^\star=(2,\,0{,}5)$ a~kontury ztráty jsou kružnice $\|\mathbf{w}-\mathbf{w}^\star\|^2=\text{konst.}$ (izotropní). Omezíme řešení podmínkou $\|\mathbf{w}\|\le 1$ ($L^2$), resp.\ $\|\mathbf{w}\|_1\le 1$ ($L^1$). Které řešení je řídké?

\smallskip
\textbf{Řešení.} Při \emph{kruhových} konturách je penalizované řešení průmět $\mathbf{w}^\star$ na omezující oblast (nejbližší bod oblasti k~$\mathbf{w}^\star$).
\emph{$L^2$ (kruh poloměru 1):} nejbližší bod je ve směru $\mathbf{w}^\star$, tj.\ $\hat{\mathbf{w}}=\mathbf{w}^\star/\|\mathbf{w}^\star\| = (2,\,0{,}5)/\sqrt{4{,}25}\approx(0{,}970,\,0{,}243)$ -- \emph{obě} složky nenulové.
\emph{$L^1$ (kosočtverec $|w_1|+|w_2|\le1$):} nejbližší bod ke vzdálenému $\mathbf{w}^\star=(2,\,0{,}5)$ je \emph{roh} $(1,\,0)$ (ověření: vzdálenost $\mathbf{w}^\star$ od $(1,0)$ je $\sqrt{1+0{,}25}\approx1{,}118$ a~každý jiný bod kosočtverce je dál), takže $\hat{\mathbf{w}}=(1,\,0)$ -- druhá váha je \emph{přesně nula}. $L^1$ tedy příznak $2$ vyřadil (řídké řešení), $L^2$ ho jen zmenšil. To je kvantitativní ukázka geometrie z~obrázku výše.
\end{priklad}

\begin{priklad}[SZZ podzim 2022 č.~18: L2 regularizace, porovnání modelů]
\szzotazka{podzim 2022}{18}{Přeučení a regularizace}\par
\textbf{Zadání.} Trénujeme lineární model pro regresi. (1)~Definujte střední čtvercovou chybu (MSE) a~upravenou chybovou funkci pro $L^2$ regularizaci. (2)~Dva lineární modely pro data se čtyřmi atributy,
\[
  M_1(\mathbf{x}) = x_1 - 2x_2 + 4x_3 - x_4 + 3, \qquad
  M_2(\mathbf{x}) = x_2 - 5x_3 + x_4 - 1,
\]
mají na daných datech \emph{stejnou} MSE. Který má menší hodnotu $L^2$-regularizované chyby? (3)~Uveďte alespoň jeden další způsob, jak zlepšit generalizaci.

\smallskip
\textbf{Řešení.} \emph{(1)} Střední čtvercová chyba je
\[
  \mathrm{MSE} = \frac{1}{n}\sum_{i=1}^{n} (y_i - \hat y_i)^2,
\]
kde $\hat y_i$ je predikce a~$y_i$ skutečná hodnota. Při $L^2$ regularizaci přidáme k~MSE penalizaci za velikost vah,
\[
  \mathrm{MSE} + \lambda \sum_{j} w_j^2,\qquad \lambda>0,
\]
přičemž absolutní člen (bias) se \emph{ne}regularizuje.

\emph{(2)} Regularizační člen závisí jen na vahách (bez biasu). Pro $M_1$ jsou váhy $(1,-2,4,-1)$, takže $\sum_j w_j^2 = 1+4+16+1 = \boxed{22}$; pro $M_2$ jsou váhy $(0,1,-5,1)$, takže $\sum_j w_j^2 = 0+1+25+1 = \boxed{27}$. Při stejné MSE má menší regularizovanou chybu model $\boldsymbol{M_1}$ ($22<27$).

\smallskip
\emph{Pozn.} Kdyby se regularizoval i~bias, vyšlo by $M_1$: $22+3^2=31$ a~$M_2$: $27+(-1)^2=28$, a~vyhrál by $M_2$. Právě proto je podstatné, že se bias neregularizuje.

\emph{(3)} Generalizaci lze zlepšit i~více trénovacími daty, výběrem/redukcí příznaků, předčasným zastavením (early stopping), dropoutem (u~neuronových sítí), kombinováním modelů (ensembling), $L^1$ regularizací nebo laděním hyperparametrů (např.\ $\lambda$) křížovou validací.
\end{priklad}

\section{Učení založené na příkladech: $k$ nejbližších sousedů}\label{sec:knn}

Předchozí dvě sekce zavedly evaluaci a~overfitting obecně; teď přichází \emph{první konkrétní model}. \textbf{$k$ nejbližších sousedů} ($k$-NN, $k$-nearest neighbors) je nejjednodušší příklad \emph{učení založeného na příkladech} (instance-based learning): predikci nového vstupu složí přímo z~\emph{uložených trénovacích příkladů, které jsou mu nejpodobnější}. Žádný globální model dat se nestaví.

\begin{poznamka}[Vztah k~S1]
Okruh S1 představuje $k$-NN jako \emph{typický neparametrický model} (drží data místo vah) na úrovni přehledu. Zde to nezdvojujeme -- princip jen krátce zrekapitulujeme a~rovnou jdeme do hloubky, kterou S1 přenechává S3: váhování sousedů, volba metriky a~normalizace, volba $k$ křížovou validací a~výpočetní cena predikce.
\end{poznamka}

\subsection{Princip: líný, neparametrický model}

\begin{definice}[$k$ nejbližších sousedů]
Mějme trénovací množinu $\{(\mathbf{x}_i,t_i)\}_{i=1}^N$ a~\emph{metriku} (funkci vzdálenosti) $d(\cdot,\cdot)$ na vstupním prostoru. Pro nový (dotazovaný) vstup $\mathbf{x}_q$:
\begin{enumerate}
  \item spočti vzdálenosti $d(\mathbf{x}_q,\mathbf{x}_i)$ ke \emph{všem} trénovacím příkladům a~vyber $k$ příkladů s~nejmenší vzdáleností -- jeho \emph{$k$ nejbližších sousedů};
  \item predikci $y(\mathbf{x}_q)$ slož z~cílů těchto sousedů: u~\textbf{klasifikace} většinovým hlasováním, u~\textbf{regrese} (váženým) průměrem (viz níže).
\end{enumerate}
\end{definice}

\begin{intuice}
$k$-NN je \emph{líný} (lazy) model: \textbf{trénování je triviální -- jen se uloží celá trénovací množina}, žádné parametry se neoptimalizují. Veškerá práce se odkládá na okamžik predikce (proto „líný``). Je to zároveň model \emph{neparametrický}: nepředpokládá pevný tvar rozhodovací funkce s~daným počtem vah; „složitost`` hypotézy roste s~počtem dat (čím víc příkladů, tím jemnější hranice umí popsat). Předpoklad za metodou je prostý: \emph{podobné vstupy mají podobné cíle}, tedy lokální hladkost cílové funkce.
\end{intuice}

\begin{algo}[$k$-NN: trénování a~predikce]
\textbf{Trénování}$(\{(\mathbf{x}_i,t_i)\})$:\\
\hspace*{1.5em} ulož celou trénovací množinu \cmt{nic víc -- líné učení}\\[3pt]
\textbf{Predikce}$(\mathbf{x}_q,\,k,\,d)$:\\
\hspace*{1.5em} \textbf{pro} každý trénovací příklad $i=1,\dots,N$:\\
\hspace*{3em} spočti $d_i \leftarrow d(\mathbf{x}_q,\mathbf{x}_i)$ \cmt{vzdálenost k~dotazu}\\
\hspace*{1.5em} $S \leftarrow$ indexy $k$ nejmenších $d_i$ \cmt{$k$ nejbližších sousedů}\\
\hspace*{1.5em} \textbf{klasifikace:} vrať nejčastější třídu mezi $\{t_i : i\in S\}$ \cmt{remízy libovolně}\\
\hspace*{1.5em} \textbf{regrese:} vrať (vážený) průměr $\{t_i : i\in S\}$
\end{algo}

\subsection{Klasifikace: většinové hlasování}

U~klasifikace do $K$ tříd se nový vstup zařadí do \textbf{nejčastější třídy} mezi jeho $k$ sousedy (\emph{majority voting}). Remízy (např.\ $2\!:\!2$ při $k=4$) se řeší libovolně -- typicky volbou liché $k$ u~dvou tříd, náhodně, nebo podle bližšího souseda.

\begin{intuice}
Hlasování lze chápat jako lokální odhad pravděpodobností tříd: podíl sousedů třídy $c$ je odhad $P(c\mid\mathbf{x}_q)$ v~okolí dotazu, a~vybíráme nejpravděpodobnější třídu. Při $k=1$ dostaneme přímo třídu nejbližšího příkladu; rozhodovací hranice je pak \emph{Voroného diagram} trénovacích bodů (každý bod „vlastní`` okolí složené z~míst, kde je nejbližší).
\end{intuice}

\begin{poznamka}[Vážené hlasování]
Sousedy lze \emph{vážit} podle blízkosti: bližší soused má větší hlas. Běžné volby vah $w_i$ jsou \emph{uniformní} ($w_i=1$, prostá většina), \emph{inverzní} ($w_i=1/d_i$, váha úměrná převrácené vzdálenosti) nebo \emph{softmax} záporných vzdáleností. Predikovaná třída je pak $\argmax_{c}\sum_{i\in S,\,t_i=c} w_i$ (váhovaný součet hlasů). Vážení zjemňuje vliv vzdálenějších, méně relevantních sousedů.
\end{poznamka}

\subsection{Regrese: (vážený) průměr sousedů}

U~regrese je cíl reálné číslo a~predikcí je \textbf{průměr cílů} $k$ sousedů. S~vahami $w_i$ (uniformní $\Rightarrow$ prostý průměr):
\[
  y(\mathbf{x}_q) \;=\; \sum_{i\in S} \frac{w_i}{\sum_{j\in S} w_j}\; t_i .
\]
Pro klasifikaci téhož vzorce: cíle $t_i$ zakódujeme \emph{one-hot} (vektor s~jedničkou na pozici třídy) a~váženě zprůměrujeme, čímž dostaneme rozdělení pravděpodobností tříd $\mathbf{t}=\sum_{i\in S}\frac{w_i}{\sum_j w_j}\mathbf{t}_i$; predikovaná třída je $\argmax_c t_{\,c}$. Uniformní váhy a~argmax pak dají právě většinové hlasování.

\subsection{Vzdálenostní metriky a~normalizace příznaků}

Výsledek $k$-NN zcela stojí na tom, \emph{jak měříme blízkost}. Pro $\mathbf{x},\mathbf{y}\in\mathbb{R}^D$ je standardní volbou \textbf{Minkowského} ($L^p$) vzdálenost
\[
  d_p(\mathbf{x},\mathbf{y}) \;=\; \|\mathbf{x}-\mathbf{y}\|_p \;=\; \Big(\sum_{i=1}^{D} |x_i-y_i|^p\Big)^{1/p},
\]
jejímiž speciálními případy jsou \emph{Manhattanská} ($p=1$, součet absolutních rozdílů), \emph{Eukleidovská} ($p=2$, „vzdušnou čarou``) a~v~limitě \emph{Čebyševova} ($p\to\infty$, $\max_i |x_i-y_i|$). Nejběžnější je Euklid ($p=2$).

\begin{intuice}[Proč normalizovat příznaky]
Metriky $L^p$ sčítají rozdíly přes všechny příznaky, takže příznak s~\emph{větším rozsahem hodnot} (např.\ příjem $0$--$100\,000$ vs.\ věk $0$--$100$) vzdálenost zcela ovládne a~ostatní příznaky se přehlasují. Před $k$-NN proto příznaky \textbf{normalizujeme} na srovnatelný rozsah: buď \emph{standardizací} (každý příznak na nulový průměr a~jednotkový rozptyl, $x'\!=\!(x-\mu)/\sigma$), nebo \emph{min-max škálováním} na $[0,1]$. Bez normalizace by „nejbližší soused`` znamenal pouze „nejbližší v~jednom dominantním příznaku``.
\end{intuice}

\begin{poznamka}[Prokletí dimenzionality]
Ve vysoké dimenzi přestává pojem „blízký soused`` dávat smysl: vzdálenosti k~nejbližšímu a~nejvzdálenějšímu bodu se k~sobě přibližují, takže všechny příklady jsou zhruba stejně daleko a~hlasování ztrácí výpovědní hodnotu. $k$-NN je proto na vysokodimenzionální data citlivý (\emph{prokletí dimenzionality}, podrobně sekce~\ref{sec:preuceni}); pomáhá redukce dimenze (PCA, sekce~\ref{sec:bezucitele}) nebo výběr příznaků.
\end{poznamka}

\subsection{Volba $k$: vliv na hranici, overfitting, křížová validace}

Počet sousedů $k$ je \textbf{hyperparametr} -- neučí se z~dat, volíme ho zvenčí. Řídí kompromis mezi \emph{overfittingem} (přeučení, velký rozptyl) a~\emph{underfittingem} (podučení, velké vychýlení); viz sekce~\ref{sec:preuceni}:

\begin{itemize}
  \item \textbf{malé $k$} (v~krajnosti $k=1$): predikce závisí na jednom či pár nejbližších bodech, je velmi citlivá na šum a~odlehlé příklady. Rozhodovací hranice je \emph{členitá, „roztřepená``}, s~ostrůvky kolem jednotlivých bodů -- model má \emph{velký rozptyl} a~\emph{overfittuje} (trénovací chyba u~$k=1$ je nulová, ale generalizace špatná);
  \item \textbf{velké $k$}: hlasuje hodně sousedů, predikce je vyhlazená a~odolná vůči šumu, ale hranice je \emph{hladká až příliš} a~ztrácí lokální detail -- model má \emph{velké vychýlení} a~\emph{underfittuje}. V~krajnosti $k=N$ vrací vždy globálně nejčastější třídu (resp.\ celkový průměr) bez ohledu na vstup.
\end{itemize}

\begin{center}
\begin{tikzpicture}[scale=0.92]
  % ---- LEVÝ PANEL: malé k (členitá hranice) ----
  \begin{scope}
    \def\xmax{4.0} \def\ymaxv{3.4}
    \draw[osa] (0,0) -- (\xmax+0.25,0) node[right,font=\scriptsize]{$x_1$};
    \draw[osa] (0,0) -- (0,\ymaxv+0.25) node[above,font=\scriptsize]{$x_2$};
    % roztřepená hranice (silně lokální)
    \draw[hranice]
      (0,2.55) .. controls (0.7,2.2) and (1.0,2.9) .. (1.5,2.45)
      .. controls (1.9,2.1) and (2.0,2.75) .. (2.5,2.3)
      .. controls (2.9,1.9) and (3.3,2.6) .. (\xmax,2.15);
    % izolovaný ostrůvek (overfitting na šumový bod)
    \draw[hranice] (1.55,1.0) circle (0.42);
    \node[cls1] at (1.5,2.9) {}; \node[cls1] at (0.6,2.95) {};
    \node[cls1] at (1.1,3.1) {}; \node[cls1] at (2.1,3.05) {};
    \node[cls1] at (2.8,2.85) {}; \node[cls1] at (3.4,2.95) {};
    \node[cls1] at (0.35,3.15) {}; \node[cls1] at (3.6,2.55) {};
    \node[cls1] at (1.55,1.0) {}; % šumový modrý bod dole -> ostrůvek
    \node[cls2] at (0.7,1.6) {}; \node[cls2] at (1.5,1.7) {};
    \node[cls2] at (2.4,1.5) {}; \node[cls2] at (3.1,1.55) {};
    \node[cls2] at (0.5,0.7) {}; \node[cls2] at (1.2,0.6) {};
    \node[cls2] at (2.0,0.8) {}; \node[cls2] at (2.8,0.65) {};
    \node[cls2] at (3.5,1.0) {}; \node[cls2] at (2.2,0.3) {};
    \node[font=\small,anchor=north] at (\xmax/2,-0.25) {malé $k$ (např.\ $k=1$)};
  \end{scope}
  % ---- PRAVÝ PANEL: velké k (hladká hranice) ----
  \begin{scope}[xshift=5.6cm]
    \def\xmax{4.0} \def\ymaxv{3.4}
    \draw[osa] (0,0) -- (\xmax+0.25,0) node[right,font=\scriptsize]{$x_1$};
    \draw[osa] (0,0) -- (0,\ymaxv+0.25) node[above,font=\scriptsize]{$x_2$};
    % hladká hranice
    \draw[hranice] (0,2.4) .. controls (1.3,2.2) and (2.7,2.1) .. (\xmax,2.0);
    % stejné body, žádný ostrůvek (šumový bod přehlasován)
    \node[cls1] at (1.5,2.9) {}; \node[cls1] at (0.6,2.95) {};
    \node[cls1] at (1.1,3.1) {}; \node[cls1] at (2.1,3.05) {};
    \node[cls1] at (2.8,2.85) {}; \node[cls1] at (3.4,2.95) {};
    \node[cls1] at (0.35,3.15) {}; \node[cls1] at (3.6,2.55) {};
    \node[cls1] at (1.55,1.0) {}; % týž šumový bod, teď bez ostrůvku
    \node[cls2] at (0.7,1.6) {}; \node[cls2] at (1.5,1.7) {};
    \node[cls2] at (2.4,1.5) {}; \node[cls2] at (3.1,1.55) {};
    \node[cls2] at (0.5,0.7) {}; \node[cls2] at (1.2,0.6) {};
    \node[cls2] at (2.0,0.8) {}; \node[cls2] at (2.8,0.65) {};
    \node[cls2] at (3.5,1.0) {}; \node[cls2] at (2.2,0.3) {};
    \node[font=\small,anchor=north] at (\xmax/2,-0.25) {velké $k$};
  \end{scope}
\end{tikzpicture}
\obrpopis{Vliv $k$ na rozhodovací hranici $k$-NN (dvě třídy: \tikz[baseline={(0,-0.5ex)}]{\node[cls1]{};}~modré, \tikz[baseline={(0,-0.5ex)}]{\node[cls2]{};}~červené). Při \emph{malém} $k$ (vlevo) je hranice členitá a~kolem osamělého šumového bodu vznikne ostrůvek -- model overfittuje. Při \emph{velkém} $k$ (vpravo) je hranice hladká, šumový bod je sousedy přehlasován; přílišné $k$ však ztrácí lokální detail (underfitting). Optimální $k$ leží mezi.}
\end{center}

\begin{intuice}[Jak vybrat $k$]
Vhodné $k$ nehádáme -- vybereme ho \textbf{křížovou validací} (sekce~\ref{sec:evaluace}): pro kandidátní hodnoty $k=1,3,5,\dots$ změříme $k$-násobnou křížovou validací chybu a~zvolíme $k$ s~nejlepším \emph{validačním} skóre. Trénovací chyba k~volbě \emph{nestačí} (u~$k=1$ je vždy nulová, přesto overfittuje). U~dvou tříd se volí \emph{liché} $k$, aby se předešlo remízám v~hlasování. Závislost validační chyby na $k$ má typicky tvar písmene U: zprvu klesá (mizí overfitting), pak roste (underfitting).
\end{intuice}

\subsection{Výpočetní cena predikce}

Líné učení má svou cenu na druhé straně. \textbf{Naivní predikce je drahá}: pro každý dotaz se spočítá vzdálenost ke \emph{všem} $N$ trénovacím bodům ($D$-rozměrným), tedy $\mathcal{O}(N D)$ na výpočet vzdáleností plus výběr $k$ nejmenších ($\mathcal{O}(N\log k)$ haldou, nebo $\mathcal{O}(N\log N)$ tříděním). Trénování je $\mathcal{O}(1)$ (jen uložení), ale predikce roste lineárně s~velikostí trénovací množiny -- opak parametrických modelů, kde je drahé trénování a~levná predikce.

\begin{poznamka}[Zrychlení vyhledávání]
Predikci lze urychlit prostorovými datovými strukturami, které předpočítají organizaci bodů: \textbf{$k$-d~stromy} (dotaz v~průměru v~logaritmickém čase pro nízkou dimenzi, ale ve vysoké dimenzi degradují), \emph{ball trees}, \emph{R-stromy} apod. Ve vysoké dimenzi nebo na obrovských datech se používá \emph{přibližné} hledání nejbližších sousedů (ANN). Tyto struktury jsou nad rámec požadavků; podstatné je, že naivní predikce je $\mathcal{O}(ND)$ na dotaz a~že existují stromové struktury pro zrychlení.
\end{poznamka}

\subsection{Řešený příklad}

\begin{priklad}[SZZ jaro 2025 č.~28: $k$ nejbližších sousedů]
\szzotazka{jaro 2025}{28}{k nejbližších sousedů}\par
\textbf{Zadání.} (1)~Popište metodu $k$-NN pro klasifikaci i~regresi; jak probíhá trénování a~predikce? (2)~Data se dvěma atributy se klasifikují do dvou tříd ($\circ$, $\times$) podle obrázku. Pro $k=3$ klasifikujte nový vstup $(3,2)$. Změnila by se klasifikace při jiném $k$? (3)~Jak nastavit vhodné $k$?

\smallskip
\textbf{Řešení.}

\emph{(1)} Při \textbf{trénování} se pouze uloží celá trénovací množina (líné učení). Při \textbf{predikci} se k~dotazu najde $k$ nejbližších uložených příkladů (zvolenou metrikou, typicky Euklid); pro \emph{klasifikaci} se vrátí jejich nejčastější třída (většinové hlasování), pro \emph{regresi} (vážený) průměr jejich cílů.

\emph{(2)} Dotaz $(3,2)$ leží na rozhraní obou tříd (viz obrázek). Dva nejbližší body jsou \emph{kroužky} ($\circ$) těsně nad ním, $(3{,}25;2{,}29)$ a~$(2{,}86;2{,}37)$ ve vzdálenostech $\approx 0{,}38$ a~$0{,}40$. Třetím sousedem je křížek: dva křížky pod dotazem, $(2{,}81;1{,}46)$ a~$(2{,}68;1{,}52)$, leží skoro stejně daleko ($\approx 0{,}58$), takže nezáleží na tom, který z~nich je třetí. Hlasování $2\!:\!1$ klasifikuje vstup jako $\boxed{\circ}$.

Při \emph{změně} $k$ se výsledek může změnit: dotaz je blízko hranice a~hned za dvěma nejbližšími kroužky následují křížky. Pro $k=5$ rozhoduje téměř remíza vzdáleností: pátý a~šestý soused, kroužek $(2{,}84;2{,}65)$ a~křížek $(2{,}35;1{,}82)$, jsou oba ve vzdálenosti $\approx 0{,}67$ a~rozdíl je pod přesností odečtu z~obrázku (poměr hlasů $3\!:\!2$ nebo $2\!:\!3$). Pro $k=7$ jsou sousedé jednoznačně $3\times\circ$, $4\times\times$ (sedmý je křížek ve vzdálenosti $\approx 0{,}70$, osmý až $\approx 0{,}89$), takže hlasování přejde na $\boxed{\times}$. Větší $k$ tu „dosáhne`` do husté oblasti křížků pod dotazem a~klasifikaci překlopí -- ukázka citlivosti na volbu $k$ u~bodu blízko rozhodovací hranice.

\emph{(3)} Vhodné $k$ zvolíme \textbf{křížovou validací}: pro $k=1,3,5,\dots$ změříme křížově validovanou chybu a~vybereme $k$ s~nejlepším validačním skóre. U~dvou tříd volíme liché $k$ kvůli remízám. Malé $k$ hrozí overfittingem (citlivost na šum), velké $k$ underfittingem (hranice příliš hladká).

\smallskip
\begin{center}
% Souřadnice = rekonstrukce dat SZZ jaro 2025 č.28 (gen_s3_knn.py); okolí dotazu přeodečteno 26. 9. 2026.
% Jednotka os zadaná přes x=/y=..cm, aby šly použít celočíselné i záporné
% souřadnice bez aritmetiky (vyhneme se "Missing number" pod babel-czech).
\begin{tikzpicture}[x=1.1cm,y=1.1cm]
  \def\xlo{-1.9} \def\xhi{5.0} \def\ylo{-1.3} \def\yhi{5.3}
  % osy a popisky dělení (osa x kreslena u y=ylo, aby byly vidět i body se záporným x2)
  \draw[osa] (\xlo,\ylo) -- (\xhi,\ylo) node[right,font=\scriptsize]{$x_1$};
  \draw[osa] (\xlo,\ylo) -- (\xlo,\yhi) node[above,font=\scriptsize]{$x_2$};
  \foreach \t in {-1,0,1,2,3,4}{\draw (\t,\ylo+0.07)--(\t,\ylo-0.07) node[below,font=\tiny]{$\t$};}
  \foreach \t in {-1,0,1,2,3,4,5}{\draw (\xlo+0.07,\t)--(\xlo-0.07,\t) node[left,font=\tiny]{$\t$};}
  % --- třída o (kroužky), horní/levá oblast: věrná data ---
  \foreach \p in {(-1.6,3.87),(-0.7,4.83),(-0.55,3.73),(-0.6,3.22),(-0.3,4.05),
                  (-0.25,4.08),(-0.1,4.15),(0.05,4.15),(0.1,3.4),(0.15,2.95),
                  (0.3,3.45),(0.35,3.1),(0.4,3.1),(0.55,4.3),(0.6,3.12),
                  (0.65,2.5),(0.9,3.68),(1,4.73),(1.05,3.62),(1.1,3.68),
                  (1.15,4.48),(1.2,3.65),(1.25,4.1),(1.3,3.08),(1.35,2.87),
                  (1.4,3.9),(1.55,3.63),(1.7,3.25),(1.75,2.82),(2,2.45),
                  (2.2,4.27),(2.35,3.25),(2.45,4.8),(2.5,4.47),(2.55,3.22),
                  (2.84,2.65),(2.86,2.37),(3.25,2.29),(0.15,2.32),(0.25,2.1),(0.25,1.9)}
    {\node[cls1] at \p {};}
  % --- třída x (křížky), dolní/pravá oblast: věrná data ---
  \foreach \p in {(0.6,1.13),(0.65,1.45),(0.7,0.5),(0.75,0.5),(0.9,2.27),
                  (0.95,0.87),(1,1.35),(1.1,1.67),(1.2,0.78),(1.25,0.75),
                  (1.3,2.22),(1.3,0.25),(1.35,1.97),(1.35,1.27),(1.4,2.27),
                  (1.45,0.68),(1.5,0.27),(1.6,1.43),(1.7,2.22),(1.75,-0.2),
                  (2,2.1),(2,0.33),(2.05,-0.05),(2.1,1.3),(2.35,1.82),
                  (2.68,1.52),(2.6,0.67),(2.81,1.46),(2.7,0.32),(2.7,-0.38),
                  (2.95,1.12),(3,0.9),(3,0.7),(3.05,1.07),(3.16,1.32),
                  (3.2,1.08),(3.2,0),(2.45,0),(1.35,-0.85),(1.35,-0.33),
                  (3.95,1.93),(4.1,1.83),(4.2,1.5),(4.5,1.53)}
    {\node[cls2] at \p {};}
  % --- dotaz a kružnice k=3 (3. a 4. soused jsou křížky ve vzdálenosti 0,577 a 0,581; kružnice prochází oběma) ---
  \draw[blue!55!black,dashed,thick] (3,2) circle [radius=0.58];
  \fill[black] (3,2) circle [radius=1.7pt];
  % popisek dotazu s vynašecí čárou do volné mezery vpravo dole (mimo body i kružnici)
  \draw[black!55,line width=0.4pt] (3.06,1.92) -- (3.5,0.85);
  \node[font=\scriptsize,anchor=west] at (3.45,0.75) {$q=(3,2)$};
  \node[font=\scriptsize,blue!55!black,anchor=south west] at (3.18,2.6) {$k=3$};
\end{tikzpicture}
\obrpopis{Data SZZ jaro 2025 č.~28 (rekonstrukce z~obrázku zadání, okolí dotazu odečteno přesně): třída \tikz[baseline={(0,-0.5ex)}]{\node[cls1]{};}~$\circ$ nahoře vlevo, třída \tikz[baseline={(0,-0.5ex)}]{\node[cls2]{};}~$\times$ dole vpravo. Čárkovaná kružnice kolem dotazu $q=(3,2)$ má poloměr $\approx 0{,}58$: uvnitř jsou dva $\circ$ (vzdálenosti $\approx 0{,}38$ a~$0{,}40$), kružnice prochází dvěma $\times$ ve vzdálenosti $\approx 0{,}58$, z~nichž jeden je třetím sousedem. Hlasování $2\!:\!1$ klasifikuje $q$ jako $\circ$. Při větším $k$ (kolem $k=7$) se do okolí dostanou křížky pod dotazem a~klasifikace se překlopí na $\times$.}
\end{center}
\end{priklad}

\begin{priklad}[SZZ podzim 2022 č.~20: k-NN s Manhattanskou vzdáleností]
\szzotazka{podzim 2022}{20}{k nejbližších sousedů}\par
\textbf{Zadání.} (1)~Popište fungování $k$-NN pro klasifikaci i~regresi -- jak probíhá trénování a~predikce? (2)~Pro data níže (příznaky $A_1,A_2$, třída $C\in\{0,1\}$) ohodnoťte $k$-NN s~$k=3$ a~\emph{Manhattanskou} (L1) vzdáleností nový vstup $(5,5)$. (3)~Změní se odpověď při jiném $k$?

\smallskip
\begin{center}\small
\begin{tabular}{ccc}
\toprule
$A_1$ & $A_2$ & $C$ \\
\midrule
0 & 3 & 0 \\ 2 & 2 & 0 \\ 3 & 0 & 0 \\ 4 & 2 & 0 \\
6 & 6 & 1 \\ 4 & 5 & 1 \\ 1 & 1 & 1 \\
\bottomrule
\end{tabular}
\end{center}

\smallskip
\textbf{Řešení.} \emph{(1)} $k$-NN je líný (instance-based) model: trénování pouze uloží trénovací data. Predikce nového bodu najde jeho $k$ nejbližších sousedů podle zvolené metriky; pro \emph{klasifikaci} vrátí jejich většinovou třídu, pro \emph{regresi} (vážený) průměr jejich cílů.

\emph{(2)} Manhattanské vzdálenosti bodu $(5,5)$ od dat (seřazeno):
\begin{center}\small
\begin{tabular}{cc}
\toprule
bod & L1 vzdálenost \\
\midrule
$(4,5)$ & $1$ \\ $(6,6)$ & $2$ \\ $(4,2)$ & $4$ \\ $(2,2)$ & $6$ \\
$(0,3)$ & $7$ \\ $(3,0)$ & $7$ \\ $(1,1)$ & $8$ \\
\bottomrule
\end{tabular}
\end{center}
Tři nejbližší jsou $(4,5)$ s~$C{=}1$, $(6,6)$ s~$C{=}1$ a~$(4,2)$ s~$C{=}0$, tedy třídy $\{1,1,0\}$. Pro $k=3$ je $(5,5)$ klasifikován jako třída $\boxed{1}$.

\emph{(3)} Ano, odpověď se mění. Pro $k=5$ je pět nejbližších $(4,5){:}1$, $(6,6){:}1$, $(4,2){:}0$, $(2,2){:}0$, $(0,3){:}0$, tedy $\{1,1,0,0,0\}$, většina $\boxed{0}$ (a~stejně tak $k=7$). Na pátém místě je sice remíza vzdáleností ($(0,3)$ i~$(3,0)$ mají vzdálenost $7$), oba body ale patří do třídy $0$, takže na výsledku nezáleží. Bod leží blízko rozhraní tříd, proto je citlivý na volbu $k$.
\end{priklad}

\section{Lineární regrese}\label{sec:linreg}

Lineární regrese je nejjednodušší regresní model: predikuje cíl jako \emph{lineární kombinaci} příznaků. Přesto je klíčová -- má \textbf{analytické (uzavřené) řešení} metodou nejmenších čtverců, dá se trénovat i~\textbf{gradientním sestupem} (SGD) a~na jejím příkladu se nejlépe ukáže $L^2$ regularizace. Okruh S1 ji zmiňuje přehledově; zde rozebíráme obojí řešení do hloubky, kterou S1 přenechává tomuto okruhu.

\subsection{Model a chybová funkce}

\begin{definice}[Model lineární regrese]
Pro vstup $\mathbf{x}\in\mathbb{R}^D$ predikuje \textbf{lineární regrese}
\[ y(\mathbf{x};\mathbf{w},b)= w_1 x_1 + w_2 x_2 + \dots + w_D x_D + b = \mathbf{x}^{\!\top}\mathbf{w} + b. \]
Vektor $\mathbf{w}\in\mathbb{R}^D$ jsou \emph{váhy}, $b$ je \emph{bias} (posun, absolutní člen). Ve statistickém značení regrese se píší jako koeficienty $\beta_1,\dots,\beta_D$ (váhy) a~$\beta_0$ (bias); jde o~totéž.
\end{definice}

\begin{intuice}[Bias jako další váha (padding jedničkou)]
S~biasem se nepohodlně počítá zvlášť. Trik: každý vstup \textbf{rozšíříme o~konstantní příznak~$1$}, tj.\ $\mathbf{x}\mapsto(x_1,\dots,x_D,1)$, a~bias se stane poslední váhou. Pak píšeme jednotně $y(\mathbf{x})=\mathbf{x}^{\!\top}\mathbf{w}$ a~„váhy`` už zahrnují i~bias. V~maticovém zápisu to znamená přidat do datové matice $\mathbf{X}$ \emph{sloupec jedniček}. Této úmluvy se držíme v~celé sekci.
\end{intuice}

\begin{definice}[Chybová funkce: MSE a součet čtverců]
Máme trénovací data $\mathbf{x}_1,\dots,\mathbf{x}_N$ s~cíli $t_1,\dots,t_N$. Kvalitu vah měříme \textbf{střední kvadratickou chybou} (mean squared error)
\[ \mathrm{MSE}(\mathbf{w})=\frac1N\sum_{i=1}^N\big(y(\mathbf{x}_i;\mathbf{w})-t_i\big)^2. \]
Často se místo ní minimalizuje \textbf{součet čtverců} (sum of squares)
\[ E(\mathbf{w})=\frac12\sum_{i=1}^N\big(y(\mathbf{x}_i;\mathbf{w})-t_i\big)^2. \]
Obě mají \emph{stejné minimum} (MSE má před součtem faktor $\tfrac1N$, součet čtverců faktor $\tfrac12$ -- liší se tedy jen kladnou multiplikativní konstantou), ale součet čtverců se příjemněji derivuje -- faktor $\tfrac12$ se vykrátí proti dvojce z~derivace kvadrátu.
\end{definice}

\begin{intuice}
Hledáme přímku (obecně nadrovinu), která prochází mraku bodů „co nejblíž`` -- tak, aby \emph{součet čtverců svislých vzdáleností} bodů od přímky byl nejmenší. Svislá vzdálenost $\hat t_i - t_i$ se nazývá \emph{reziduum}. Kvadrát (místo absolutní hodnoty) penalizuje velké odchylky silněji a~činí funkci hladkou a~konvexní, takže každé lokální minimum je globální (a~jediné, je-li $\mathbf{X}^{\!\top}\mathbf{X}$ regulární, viz níže). Vztah k~věrohodnosti: pro normálně rozdělený šum kolem přímky dává metoda nejmenších čtverců přesně maximum likelihood estimation (MLE, maximálně věrohodný odhad; viz sekce~\ref{sec:evaluace}).
\end{intuice}

\subsection{Analytické řešení: metoda nejmenších čtverců}

Vyjádříme součet čtverců maticově. Datová matice $\mathbf{X}\in\mathbb{R}^{N\times D}$ má příklad $\mathbf{x}_i^{\!\top}$ na $i$-tém řádku (po paddingu jedničkou má $D$ sloupců včetně biasu), $\mathbf{t}\in\mathbb{R}^N$ je vektor cílů. Pak $(\mathbf{X}\mathbf{w})_i=\mathbf{x}_i^{\!\top}\mathbf{w}$ a
\[ E(\mathbf{w})=\frac12\sum_{i=1}^N(\mathbf{x}_i^{\!\top}\mathbf{w}-t_i)^2=\frac12\,\|\mathbf{X}\mathbf{w}-\mathbf{t}\|^2. \]

\begin{veta}[Normální rovnice nejmenších čtverců]
Minimum součtu čtverců $E(\mathbf{w})=\tfrac12\|\mathbf{X}\mathbf{w}-\mathbf{t}\|^2$ splňuje \textbf{normální rovnice}
\[ \mathbf{X}^{\!\top}\mathbf{X}\,\mathbf{w}=\mathbf{X}^{\!\top}\mathbf{t}. \]
Je-li matice $\mathbf{X}^{\!\top}\mathbf{X}$ (rozměru $D\times D$) regulární, má řešení tvar
\[ \boxed{\;\mathbf{w}=(\mathbf{X}^{\!\top}\mathbf{X})^{-1}\mathbf{X}^{\!\top}\mathbf{t}.\;} \]
\end{veta}

\begin{intuice}[Odvození přes nulový gradient (Fermat)]
Toto \emph{není} formální důkaz, jen výpočet -- ukázka, jak se ke vzorci dojde. Funkce $E$ je hladká a~konvexní, takže její minimum je tam, kde je gradient nulový (Fermatova věta o~vnitřním extrému, vícerozměrná verze: $\nabla_{\mathbf{w}}E=\mathbf{0}$). Krok po kroku:
\begin{enumerate}
  \item \textbf{Parciální derivace podle jedné váhy $w_j$.} Z~$E=\tfrac12\sum_i(\mathbf{x}_i^{\!\top}\mathbf{w}-t_i)^2$ řetízkovým pravidlem
  \[ \frac{\partial E}{\partial w_j}=\frac12\sum_{i=1}^N 2(\mathbf{x}_i^{\!\top}\mathbf{w}-t_i)\,x_{ij}=\sum_{i=1}^N x_{ij}(\mathbf{x}_i^{\!\top}\mathbf{w}-t_i). \]
  \item \textbf{Položíme rovno nule pro všechna $j$.} Suma $\sum_i x_{ij}(\dots)$ je skalární součin $j$-tého sloupce $\mathbf{X}$ s~vektorem $(\mathbf{X}\mathbf{w}-\mathbf{t})$, tj.\ $\mathbf{X}_{*,j}^{\!\top}(\mathbf{X}\mathbf{w}-\mathbf{t})=0$.
  \item \textbf{Sebereme rovnice přes všechna $j$ do matice.} Dostaneme $\mathbf{X}^{\!\top}(\mathbf{X}\mathbf{w}-\mathbf{t})=\mathbf{0}$, tj.\ $\mathbf{X}^{\!\top}\mathbf{X}\mathbf{w}=\mathbf{X}^{\!\top}\mathbf{t}$.
  \item \textbf{Invertujeme} (je-li $\mathbf{X}^{\!\top}\mathbf{X}$ regulární): $\mathbf{w}=(\mathbf{X}^{\!\top}\mathbf{X})^{-1}\mathbf{X}^{\!\top}\mathbf{t}$.
\end{enumerate}
Že jde o~minimum (ne maximum/sedlo), plyne z~konvexity $E$: její Hessián $\mathbf{X}^{\!\top}\mathbf{X}$ je pozitivně semidefinitní.
\end{intuice}

\begin{poznamka}[Singularita a SVD]
Matice $\mathbf{X}^{\!\top}\mathbf{X}$ je singulární (neinvertibilní), pokud sloupce $\mathbf{X}$ jsou lineárně závislé -- typicky když je příznaků víc než příkladů ($D>N$) nebo jsou některé příznaky kolineární. Pak normální rovnice nemá jednoznačné řešení. Místo přímé inverze se používá \textbf{SVD} (singulární rozklad) $\mathbf{X}$, který spolehlivě najde řešení (přes pseudoinverzi) i~v~singulárním a~numericky špatně podmíněném případě; je to i~standardní postup uvnitř knihovních funkcí (\texttt{numpy.linalg.lstsq}). Alternativně singularitu odstraní $L^2$ regularizace (níže).
\end{poznamka}

\begin{algo}[Analytické řešení lineární regrese]
\textbf{Vstup:} data $\mathbf{X}\in\mathbb{R}^{N\times D}$ (se sloupcem jedniček), cíle $\mathbf{t}\in\mathbb{R}^N$.\\
\textbf{Výstup:} váhy $\mathbf{w}$ minimalizující MSE.
\begin{enumerate}
  \item Sestav $\mathbf{X}^{\!\top}\mathbf{X}$ a~$\mathbf{X}^{\!\top}\mathbf{t}$. \cmt{$D\times D$ matice, $D$-vektor}
  \item Je-li $\mathbf{X}^{\!\top}\mathbf{X}$ regulární: $\mathbf{w}\leftarrow(\mathbf{X}^{\!\top}\mathbf{X})^{-1}\mathbf{X}^{\!\top}\mathbf{t}$. \cmt{jinak řeš přes SVD}
\end{enumerate}
Složitost $\mathcal{O}(N D^2)$ (sestavení $\mathbf{X}^{\!\top}\mathbf{X}$) za předpokladu $N\ge D$.
\end{algo}

\begin{priklad}[SZZ léto 2024 č.~31: Metoda nejmenších čtverců]
\szzotazka{léto 2024}{31}{Metoda nejmenších čtverců}\par
\textbf{Zadání.} Data s~jedním číselným atributem $x$ a~cílem $y$: $(1,1),(2,3),(3,2)$. (1)~Popište lineární model a~metodu nejmenších čtverců. (2)~Odhadněte koeficienty modelu. (3)~Spočtěte střední kvadratickou chybu modelu.

\smallskip
\textbf{Řešení.} (1)~Model je přímka $\hat y=\beta_0+\beta_1 x$ ($\beta_0$ posun, $\beta_1$ směrnice); koeficienty odhadneme tak, aby \emph{součet čtverců reziduí} $\sum_i(\beta_0+\beta_1 x_i-y_i)^2$ byl nejmenší (nulový gradient $\Rightarrow$ normální rovnice).

(2)~Pro jednu proměnnou stačí uzavřené vzorce s~průměry $\bar x=2$, $\bar y=2$:
\[ \beta_1=\frac{\sum_i(x_i-\bar x)(y_i-\bar y)}{\sum_i(x_i-\bar x)^2}=\frac{(-1)(-1)+0\cdot1+1\cdot0}{(-1)^2+0^2+1^2}=\frac{1}{2},\qquad \beta_0=\bar y-\beta_1\bar x=2-\tfrac12\cdot2=1. \]
Tedy $\boxed{\hat y=1+\tfrac12 x}$ (totéž dá maticově $(\mathbf{X}^{\!\top}\mathbf{X})^{-1}\mathbf{X}^{\!\top}\mathbf{t}$ se sloupcem jedniček).

(3)~Predikce $\hat y=(1{,}5;\,2;\,2{,}5)$, rezidua $y-\hat y=(-0{,}5;\,1;\,-0{,}5)$. Pak
\[ \mathrm{MSE}=\frac13\big((-0{,}5)^2+1^2+(-0{,}5)^2\big)=\frac{0{,}25+1+0{,}25}{3}=\frac{1{,}5}{3}=\boxed{0{,}5}. \]
\emph{(Ověřeno \texttt{numpy}: \texttt{lstsq} dá $(\beta_1,\beta_0)=(0{,}5;\,1)$, MSE $=0{,}5$, součet čtverců reziduí $=1{,}5$.)}

\smallskip
Geometricky: přímka prokládá tři body s~minimálním součtem čtverců svislých úseček (reziduí).
\begin{center}
\begin{tikzpicture}[x=2.0cm,y=1.6cm]
  \draw[osa] (0,0) -- (3.6,0) node[right,font=\scriptsize]{$x$};
  \draw[osa] (0,0) -- (0,3.4) node[above,font=\scriptsize]{$y$};
  \foreach \xv in {1,2,3}{\draw (\xv,0.03)--(\xv,-0.03) node[below,font=\tiny]{\xv};}
  \foreach \yv in {1,2,3}{\draw (0.02,\yv)--(-0.02,\yv) node[left,font=\tiny]{\yv};}
  % regresni primka y=1+0.5x na intervalu x in [0.4,3.4]
  \draw[hranice,blue!60!black] (0.4,1.2) -- (3.4,2.7);
  \node[blue!60!black,font=\scriptsize,anchor=south west] at (3.0,2.55) {$\hat y=1+\tfrac12 x$};
  % rezidua (svisle usecky od bodu k primce)
  \draw[red!70!black,thick] (1,1) -- (1,1.5);
  \draw[red!70!black,thick] (2,3) -- (2,2);
  \draw[red!70!black,thick] (3,2) -- (3,2.5);
  % popisek rezidua (k nejdelsimu, vlevo od usecky, mimo primku i body)
  \node[red!70!black,font=\scriptsize,anchor=east] at (1.9,2.6) {reziduum};
  % body
  \node[cls1] at (1,1) {}; \node[cls1] at (2,3) {}; \node[cls1] at (3,2) {};
  % predikce na primce (prazdne)
  \node[circle,draw=blue!60!black,fill=white,inner sep=0pt,minimum size=3pt] at (1,1.5) {};
  \node[circle,draw=blue!60!black,fill=white,inner sep=0pt,minimum size=3pt] at (2,2) {};
  \node[circle,draw=blue!60!black,fill=white,inner sep=0pt,minimum size=3pt] at (3,2.5) {};
\end{tikzpicture}
\obrpopis{Proložení nejmenšími čtverci dat $(1,1),(2,3),(3,2)$ přímkou $\hat y=1+\tfrac12 x$. Plné body jsou data, prázdné kroužky predikce na přímce, červené svislé úsečky jsou rezidua $y_i-\hat y_i=(-0{,}5;\,1;\,-0{,}5)$. Metoda minimalizuje součet jejich čtverců ($=1{,}5$).}
\end{center}
\end{priklad}

\begin{priklad}[SZZ podzim 2025 č.~30: Lineární regrese a~L2 (centrovaná data)]
\szzotazka{podzim 2025}{30}{Lineární regrese, L2 regularizace}\par
\textbf{Zadání.} Data $(x,y)$: $(-1,1),(0,3),(1,2)$. (1)~Popište model a~MNČ. (2)~Spočtěte OLS koeficienty (bez regularizace). (3)~Popište cílovou funkci s~$L^2$ a~určete \emph{směr} změny koeficientu u~$x$ oproti OLS.

\smallskip
\textbf{Řešení.} (1)~Model $f(x)=\beta_0+\beta_1 x$, MNČ minimalizuje $\sum_i(f(x_i)-y_i)^2$, řešení $\boldsymbol\beta=(\mathbf{X}^{\!\top}\mathbf{X})^{-1}\mathbf{X}^{\!\top}\mathbf{y}$ se sloupcem jedniček.

(2)~Data jsou \emph{centrovaná} ($\bar x=0$), takže $\mathbf{X}^{\!\top}\mathbf{X}$ vyjde diagonální a~výpočet je krátký. Se sloupcem jedniček $\mathbf{X}=\begin{psmallmatrix}1&-1\\1&0\\1&1\end{psmallmatrix}$:
\[ \mathbf{X}^{\!\top}\mathbf{X}=\begin{psmallmatrix}3&0\\0&2\end{psmallmatrix},\qquad \mathbf{X}^{\!\top}\mathbf{y}=\begin{psmallmatrix}1+3+2\\-1\cdot1+0\cdot3+1\cdot2\end{psmallmatrix}=\begin{psmallmatrix}6\\1\end{psmallmatrix},\qquad \boldsymbol\beta=\begin{psmallmatrix}3&0\\0&2\end{psmallmatrix}^{-1}\begin{psmallmatrix}6\\1\end{psmallmatrix}=\begin{psmallmatrix}\boxed{2}\\[1pt]\boxed{1/2}\end{psmallmatrix}. \]
Tedy $\beta_0=2$, $\beta_1=\tfrac12$, model $\hat y=2+\tfrac12 x$. \emph{(Ověřeno \texttt{numpy}: $\boldsymbol\beta=(2;\,0{,}5)$.)}

(3)~$L^2$ minimalizuje $\sum_i(f(x_i)-y_i)^2+\lambda\beta_1^2$ ($\lambda>0$, bias $\beta_0$ se nepenalizuje). Penalizace táhne $\beta_1$ k~nule, takže \emph{$|\beta_1|$ klesne} (přímka se zploští), znaménko zůstane kladné. Intercept $\beta_0$ se zde nezmění (data jsou centrovaná, $\bar x=0$); obecně se měnit může.
\end{priklad}

\subsection{Trénování pomocí SGD}

Analytické řešení vyžaduje invertovat $D\times D$ matici, což je pro velké $D$ pomalé a~paměťově náročné. Iterativní alternativa je \textbf{gradientní sestup}: opakovaně se posouváme proti gradientu chyby.

\begin{definice}[Gradientní sestup, SGD, minibatch]
Minimalizujeme-li chybu $E(\mathbf{w})$, \textbf{gradientní sestup} opakuje krok
\[ \mathbf{w}\leftarrow\mathbf{w}-\alpha\,\nabla_{\mathbf{w}}E(\mathbf{w}), \]
kde $\alpha>0$ je \textbf{rychlost učení} (learning rate, délka kroku). Podle toho, z~kolika příkladů gradient odhadujeme:
\begin{itemize}
  \item \textbf{(dávkový) gradientní sestup} -- z~\emph{celé} trénovací množiny (přesný gradient, drahý krok);
  \item \textbf{stochastický gradientní sestup (SGD)} -- z~\emph{jednoho náhodného} příkladu (nestranný, ale zašuměný odhad gradientu, levný krok);
  \item \textbf{minibatch SGD} -- kompromis: z~náhodné podmnožiny $\mathbb{B}$ o~$B$ příkladech.
\end{itemize}
Jeden průchod všemi trénovacími daty (rozdělenými na minibatche náhodnou permutací) se nazývá \textbf{epocha}.
\end{definice}

Pro lineární regresi s~chybou $\tfrac12(\mathbf{x}^{\!\top}\mathbf{w}-t)^2$ (a~volitelně $L^2$ členem $\tfrac\lambda2\|\mathbf{w}\|^2$) je gradient přes minibatch
\[ \nabla_{\mathbf{w}}E(\mathbf{w})\approx\frac{1}{|\mathbb{B}|}\sum_{i\in\mathbb{B}}\big(\mathbf{x}_i^{\!\top}\mathbf{w}-t_i\big)\mathbf{x}_i+\lambda\mathbf{w}. \]
Reziduum $(\mathbf{x}_i^{\!\top}\mathbf{w}-t_i)$ škáluje příspěvek vstupu $\mathbf{x}_i$; $L^2$ člen přidá $\lambda\mathbf{w}$ (tlak vah k~nule, \emph{weight decay}); složka odpovídající biasu se v~praxi obvykle nepenalizuje.

\begin{algo}[Lineární regrese minibatch SGD]
\textbf{Vstup:} data $(\mathbf{X},\mathbf{t})$, rychlost učení $\alpha>0$, síla $L^2$ $\lambda\ge0$.\\
\textbf{Výstup:} váhy $\mathbf{w}$ (přibližně) minimalizující regularizovanou MSE.
\begin{enumerate}
  \item $\mathbf{w}\leftarrow\mathbf{0}$ (nebo náhodně). \cmt{inicializace}
  \item Opakuj do konvergence (přes epochy):
  \begin{enumerate}
    \item navzorkuj minibatch indexů $\mathbb{B}$; \cmt{náhodná permutace $\to$ chunky velikosti $B$ = epocha}
    \item $\mathbf{w}\leftarrow\mathbf{w}-\alpha\,\dfrac{1}{|\mathbb{B}|}\displaystyle\sum_{i\in\mathbb{B}}\big(\mathbf{x}_i^{\!\top}\mathbf{w}-t_i\big)\mathbf{x}_i-\alpha\lambda\mathbf{w}.$ \cmt{krok proti gradientu}
  \end{enumerate}
\end{enumerate}
\end{algo}

\begin{poznamka}[Konvergence a normalizace]
Součet čtverců je \emph{konvexní} funkce vah, takže gradientní sestup s~dostatečně malým krokem konverguje ke globálnímu optimu (pro konvexní hladkou ztrátu; přesněji SGD konverguje skoro jistě, klesá-li krok dostatečně, ale ne moc rychle -- podmínky $\sum_i\alpha_i=\infty$, $\sum_i\alpha_i^2<\infty$). Příznaky v~různých škálách potřebují různé kroky; proto se před SGD obvykle \textbf{normalizují} -- min--max škálováním do $[0,1]$, nebo \emph{standardizací} $x'=\frac{x-\hat\mu}{\hat\sigma}$ (nulový průměr, jednotkový rozptyl). $L^2$ regularizaci a~výběr $\lambda$/stupně polynomu coby hyperparametrů řešíme obecně v~sekci~\ref{sec:preuceni}.
\end{poznamka}

\begin{priklad}[SZZ jaro 2024 č.~27: Krok SGD s~L2 regularizací]
\szzotazka{jaro 2024}{27}{Lineární regrese a~regularizace}\par
\textbf{Zadání.} Model $y=\beta_1 x_1+\beta_2 x_2+\beta_0$, data $\{(0,2;3),(2,3;9),(1,1;2),(2,0;2)\}$ (řádky $x_1,x_2;y$). Učící konstanta $\alpha=0{,}1$, $L^2$ síla $\lambda=0{,}5$. (1)~Definujte chybovou funkci (součet čtverců $+$ $L^2$). (2)~Z~odhadu $\beta_1=1,\beta_2=3,\beta_0=-1$ proveďte \emph{jeden} krok gradientní metody.

\smallskip
\textbf{Řešení.} (1)~$E=\sum_{i=1}^{N}(y_i-\hat y_i)^2+\lambda\sum_{j\ge1}\beta_j^2$, kde $\hat y_i=\beta_1 x_{i1}+\beta_2 x_{i2}+\beta_0$; bias $\beta_0$ se do regularizace nezahrnuje.

(2)~Gradient počítáme jako \emph{součet přes všechny čtyři vzory} (zde dávkový krok). Pro tuto formu chyby (\emph{bez} faktoru $\tfrac12$) je $\frac{\partial E}{\partial\beta_0}=\sum_i -2(y_i-\hat y_i)$, $\frac{\partial E}{\partial\beta_j}=\sum_i -2(y_i-\hat y_i)x_{ij}+2\lambda\beta_j$. Spočteme predikce a~rezidua $r_i=\hat y_i-y_i$:
\[
\begin{array}{c|cc|c|c|c}
i & x_1 & x_2 & y & \hat y=\beta_1x_1+\beta_2x_2+\beta_0 & r=\hat y-y\\\hline
1 & 0 & 2 & 3 & 0+6-1=5 & 2\\
2 & 2 & 3 & 9 & 2+9-1=10 & 1\\
3 & 1 & 1 & 2 & 1+3-1=3 & 1\\
4 & 2 & 0 & 2 & 2+0-1=1 & -1
\end{array}
\]
Derivace součtu čtverců $E_{\mathrm{RSS}}$ (nepenalizovaná část chyby; pozor: $-2(y-\hat y)=2r$, sčítáme příspěvky $2r_i$, resp.\ $2r_i x_{ij}$):
\[
\frac{\partial E_{\mathrm{RSS}}}{\partial\beta_0}=2(2+1+1-1)=6,\quad
\frac{\partial E_{\mathrm{RSS}}}{\partial\beta_1}=2(2\cdot0+1\cdot2+1\cdot1-1\cdot2)=2,\quad
\frac{\partial E_{\mathrm{RSS}}}{\partial\beta_2}=2(2\cdot2+1\cdot3+1\cdot1-1\cdot0)=16.
\]
Regularizační člen přidá $\frac{\partial}{\partial\beta_j}\lambda\beta_j^2=2\lambda\beta_j=2\cdot0{,}5\cdot\beta_j=\beta_j$ (jen k~vahám, ne k~$\beta_0$): $+\beta_1=1$, $+\beta_2=3$. Celkové gradienty:
\[ g_0=6,\qquad g_1=2+1=3,\qquad g_2=16+3=19. \]
Krok $\beta\leftarrow\beta-\alpha\,g$ s~$\alpha=0{,}1$:
\[ \boxed{\beta_0=-1-0{,}1\cdot6=-1{,}6,\quad \beta_1=1-0{,}1\cdot3=0{,}7,\quad \beta_2=3-0{,}1\cdot19=1{,}1.} \]
\emph{(Ověřeno \texttt{numpy}: $g=(6;\,3;\,19)$, nové $\beta=(-1{,}6;\,0{,}7;\,1{,}1)$ -- shoduje se s~oficiálním náčrtem.)}
\end{priklad}

\subsection{OLS vs. SGD a krátce o nelinearitě}

\begin{priklad}[SZZ léto 2025 č.~30: Lineární regrese -- analyticky vs.~SGD]
\szzotazka{léto 2025}{30}{Lineární regrese}\par
\textbf{Zadání (zkráceně).} (1)~Zformulujte model a~chybovou funkci. (2)~Jak získat optimum \emph{explicitně}? (3)~Zformulujte optimalizaci pomocí SGD. (4)~Výhody a~nevýhody obou, zejména z~hlediska kontroly chyby na validačních datech.

\smallskip
\textbf{Řešení.} (1)~$y(\mathbf{x})=\mathbf{x}^{\!\top}\mathbf{w}$ (bias paddingem 1), minimalizujeme součet čtverců $\tfrac12\|\mathbf{X}\mathbf{w}-\mathbf{t}\|^2$. (2)~\emph{Analyticky}: položíme gradient roven nule, dostaneme normální rovnice $\mathbf{X}^{\!\top}\mathbf{X}\mathbf{w}=\mathbf{X}^{\!\top}\mathbf{t}$ a~řešíme přes inverzi $\mathbf{w}=(\mathbf{X}^{\!\top}\mathbf{X})^{-1}\mathbf{X}^{\!\top}\mathbf{t}$, resp.\ přes SVD při singularitě. (3)~\emph{SGD}: opakuj $\mathbf{w}\leftarrow\mathbf{w}-\alpha\frac{1}{|\mathbb{B}|}\sum_{i\in\mathbb{B}}(\mathbf{x}_i^{\!\top}\mathbf{w}-t_i)\mathbf{x}_i$ po minibatchích/epochách.

(4)~\textbf{Analytické řešení:} přesné, v~jednom kroku, bez ladění; ale inverze $D\times D$ matice je drahá a~paměťově náročná pro velké $D$, a~přesné optimum na trénovacích datech může způsobit \emph{overfitting} (přeučení). \textbf{SGD:} škáluje na velká data a~málo paměti, snadno přidá regularizaci -- a~hlavně \emph{můžeme v~průběhu sledovat chybu na validační množině} a~zastavit, než nastane overfitting (\textbf{early stopping}); naopak vyžaduje ladit $\alpha$, počet epoch a~konverguje jen přibližně. Stručně: \emph{malé $D$ a~potřeba přesnosti $\to$ analyticky; velká data a~kontrola overfittingu $\to$ SGD}.
\end{priklad}

\begin{poznamka}[Polynomiální příznaky: most k~nelinearitě]
Lineární regrese modeluje jen přímky/nadroviny, ale „lineární`` se vztahuje k~\emph{vahám}, ne ke vstupu. Rozšíříme-li vstup $x$ o~mocniny, příznaky $(x^0,x^1,\dots,x^M)$, predikce $w_0+w_1x+\dots+w_Mx^M$ je polynom stupně $M$ -- model zůstává lineární ve~$\mathbf{w}$, takže platí všechny vzorce výše. Vyšší $M$ zvyšuje kapacitu (riziko overfittingu, viz sekce~\ref{sec:preuceni}). Obecněji se konstruují \emph{příznaky} (feature engineering), např.\ one-hot kódování kategorických hodnot.
\end{poznamka}

\section{Logistická regrese}\label{sec:logisticka}

Logistická regrese je přes svůj název \emph{klasifikační} model: lineární funkci vstupu protáhne \emph{logistickou (sigmoidní)} funkcí a~výstup čte jako \emph{pravděpodobnost} třídy. Je to jedna z~nejpoužívanějších klasifikačních technik a~zároveň nejjednodušší \emph{pravděpodobnostní} klasifikátor trénovaný \emph{maximum likelihood estimation} (MLE, maximální věrohodnost; sekce~\ref{sec:evaluace}). V~této sekci ji zavedeme pro binární klasifikaci, odvodíme její ztrátovou funkci i~gradient a~zobecníme ji softmaxem na více tříd.

\begin{poznamka}[Vztah k~S1]
Okruh S1 (sekce \emph{Strojové učení}) zmiňuje logistickou regresi přehledově jako „změkčený tvrdý práh`` s~pravděpodobnostním výstupem. Zde to \textbf{nezdvojujeme} -- jdeme do hloubky, kterou S1 přenechává S3: parametrizace přes sigmoid, ztráta jako negativní log-věrohodnost, odvození gradientu a~SGD, multiclass softmax.
\end{poznamka}

\subsection{Binární logistická regrese}

\subsubsection*{Od tvrdého prahu k~sigmoidě}

Lineární klasifikátor počítá skóre $\bar y(\mathbf{x};\mathbf{w})=\mathbf{w}^{\!\top}\mathbf{x}+b$ a~třídu přiřadí podle znaménka (prahu). \emph{Perceptron} (historicky nejstarší algoritmus, Rosenblatt 1958) používá \emph{tvrdý} práh $\sign(\bar y)$: vrací jen třídu $\{-1,+1\}$, na lineárně neseparabilních datech nekonverguje a~nedává pravděpodobnosti. Logistická regrese tvrdý práh \emph{změkčí} -- skóre $\bar y$ protáhne \textbf{sigmoidou} a~výstup interpretuje jako $P(C_1\mid\mathbf{x})$.

\begin{definice}[Sigmoidní (logistická) funkce]
\emph{Sigmoida} je funkce
\[ \sigma(z)=\frac{1}{1+e^{-z}}. \]
Má hodnoty v~otevřeném intervalu $(0,1)$, je rostoucí, $\sigma(0)=\tfrac12$, $\sigma(-z)=1-\sigma(z)$ (středová symetrie kolem bodu $(0,\tfrac12)$) a~asymptoticky $\sigma(z)\to 1$ pro $z\to+\infty$, $\sigma(z)\to 0$ pro $z\to-\infty$. Její derivace je
\[ \sigma'(z)=\sigma(z)\big(1-\sigma(z)\big),\qquad \sigma'(0)=\tfrac14. \]
\end{definice}

\begin{center}
\begin{tikzpicture}
\begin{axis}[width=10cm,height=5.4cm,
  xlabel={$z=\mathbf{w}^{\!\top}\mathbf{x}+b$},ylabel={$\sigma(z)$},
  xmin=-6,xmax=6,ymin=-0.08,ymax=1.12,
  xtick={-6,-4,-2,0,2,4,6},ytick={0,0.5,1},
  axis lines=middle,every axis plot/.append style={thick},clip=false,
  enlargelimits=false]
  % asymptoty
  \addplot[black!45,dashed,domain=-6:6,samples=2]{1};
  \addplot[black!45,dashed,domain=-6:6,samples=2]{0};
  \node[font=\scriptsize,black!55,anchor=south west] at (axis cs:-5.9,1) {asymptota $1$};
  \node[font=\scriptsize,black!55,anchor=south west] at (axis cs:-5.9,0.04) {asymptota $0$};
  % sigmoida
  \addplot[blue!60!black,domain=-6:6,samples=120]{1/(1+exp(-x))};
  % bod 0.5
  \addplot[only marks,mark=*,mark size=1.6pt,red!70!black] coordinates {(0,0.5)};
  \node[font=\scriptsize,red!70!black,anchor=west] at (axis cs:0.25,0.55) {$\sigma(0)=0{,}5$ (rozhodovací práh)};
\end{axis}
\end{tikzpicture}
\obrpopis{Sigmoida $\sigma(z)=1/(1+e^{-z})$. Hladce přechází z~$0$ do~$1$, hodnotu $0{,}5$ nabývá v~$z=0$ (rozhodovací práh) a~k~asymptotám $0$ a~$1$ se blíží pro velká záporná, resp.\ kladná $z$. Záporné $z$ tedy znamená pravděpodobnost třídy $C_1$ pod $50\,\%$, kladné nad $50\,\%$.}
\end{center}

\begin{definice}[Model binární logistické regrese]
Pro binární klasifikaci ($t\in\{0,1\}$) logistická regrese modeluje podmíněnou pravděpodobnost pozitivní třídy $C_1$ pomocí lineární funkce protažené sigmoidou. Označíme \emph{lineární část} $\bar y(\mathbf{x};\mathbf{w})=\mathbf{w}^{\!\top}\mathbf{x}$ a~celkový výstup
\[ y(\mathbf{x};\mathbf{w})=\sigma\big(\mathbf{w}^{\!\top}\mathbf{x}\big),\qquad
   P(C_1\mid\mathbf{x})=y(\mathbf{x};\mathbf{w}),\qquad
   P(C_0\mid\mathbf{x})=1-y(\mathbf{x};\mathbf{w}). \]
\textbf{Predikce.} Vstup zařadíme do třídy $C_1$, je-li $y(\mathbf{x};\mathbf{w})\ge 0{,}5$ (\emph{rozhodovací práh} $0{,}5$), jinak do $C_0$. Protože $\sigma(z)\ge\tfrac12 \iff z\ge 0$, je rozhodnutí ekvivalentní podmínce $\mathbf{w}^{\!\top}\mathbf{x}\ge 0$.
\end{definice}

\begin{intuice}[Bias jako práh; lineární rozhodovací hranice]
\emph{Bias} $b$ (absorbovaný do $\mathbf{w}$ přidáním stálého příznaku $x_0=1$, jak děláme dál) hraje roli trénovatelného prahu, proto stačí porovnávat skóre s~nulou. \textbf{Rozhodovací hranice} je množina bodů s~$y=0{,}5$, tj.\ $\mathbf{w}^{\!\top}\mathbf{x}=0$ -- to je \emph{nadrovina} (přímka v~2D), takže logistická regrese je \emph{lineární} klasifikátor. Sigmoida mění jen \emph{tvar} přechodu mezi třídami (jak rychle pravděpodobnost roste přes hranici), ne jeho \emph{polohu}.
\end{intuice}

\begin{poznamka}[Logit jako log-šance]
Inverze sigmoidy dává pro $p=P(C_1\mid\mathbf{x})$ lineární část jako \emph{logit}:
\[ \bar y(\mathbf{x};\mathbf{w})=\log\frac{p}{1-p}=\log\frac{P(C_1\mid\mathbf{x})}{P(C_0\mid\mathbf{x})}. \]
Lineární skóre tak má význam \emph{logaritmu šance} (log-odds) mezi oběma třídami: $\bar y>0$ znamená šanci ve prospěch $C_1$, $\bar y<0$ ve prospěch $C_0$.
\end{poznamka}

\subsection{Ztrátová funkce, gradient a~SGD}

\subsubsection*{Ztráta = negativní log-věrohodnost (křížová entropie)}

Logistická regrese se trénuje \emph{maximum likelihood estimation} (MLE) (sekce~\ref{sec:evaluace}). Výstup $y(\mathbf{x};\mathbf{w})$ je přímo pravděpodobnost třídy $C_1$, takže cíl $t$ pochází z~\emph{Bernoulliho} rozdělení s~parametrem $\varphi=y$: $p(t\mid\mathbf{x};\mathbf{w})=y^{\,t}(1-y)^{1-t}$. Minimalizujeme \emph{negativní logaritmickou věrohodnost} (NLL); pro minibatch $N$ příkladů je ztráta
\[ E(\mathbf{w})=\frac1N\sum_{i=1}^N -\log p\big(C_{t_i}\mid\mathbf{x}_i;\mathbf{w}\big)
   =\frac1N\sum_{i=1}^N \Big(-t_i\log y(\mathbf{x}_i;\mathbf{w})-(1-t_i)\log\big(1-y(\mathbf{x}_i;\mathbf{w})\big)\Big). \]

\begin{definice}[Binární křížová entropie]
Pro jeden příklad s~výstupem $y=\sigma(\mathbf{w}^{\!\top}\mathbf{x})$ a~cílem $t\in\{0,1\}$ je ztráta (binární křížová entropie, $=$ NLL Bernoulliho modelu)
\[ E = -t\log y-(1-t)\log(1-y). \]
\end{definice}

\begin{intuice}
Pro $t=1$ se z~ní stane $-\log y$ (chceme $y\to 1$), pro $t=0$ je to $-\log(1-y)$ (chceme $y\to 0$). Jistá \emph{a~správná} predikce dá ztrátu $0$; jistá \emph{a~chybná} predikce ($y\to 0$ při $t=1$) ztrátu $+\infty$ -- model je za sebevědomý omyl tvrdě penalizován. Název „křížová entropie`` plyne z~obecného rámce MLE: NLL $=$ křížová entropie mezi empirickým rozdělením cílů a~modelem (sekce~\ref{sec:evaluace}). Pozn.: kvadratická (MSE) ztráta se zde \emph{nepoužívá} -- vede k~nekonvexnímu problému a~mizejícím gradientům v~saturaci sigmoidy.
\end{intuice}

\subsubsection*{Odvození gradientu}

Toto odvození žádá otázka SZZ jaro 2025 (č.~29, část~2). Provedeme derivaci krok po kroku pro \emph{jeden} příklad; značme $z=\mathbf{w}^{\!\top}\mathbf{x}$, takže $y=\sigma(z)$ a~$E=-t\log\sigma(z)-(1-t)\log\big(1-\sigma(z)\big)$. Cílem je $\nabla_{\mathbf{w}}E$.

\begin{veta}[Gradient ztráty logistické regrese]
Pro jeden příklad $(\mathbf{x},t)$ a~ztrátu $E=-t\log\sigma(\mathbf{w}^{\!\top}\mathbf{x})-(1-t)\log\big(1-\sigma(\mathbf{w}^{\!\top}\mathbf{x})\big)$ platí
\[ \frac{\partial E}{\partial z}=\sigma(z)-t,\qquad\text{a~tedy}\qquad
   \boxed{\ \nabla_{\mathbf{w}}E=\big(\sigma(\mathbf{w}^{\!\top}\mathbf{x})-t\big)\,\mathbf{x}\ }. \]
\end{veta}

\begin{proof}[Odvození \textnormal{(krok po kroku)}]
\hfill\par
\begin{enumerate}[label=(\arabic*)]
  \item \textbf{Derivace sigmoidy.} Z~$\sigma(z)=\big(1+e^{-z}\big)^{-1}$ řetízkovým pravidlem
    \[ \sigma'(z)=-\big(1+e^{-z}\big)^{-2}\cdot(-e^{-z})=\frac{e^{-z}}{(1+e^{-z})^2}
       =\frac{1}{1+e^{-z}}\cdot\frac{e^{-z}}{1+e^{-z}}=\sigma(z)\big(1-\sigma(z)\big), \]
    kde poslední rovnost využívá $\dfrac{e^{-z}}{1+e^{-z}}=\dfrac{(1+e^{-z})-1}{1+e^{-z}}=1-\sigma(z)$.
  \item \textbf{Derivace ztráty podle $z$.} Použijeme $\dfrac{\partial}{\partial z}\log\sigma=\dfrac{\sigma'}{\sigma}$ a~$\dfrac{\partial}{\partial z}\log(1-\sigma)=\dfrac{-\sigma'}{1-\sigma}$:
    \[ \frac{\partial E}{\partial z}
       = -t\cdot\frac{\sigma'(z)}{\sigma(z)}-(1-t)\cdot\frac{-\sigma'(z)}{1-\sigma(z)}
       = -t\,\frac{\sigma'(z)}{\sigma(z)}+(1-t)\,\frac{\sigma'(z)}{1-\sigma(z)}. \]
  \item \textbf{Dosazení $\sigma'=\sigma(1-\sigma)$.} V~prvním zlomku se zkrátí $\sigma$, ve druhém $(1-\sigma)$:
    \[ \frac{\partial E}{\partial z}
       = -t\,\big(1-\sigma(z)\big)+(1-t)\,\sigma(z). \]
  \item \textbf{Roznásobení a~zjednodušení.}
    \[ \frac{\partial E}{\partial z}
       = -t+t\,\sigma(z)+\sigma(z)-t\,\sigma(z)
       = \sigma(z)-t. \]
    Členy $\pm t\,\sigma(z)$ se vyruší; zbývá rozdíl predikce a~cíle.
  \item \textbf{Přechod ke gradientu podle $\mathbf{w}$.} Protože $z=\mathbf{w}^{\!\top}\mathbf{x}$, je $\nabla_{\mathbf{w}}z=\mathbf{x}$. Řetízkovým pravidlem
    \[ \nabla_{\mathbf{w}}E=\frac{\partial E}{\partial z}\,\nabla_{\mathbf{w}}z=\big(\sigma(z)-t\big)\,\mathbf{x}=\big(y(\mathbf{x};\mathbf{w})-t\big)\,\mathbf{x}. \qedhere \]
\end{enumerate}
\end{proof}

\begin{intuice}[Elegantní tvar a~jednotnost s~lineární regresí]
Gradient má pozoruhodně jednoduchý tvar \emph{(chyba predikce)} $\times$ \emph{(vstup)}: $\big(y-t\big)\mathbf{x}$. Je-li predikce přesná ($y=t$), gradient je nulový a~váhy se nemění; čím větší chyba, tím větší krok, a~to ve směru příznaků $\mathbf{x}$. \emph{Stejný} tvar $\big(y-t\big)\mathbf{x}$ má gradient \emph{lineární} regrese s~MSE ztrátou (sekce~\ref{sec:linreg}) i~\emph{softmaxové} regrese (níže) -- to není náhoda, ale důsledek \emph{zobecněných lineárních modelů}: vhodná aktivace (identita / sigmoida / softmax) s~odpovídajícím rozdělením (normální / Bernoulli / kategorické) a~NLL ztrátou vždy vede k~témuž tvaru gradientu.
\end{intuice}

\subsubsection*{Stochastický gradientní sestup}

Minimum ztráty nemá logistická regrese (na rozdíl od lineární) v~uzavřeném tvaru, hledáme ho proto \emph{gradientním sestupem}. V~praxi \emph{stochastickým} (SGD): místo plného gradientu přes všechna data použijeme odhad z~jednoho příkladu nebo malé dávky (minibatch).

\begin{algo}[Trénování logistické regrese (minibatch SGD)]
\textbf{Vstup:} data $(\mathbf{X}\in\mathbb{R}^{N\times D},\ \mathbf{t}\in\{0,1\}^N)$, rychlost učení $\alpha>0$.\par
\textbf{Výstup:} váhy $\mathbf{w}$.
\begin{enumerate}[label=\arabic*.,leftmargin=2.2em]
  \item $\mathbf{w}\leftarrow\mathbf{0}$ (nebo malé náhodné hodnoty) \cmt{inicializace}
  \item \textbf{opakuj} do konvergence (nebo vyčerpání trpělivosti) pro minibatch $\mathbb{B}$:
  \begin{enumerate}[label=\alph*),leftmargin=2.2em]
    \item $\displaystyle \mathbf{g}\leftarrow\frac{1}{|\mathbb{B}|}\sum_{i\in\mathbb{B}}\big(\sigma(\mathbf{w}^{\!\top}\mathbf{x}_i)-t_i\big)\,\mathbf{x}_i$ \cmt{průměrný gradient na dávce}
    \item $\mathbf{w}\leftarrow\mathbf{w}-\alpha\,\mathbf{g}$ \cmt{krok proti gradientu}
  \end{enumerate}
\end{enumerate}
\end{algo}

\noindent Pro jeden příklad (čistý SGD, $|\mathbb{B}|=1$) je krok přímo
\[ \mathbf{w}\leftarrow\mathbf{w}-\alpha\big(\sigma(\mathbf{w}^{\!\top}\mathbf{x})-t\big)\,\mathbf{x}. \]

\begin{poznamka}[Prevence overfittingu]
Vše, co platí o~\emph{příznacích} a~$L^2$ \emph{regularizaci} obecně, platí i~zde (sekce~\ref{sec:preuceni}). Overfittingu (přeučení) při SGD bráníme zejména:
\begin{itemize}
  \item \textbf{$L^2$ regularizací} -- ke ztrátě přidáme $\tfrac{\lambda}{2}\|\mathbf{w}\|^2$, čímž do kroku přibude člen $-\alpha\lambda\mathbf{w}$ (váhy se v~každém kroku mírně „smršťují`` k~nule);
  \item \textbf{early stopping} (včasné zastavení) -- sledujeme \emph{validační} chybu a~trénování ukončíme, jakmile začne růst (model se přestává učit signál a~začíná se učit šum).
\end{itemize}
Pomáhá i~rozumný \emph{learning-rate schedule} (postupné snižování $\alpha$). Podrobnosti a~obrázek validační vs.\ trénovací chyby viz sekce~\ref{sec:preuceni}.
\end{poznamka}

\subsection{Klasifikace do více tříd: softmax}

Pro $K>2$ tříd zobecníme model dvěma kroky: vygenerujeme $K$ lineárních výstupů (každý se svými vahami) a~sigmoidu nahradíme \textbf{softmaxem}, který $K$ čísel převede na rozdělení pravděpodobností přes třídy.

\begin{definice}[Multinomiální logistická regrese (softmax)]
Pro $K$ tříd máme matici vah $\mathbf{W}\in\mathbb{R}^{D\times K}$ (sloupec na třídu). Lineární část je vektor $\bar{\mathbf{y}}(\mathbf{x};\mathbf{W})=\mathbf{x}^{\!\top}\mathbf{W}\in\mathbb{R}^K$ (jeho $i$-tá složka je skóre třídy $C_i$). \emph{Softmax} z~něj dělá pravděpodobnosti:
\[ P(C_i\mid\mathbf{x};\mathbf{W})=\operatorname{softmax}\big(\bar{\mathbf{y}}\big)_i=\frac{e^{\bar y_i}}{\sum_{j=0}^{K-1} e^{\bar y_j}}. \]
Výsledné hodnoty jsou v~$(0,1)$ a~sčítají se na~$1$. Klasifikujeme do třídy s~\emph{nejvyšší} pravděpodobností, $\hat c=\argmax_i P(C_i\mid\mathbf{x};\mathbf{W})$. Model se též nazývá \emph{softmax regrese} nebo \emph{maximum entropy} klasifikátor.
\end{definice}

\begin{intuice}[Softmax zobecňuje sigmoidu]
Softmax je „měkký argmax``: zvýrazní největší skóre, ale zachová pravděpodobnostní rozdělení (nejen tvrdé vítězství). Pro $K=2$ se redukuje na sigmoidu, $\sigma(z)=\operatorname{softmax}\big([z,0]\big)_0=\dfrac{e^{z}}{e^{z}+e^{0}}=\dfrac{1}{1+e^{-z}}$. Exponenciála je vždy kladná, dělení normalizuje na součet~$1$. Softmax je invariantní vůči přičtení konstanty ke všem skóre ($\operatorname{softmax}(\bar{\mathbf{y}}+c)=\operatorname{softmax}(\bar{\mathbf{y}})$), takže model je mírně \emph{přeurčený} -- pro $K$ tříd by stačilo $K-1$ sad vah.
\end{intuice}

\begin{poznamka}[Ztráta a~trénování -- analogicky binárnímu případu]
Cíl $t\in\{0,1,\dots,K-1\}$ reprezentujeme \emph{one-hot} vektorem $\mathbf{1}_t$ ($1$ na pozici správné třídy, jinak $0$). Cíl pochází z~\emph{kategorického} rozdělení, ztráta je opět NLL $=$ (kategorická) křížová entropie:
\[ E(\mathbf{W})=\frac1N\sum_{i=1}^N -\log P\big(C_{t_i}\mid\mathbf{x}_i;\mathbf{W}\big). \]
Trénuje se stejným minibatch SGD jako binární případ. Gradient ztráty podle vah třídy má opět tvar \emph{(chyba) $\times$ (vstup)},
\[ \nabla_{\mathbf{W}}E=\mathbf{x}\,\big(\mathbf{y}(\mathbf{x})-\mathbf{1}_t\big)^{\!\top}\in\mathbb{R}^{D\times K}\quad\text{(pro jeden příklad; sloupec třídy $i$ je $(y_i-[t{=}i])\,\mathbf{x}$)}, \]
kde $\mathbf{y}(\mathbf{x})=\operatorname{softmax}(\mathbf{x}^{\!\top}\mathbf{W})$ je vektor predikovaných pravděpodobností. Odvození je analogické binárnímu případu: derivace křížové entropie podle skóre softmaxu vyjde $\partial E/\partial\bar{\mathbf{y}}=\mathbf{y}(\mathbf{x})-\mathbf{1}_t$. Rozhodovací oblasti softmaxu jsou \emph{konvexní} (a~tedy souvislé): hranice mezi dvojicemi tříd jsou lineární.
\end{poznamka}

\subsection{Řešené příklady}

\begin{priklad}[SZZ léto 2022 č.~10: Predikční funkce logistické regrese]
\szzotazka{léto 2022}{10}{Logistická regrese}\par
\textbf{Zadání.} \textbf{(1)}~Napište funkci, kterou predikuje model logistické regrese. \textbf{(2)}~Jakých hodnot tato funkce nabývá a~jaká je jejich interpretace?

\smallskip
\textbf{Řešení.} \textbf{(1)}~Model spočítá lineární kombinaci příznaků $z=\mathbf{w}^{\!\top}\mathbf{x}+b$ a~aplikuje na ni logistickou (sigmoidní) funkci $\sigma(z)=\tfrac{1}{1+e^{-z}}$, takže predikuje
\[ \hat y=\sigma\!\big(\mathbf{w}^{\!\top}\mathbf{x}+b\big)=\frac{1}{1+e^{-(\mathbf{w}^{\!\top}\mathbf{x}+b)}}. \]
\textbf{(2)}~Hodnoty leží v~intervalu $(0,1)$ a~interpretují se jako pravděpodobnost příslušnosti k~pozitivní třídě, $P(y{=}1\mid\mathbf{x})$. Vlastní klasifikaci dostaneme prahováním (typicky práh $0{,}5$: třída $1$ pro $\hat y\ge 0{,}5$, jinak třída $0$).
\end{priklad}

\begin{priklad}[SZZ jaro 2025 č.~29: Logistická regrese pro binární klasifikaci]
\szzotazka{jaro 2025}{29}{Logistická regrese pro binární klasifikaci}\par
\textbf{Zadání.} Dataset $\mathbf{X}=(\mathbf{x}_1,\dots,\mathbf{x}_N)$ s~cíli $\mathbf{t}=(t_1,\dots,t_N)$. \textbf{(1)}~Jak použít logistickou regresi pro binární klasifikaci? Jak vypadá predikce? Jak nezávisle odhadnout chybu na neviděných datech? \textbf{(2)}~Jak je definována ztrátová funkce? \emph{Odvoďte} její derivaci a~ukažte SGD. \textbf{(3)}~Jak při SGD zabráníte overfittingu?

\smallskip
\textbf{Řešení.}

\textbf{(1)~Model a~predikce.} Logistická regrese modeluje $P(C_1\mid\mathbf{x})=y(\mathbf{x};\mathbf{w})=\sigma(\mathbf{w}^{\!\top}\mathbf{x})$ s~$\sigma(z)=1/(1+e^{-z})$ a~$P(C_0\mid\mathbf{x})=1-y$. \emph{Predikce:} třída $1$, je-li $y\ge 0{,}5$ (ekvivalentně $\mathbf{w}^{\!\top}\mathbf{x}\ge 0$), jinak třída $0$. \emph{Odhad chyby na neviděných datech:} data rozdělíme na \emph{trénovací} (učení vah), \emph{validační} (ladění hyperparametrů, např.\ $\lambda$, rozhodnutí o~early stopping) a~\emph{testovací} (jediné finální měření). Při málu dat použijeme $k$-násobnou křížovou validaci (sekce~\ref{sec:evaluace}).

\textbf{(2)~Ztráta, derivace, SGD.} Ztráta je negativní log-věrohodnost (Bernoulli), pro minibatch
\[ E(\mathbf{w})=\frac1N\sum_i -\log p\big(C_{t_i}\mid\mathbf{x}_i;\mathbf{w}\big),\qquad
   \text{pro 1 vzor}\ E=-t\log\sigma(z)-(1-t)\log\big(1-\sigma(z)\big),\ z=\mathbf{w}^{\!\top}\mathbf{x}. \]
Derivaci provedeme přes $\sigma'(z)=\sigma(z)\big(1-\sigma(z)\big)$:
\[
\frac{\partial E}{\partial z}
= -t\frac{\sigma'}{\sigma}+(1-t)\frac{\sigma'}{1-\sigma}
= -t(1-\sigma)+(1-t)\sigma
= -t+t\sigma+\sigma-t\sigma
= \sigma(z)-t,
\]
a~protože $\nabla_{\mathbf{w}}z=\mathbf{x}$, je $\nabla_{\mathbf{w}}E=\big(\sigma(\mathbf{w}^{\!\top}\mathbf{x})-t\big)\mathbf{x}$ (plné odvození viz věta výše). \emph{SGD:} inicializuj $\mathbf{w}$, opakuj $\mathbf{w}\leftarrow\mathbf{w}-\alpha\big(\sigma(\mathbf{w}^{\!\top}\mathbf{x}_i)-t_i\big)\mathbf{x}_i$ (resp.\ průměr přes minibatch).

\textbf{(3)~Prevence overfittingu.} $L^2$ regularizace (člen $\tfrac{\lambda}{2}\|\mathbf{w}\|^2$ ke ztrátě, ve kroku $-\alpha\lambda\mathbf{w}$), \emph{early stopping} (zastavit, když roste validační chyba) a~rozumný learning-rate schedule (sekce~\ref{sec:preuceni}).
\end{priklad}

\begin{priklad}[SZZ podzim 2024 č.~29: Logistická regrese (model, výstup, multiclass)]
\szzotazka{podzim 2024}{29}{Logistická regrese}\par
\textbf{Zadání.} Model má předpovídat pravděpodobnost úspěchu u~zkoušky podle doby učení; data (úspěch $1$ / neúspěch $0$ vs.\ doba učení) jsou na obrázku. \textbf{(1)}~Popište model logistické regrese: jak se spočítá výstup, jak se trénují parametry, jaká je ztrátová funkce? \textbf{(2)}~Zakreslete přibližně výstup natrénovaného modelu (nepočítejte). \textbf{(3)}~Jak by se model lišil při klasifikaci do více než dvou tříd (např.\ predikce známky)?

\smallskip
\textbf{Řešení.}

\textbf{(1)~Model.} Jeden příznak $x=$ doba učení (plus bias). Výstup $y(x;\mathbf{w})=\sigma(w_1 x+w_0)$, čtený jako $P(\text{úspěch}\mid x)$; klasifikace prahem $0{,}5$. \emph{Trénink:} maximum likelihood estimation (MLE) přes SGD, minimalizujeme \emph{negativní log-věrohodnost} (NLL, binární křížovou entropii) $E=-t\log y-(1-t)\log(1-y)$ s~gradientem $\big(y-t\big)\,\mathbf{x}$; krok $\mathbf{w}\leftarrow\mathbf{w}-\alpha\,\mathbf{g}$.

\textbf{(2)~Výstup modelu.} Sigmoidní křivka rostoucí s~dobou učení: u~krátké doby blízko~$0$, u~dlouhé blízko~$1$, hodnotu $0{,}5$ (rozhodovací práh) protíná zhruba tam, kde se v~datech mísí neúspěchy a~úspěchy.

\begin{center}
\begin{tikzpicture}
\begin{axis}[width=10cm,height=5.2cm,
  xlabel={doba učení [h]},ylabel={$P(\text{úspěch})$},
  xmin=0,xmax=10,ymin=-0.12,ymax=1.15,
  xtick={0,2,4,6,8,10},ytick={0,0.5,1},
  axis lines=left,every axis plot/.append style={thick},clip=false,
  enlargelimits=false]
  % prahova cara 0.5
  \addplot[black!40,dashed,domain=0:10,samples=2]{0.5};
  \node[font=\scriptsize,black!55,anchor=south east] at (axis cs:10,0.5) {práh $0{,}5$};
  % sigmoidni vystup: stred okolo x=5, sklon 1.1
  \addplot[blue!60!black,domain=0:10,samples=120]{1/(1+exp(-1.1*(x-5)))};
  \node[font=\scriptsize,blue!60!black,anchor=west] at (axis cs:6.2,0.78) {výstup modelu};
  % data: neuspechy (t=0) pri kratke dobe, uspechy (t=1) pri dlouhe
  \addplot[only marks,mark=*,mark size=2pt,red!70!black] coordinates {(1,0)(1.8,0)(2.4,0)(3,0)(3.6,0)(4.4,0)(5.2,0)};
  \addplot[only marks,mark=*,mark size=2pt,green!45!black] coordinates {(4.8,1)(5.6,1)(6.2,1)(7,1)(7.6,1)(8.4,1)(9.2,1)};
  \node[font=\scriptsize,red!70!black,anchor=south west] at (axis cs:0.2,0.07) {neúspěch ($t=0$)};
  \node[font=\scriptsize,green!45!black,anchor=south west] at (axis cs:0.2,1.04) {úspěch ($t=1$)};
\end{axis}
\end{tikzpicture}
\obrpopis{Logistická regrese na 1D datech (doba učení $\to$ úspěch). Body leží na $y=0$ (neúspěchy, krátká doba) a~$y=1$ (úspěchy, dlouhá doba); modrá sigmoida je výstup modelu $P(\text{úspěch}\mid x)$. Práh $0{,}5$ protíná v~pásmu, kde se obě třídy překrývají -- to je rozhodovací hranice. Data jsou schematická.}
\end{center}

\textbf{(3)~Více tříd.} Místo jednoho výstupu se sigmoidou bychom měli $K$ výstupů (sada vah na třídu) a~sigmoidu nahradili \emph{softmaxem} $P(C_i\mid\mathbf{x})=e^{\bar y_i}/\sum_j e^{\bar y_j}$, který dá rozdělení pravděpodobností přes $K$ tříd (sčítá se na~$1$). Predikce $=\argmax_i P(C_i\mid\mathbf{x})$, ztráta $=$ kategorická křížová entropie, trénink opět SGD. (Pro predikci \emph{známky} by navíc dávalo smysl využít její \emph{uspořádanost} -- ordinální regrese -- ale to už je nad rámec.)
\end{priklad}

\section{Rozhodovací stromy}\label{sec:stromy}

\textbf{Rozhodovací strom} (decision tree) je model, který rozdělí prostor vstupů na \emph{obdélníkové oblasti} (osově zarovnané řezy) a~v~každé oblasti predikuje jednu hodnotu. Inference je triviální: od kořene následujeme větvící pravidla v~uzlech, až dorazíme do \emph{listu}, jehož hodnotu vrátíme. Stromy zvládají klasifikaci i~regresi, mix číselných i~kategorických příznaků, nevyžadují škálování dat a~jsou dobře \emph{interpretovatelné} -- proto jsou populární jako stavební kámen ansámblů (random forest, boosting; viz sekce~\ref{sec:kombinace}).

\begin{poznamka}[Vztah k~S1]
Okruh S1 (sekce \emph{Strojové učení}) už rozhodovací stromy zavádí na úrovni „základní principy``: definice (vnitřní uzly testují atribut, listy dávají rozhodnutí), hladová konstrukce „rozděl a~panuj`` a~kritérium \emph{informačního zisku} s~booleovskou entropií $B(q)$ a~příkladem AIMA (restaurace, Patrons vs.\ Type). Zde to \textbf{nezdvojujeme} -- entropii a~zisk jen krátce zrekapitulujeme s~konzistentním značením a~jdeme do hloubky, kterou S1 výslovně přenechává S3: \emph{Gini index} jako alternativní kritérium, \emph{regresní stromy} (rozptylové/MSE kritérium), \emph{souvislost kritérií se ztrátovou funkcí} a~především \textbf{prořezávání} (pruning).
\end{poznamka}

\subsection{Algoritmus učení: hladově shora dolů (rozděl a~panuj)}

Nalézt globálně nejmenší strom konzistentní s~daty je NP-těžké; v~praxi se proto staví \textbf{hladově}: začneme jediným listem se všemi příklady a~opakovaně \emph{rozdělíme} list na rozhodovací uzel se dvěma (obecně více) potomky tak, aby rozdělení co nejvíc \emph{zlepšilo homogenitu} skupin. Hladová volba se nikdy nereviduje (greedy). Označme $I_{\mathcal{T}}$ množinu indexů trénovacích příkladů v~uzlu $\mathcal{T}$.

\begin{definice}[Kritérium uzlu a~větvení]
\textbf{Kritérium} $c(\mathcal{T})$ měří \emph{nehomogenitu} (nečistotu, impurity) příkladů v~uzlu $\mathcal{T}$ -- kolik nejistoty o~cíli v~uzlu zbývá. Při rozdělení uzlu $\mathcal{T}$ na potomky $\mathcal{T}_L$ a~$\mathcal{T}_R$ hledáme \emph{příznak} a~jeho \emph{práh} (resp.\ hodnotu), které nejvíc sníží celkové kritérium, tj.\ minimalizují
\[ c(\mathcal{T}_L) + c(\mathcal{T}_R) - c(\mathcal{T}) \le 0. \]
Kritéria se obvykle volí \emph{vážená velikostí uzlu}, $c(\mathcal{T}) = |I_{\mathcal{T}}|\cdot(\text{nečistota})$, takže rozdíl výše je přímo \emph{informační zisk} (až na znaménko a~násobek $|I_{\mathcal{T}}|$).
\end{definice}

\begin{algo}[Učení rozhodovacího stromu (greedy, rozděl a~panuj)]
\ttfamily\small
BUILD-TREE(uzel $\mathcal{T}$ s~příklady $I_{\mathcal{T}}$):\\
\hspace*{1.5em}\textbf{if} zastavovací podmínka($\mathcal{T}$): \textbf{return} list s~predikcí z~$I_{\mathcal{T}}$ \cmt{čistý uzel / hloubka / min.\ příkladů}\\
\hspace*{1.5em}$(j^{*}, v^{*}) \leftarrow \argmin_{j,\,v}\;\big[\,c(\mathcal{T}_L^{j,v}) + c(\mathcal{T}_R^{j,v})\,\big]$ \cmt{for přes příznaky $j$ a~jejich hodnoty $v$}\\
\hspace*{1.5em}rozděl $I_{\mathcal{T}}$ testem „$x_{j^{*}} \le v^{*}$`` na $\mathcal{T}_L$, $\mathcal{T}_R$\\
\hspace*{1.5em}BUILD-TREE($\mathcal{T}_L$); BUILD-TREE($\mathcal{T}_R$)\\
\hspace*{1.5em}\textbf{return} rozhodovací uzel s~testem $(j^{*}, v^{*})$ a~podstromy $\mathcal{T}_L$, $\mathcal{T}_R$
\end{algo}

\begin{intuice}
Hledání nejlepšího řezu je dvojitý cyklus: \emph{přes všechny příznaky} a~\emph{přes všechny prahy} (u~číselného příznaku stačí hodnoty mezi sousedními setříděnými hodnotami, u~kategorického jednotlivé hodnoty). Pro každý kandidát spočítáme kritérium obou potomků a~vybereme řez s~nejnižším součtem. Predikce listu je \emph{nejčastější třída} (klasifikace), \emph{rozdělení tříd} (pravděpodobnostní výstup) nebo \emph{průměr cílů} (regrese). Rozhodovací uzly tak postupně \emph{dělí prostor osově zarovnanými řezy} na obdélníky (viz obrázek níže).
\end{intuice}

\begin{poznamka}[Zastavovací heuristiky a~způsob stavby]
Bez omezení strom roste, dokud nejsou listy čisté (overfitting (přeučení) -- viz prořezávání). Běžné \emph{heuristiky}: maximální \emph{hloubka}; minimální počet příkladů nutný k~dělení uzlu; maximální počet listů. Strom se obvykle staví buď \emph{do hloubky} (rekurzivně dělíme každý list, dokud platí omezení), nebo -- je-li dán limit počtu listů -- dělíme vždy ten list, který nabízí největší pokles kritéria. Strom \emph{neomezené} velikosti je neparametrický model (roste s~daty); s~pevným limitem listů je parametrický.
\end{poznamka}

\begin{center}
\begin{tikzpicture}[scale=1.18]
% --- levý panel: 2D prostor rozdělený osově zarovnanými řezy ---
\begin{scope}
  \def\S{3.6}
  % osy
  \draw[osa] (0,0) -- (\S+0.4,0) node[right,font=\small]{$x_1$};
  \draw[osa] (0,0) -- (0,\S+0.4) node[above,font=\small]{$x_2$};
  % oblasti: rez x1<=1.8 (svisle); vpravo rez x2<=2.2 (vodorovne)
  % R1: x1<=1.8                -> trida 1 (modra)
  % R2: x1>1.8, x2<=2.2        -> trida 2 (cervena)
  % R3: x1>1.8, x2>2.2         -> trida 3 (zelena)
  \fill[blue!12]  (0,0) rectangle (1.8,\S);
  \fill[red!12]   (1.8,0) rectangle (\S,2.2);
  \fill[green!12] (1.8,2.2) rectangle (\S,\S);
  % rezy
  \draw[hranice] (1.8,0) -- (1.8,\S);
  \draw[hranice] (1.8,2.2) -- (\S,2.2);
  \draw (0,0) rectangle (\S,\S);
  % body
  \node[cls1] at (0.7,1.1) {}; \node[cls1] at (1.2,2.6) {}; \node[cls1] at (0.5,3.0) {};
  \node[cls2] at (2.6,0.8) {}; \node[cls2] at (3.1,1.4) {}; \node[cls2] at (2.3,1.7) {};
  \node[cls3] at (2.7,3.0) {}; \node[cls3] at (3.2,2.7) {}; \node[cls3] at (2.4,3.2) {};
  % popisky oblasti
  \node[font=\scriptsize,blue!55!black]  at (0.9,0.4) {$R_1$};
  \node[font=\scriptsize,red!65!black]   at (2.9,0.4) {$R_2$};
  \node[font=\scriptsize,green!45!black] at (2.9,3.4) {$R_3$};
  % popisky rezu
  \node[font=\scriptsize] at (1.8,-0.3) {$x_1{=}1{,}8$};
  \node[font=\scriptsize,anchor=west] at (\S+0.05,2.2) {$x_2{=}2{,}2$};
\end{scope}
% --- pravy panel: odpovidajici strom ---
\begin{scope}[xshift=6.4cm,yshift=0.2cm]
  \node[treenode] (r) at (1.4,3.4) {$x_1 \le 1{,}8$};
  \node[treeleaf,fill=blue!14] (l1) at (0,1.9) {$R_1$};
  \node[treenode] (n2) at (2.6,1.9) {$x_2 \le 2{,}2$};
  \node[treeleaf,fill=red!14]   (l2) at (1.7,0.4) {$R_2$};
  \node[treeleaf,fill=green!14] (l3) at (3.5,0.4) {$R_3$};
  \draw[vetev] (r) -- node[popis,above left]{ano} (l1);
  \draw[vetev] (r) -- node[popis,above right]{ne} (n2);
  \draw[vetev] (n2) -- node[popis,above left]{ano} (l2);
  \draw[vetev] (n2) -- node[popis,above right]{ne} (l3);
\end{scope}
\end{tikzpicture}
\obrpopis{Rozhodovací strom dělí prostor příznaků \emph{osově zarovnanými řezy} na obdélníkové oblasti, každá oblast je jeden list (jedna predikovaná třída). Vlevo dva řezy ($x_1{=}1{,}8$ svisle, pak $x_2{=}2{,}2$ vodorovně vpravo) rozdělují rovinu na tři oblasti $R_1,R_2,R_3$; vpravo strom, který tyto řezy generuje. Hranice je tedy vždy po částech kolmá na osy -- strom neumí šikmou hranici jedním řezem (potřeboval by „schody``).}
\end{center}

\subsection{Kritéria větvení}

Otázka „který řez je nejlepší`` se redukuje na „jak měřit nečistotu uzlu``. Pro \emph{klasifikaci} se používají dvě kritéria -- \textbf{entropie} (informační zisk) a~\textbf{Gini index} -- pro \emph{regresi} \textbf{rozptylové (MSE) kritérium}. Označme $p_{\mathcal{T}}(k)$ podíl příkladů třídy $k$ v~uzlu $\mathcal{T}$.

\subsubsection*{Entropie a~informační zisk (rekapitulace z~S1)}

\begin{definice}[Entropie, booleovská entropie, informační zisk]
\textbf{Entropie} rozdělení tříd v~uzlu (v~bitech, $\log=\log_2$) je
\[ H(\mathbf{p}_{\mathcal{T}}) = -\sum_{k:\,p_{\mathcal{T}}(k)>0} p_{\mathcal{T}}(k)\,\log_2 p_{\mathcal{T}}(k). \]
Pro \emph{binární} cíl (podíl pozitivních $q$) píšeme \emph{booleovskou entropii}
\[ B(q) = -\,q\log_2 q - (1-q)\log_2(1-q),\qquad B(0)=B(1)=0,\ B(\tfrac12)=1. \]
Pro množinu s~$p$ pozitivními a~$n$ negativními příklady je entropie $B\!\big(\tfrac{p}{p+n}\big)$. Atribut $A$ rozdělí příklady do skupin; \emph{zbytková} (vážená) entropie je $\mathrm{Remainder}(A)=\sum_k \tfrac{p_k+n_k}{p+n}B\!\big(\tfrac{p_k}{p_k+n_k}\big)$ a~\textbf{informační zisk}
\[ \mathrm{Gain}(A) = B\!\Big(\tfrac{p}{p+n}\Big) - \mathrm{Remainder}(A). \]
V~uzlu vybíráme atribut (řez) s~\emph{největším} ziskem.
\end{definice}

\begin{center}
\begin{tikzpicture}
\begin{axis}[width=10cm,height=6.4cm,xlabel={$q$ (podíl pozitivních)},ylabel={$B(q)$ \small[bity]},
  xmin=0,xmax=1,ymin=0,ymax=1.14,xtick={0,0.25,0.5,0.75,1},ytick={0,0.5,1},
  axis lines=left,every axis plot/.append style={very thick},clip=false]
  \addplot[blue!60!black] table[x=q,y=B]{fig/s6_entropy.dat};
  \addplot[only marks,mark=*,mark size=1.6pt,red!70!black] coordinates {(0.5,1.0)};
  \node[font=\small,anchor=south,inner sep=2pt] at (axis cs:0.5,1.0) {max $=1$};
  \addplot[only marks,mark=*,mark size=1.6pt,blue!60!black] coordinates {(0,0) (1,0)};
  \node[font=\scriptsize,anchor=north east,inner sep=2pt] at (axis cs:0,0) {$B(0){=}0$};
  \node[font=\scriptsize,anchor=north west,inner sep=2pt] at (axis cs:1,0) {$B(1){=}0$};
\end{axis}
\end{tikzpicture}
\obrpopis{Booleovská entropie $B(q)$: je \emph{konkávní}, nulová v~čistém uzlu ($q=0$ nebo $q=1$) a~maximální ($1$ bit) při vyrovnaném poměru $q=\tfrac12$. Čím nečistější uzel, tím vyšší entropie; dělení vybíráme tak, aby vážená entropie potomků byla co nejnižší. Hodnota v~$q=\tfrac13$ (i~$\tfrac23$) je $B(\tfrac13)\approx 0{,}918$.}
\end{center}

\subsubsection*{Gini index}

\begin{definice}[Gini index (Gini impurity)]
\textbf{Gini index} uzlu $\mathcal{T}$ měří pravděpodobnost, že náhodně vybraný prvek bude \emph{špatně} označen, kdybychom mu náhodně přiřadili třídu podle rozdělení $\mathbf{p}_{\mathcal{T}}$:
\[ \mathrm{Gini}(\mathcal{T}) = \sum_k p_{\mathcal{T}}(k)\,\big(1-p_{\mathcal{T}}(k)\big) = 1 - \sum_k p_{\mathcal{T}}(k)^2. \]
Pro binární cíl s~podílem $q$ pozitivních je $\mathrm{Gini}=2q(1-q)$ (maximum $\tfrac12$ v~$q=\tfrac12$, $0$ v~$q\in\{0,1\}$). Stejně jako u~entropie hledáme řez minimalizující \emph{vážený} Gini potomků $\sum \tfrac{|I_{\mathcal{T}_\cdot}|}{|I_{\mathcal{T}}|}\mathrm{Gini}(\mathcal{T}_\cdot)$.
\end{definice}

\begin{intuice}
Gini i~entropie se chovají velmi podobně: obě jsou $0$ v~čistém uzlu, maximální u~vyrovnaného a~jsou \emph{konkávní} (proto dělení nikdy nezhorší kritérium). Liší se jen mírně tvarem a~Gini nepotřebuje logaritmus, je tedy \emph{levnější na výpočet} -- proto je výchozí v~\texttt{scikit-learn}. V~praxi vedou skoro vždy na \emph{stejný strom} (jak uvidíme u~klasifikačních SZZ otázek níže). Vztah binárních verzí je $\mathrm{Gini}=2q(1-q)$ vs.\ $B(q)$; obě křivky mají stejný tvar „kopce`` mezi $0$ a~$1$.
\end{intuice}

\subsubsection*{Regresní stromy: rozptylové (MSE) kritérium}

Je-li cíl reálné číslo, predikuje list \emph{průměr} cílů svých příkladů $\hat t_{\mathcal{T}}=\tfrac{1}{|I_{\mathcal{T}}|}\sum_{i\in I_{\mathcal{T}}} t_i$. Nečistotu měříme \emph{součtem kvadrátů odchylek} od tohoto průměru.

\begin{definice}[Kritérium regresního stromu]
\[ c_{\mathrm{SE}}(\mathcal{T}) = \sum_{i\in I_{\mathcal{T}}} (t_i - \hat t_{\mathcal{T}})^2,\qquad \hat t_{\mathcal{T}} = \frac{1}{|I_{\mathcal{T}}|}\sum_{i\in I_{\mathcal{T}}} t_i. \]
Je to $|I_{\mathcal{T}}|$-násobek \emph{rozptylu} cílů v~uzlu. Řez vybíráme tak, aby $c_{\mathrm{SE}}(\mathcal{T}_L)+c_{\mathrm{SE}}(\mathcal{T}_R)$ byl co nejmenší (největší pokles rozptylu).
\end{definice}

\begin{intuice}[Kritérium plyne ze ztrátové funkce]
Není náhoda, že regrese používá kvadratickou chybu a~klasifikace entropii/Gini. Vezmeme-li \emph{ztrátovou funkci} modelu a~do ní dosadíme optimální parametr listu (ten, který ztrátu minimalizuje), dostaneme „minimální dosažitelnou ztrátu`` daného listu -- a~právě ta slouží jako kritérium. Pro list:
\begin{itemize}
  \item dosazení optima do \emph{kvadratické ztráty} dá $c_{\mathrm{SE}}$ (regrese);
  \item dosazení optima do \emph{záporné log-věrohodnosti} (NLL, sekce~\ref{sec:evaluace}) dá $|I_{\mathcal{T}}|\cdot H(\mathbf{p}_{\mathcal{T}})$, tj.\ \emph{entropické} kritérium;
  \item dosazení optima do \emph{kvadratické ztráty} na binárním cíli $t\in\{0,1\}$ dá $|I_{\mathcal{T}}|\cdot p_{\mathcal{T}}(1-p_{\mathcal{T}})$, tj.\ (až na faktor $2$) \emph{Gini}; s~dvousložkovým one-hot cílem vyjde přesně $|I_{\mathcal{T}}|\cdot\mathrm{Gini}$.
\end{itemize}
Dělením listu chceme tuto minimální dosažitelnou ztrátu snížit, proto je \emph{minimální ztráta uzlu rovnou kritériem větvení}.
\end{intuice}

\subsection{Prořezávání (pruning)}

Strom rostlý do čistých listů snadno \textbf{overfituje} (přeučí se): dokonale klasifikuje trénovací data (každý šum dostane vlastní list), ale špatně generalizuje. \textbf{Prořezávání} (pruning) zjednodušuje strom, aby lépe generalizoval -- buď ho rovnou nenechá přerůst, nebo přerostlý strom zpětně zařízne.

\begin{definice}[Pre-pruning vs.\ post-pruning]
\begin{itemize}
  \item \textbf{Pre-pruning} (early stopping): zastavíme dělení \emph{už při stavbě}, jakmile přestane být užitečné -- omezením hloubky, minimálním počtem příkladů v~uzlu, minimálním poklesem kritéria, nebo statistickým testem ($\chi^2$, sekce~\ref{sec:testy}), zda je zisk významný. Levné, ale \emph{krátkozraké}: může zastavit příliš brzy (řez, který sám nepomůže, by mohl umožnit dobré řezy o~úroveň níž).
  \item \textbf{Post-pruning}: necháme strom plně přerůst a~pak ho \emph{zpětně} zjednodušujeme -- nahrazujeme podstromy listy. Dražší, ale spolehlivější (vidíme celý strom). Dvě hlavní varianty níže.
\end{itemize}
\end{definice}

\begin{definice}[Redukce chyby pomocí validačních dat (reduced-error pruning)]
Máme samostatnou \textbf{validační množinu}. Procházíme uzly (typicky zdola nahoru) a~uzel nahradíme \emph{listem} (s~většinovou třídou jeho trénovacích příkladů) vždy, když to \emph{nezhorší} (nebo zlepší) chybu na validačních datech. Opakujeme, dokud lze. Validační data strom při stavbě neviděl; větve, jejichž odstranění validační chybu nezhorší, tedy pravděpodobně zachycovaly spíš šum trénovacích dat než obecnou zákonitost.
\end{definice}

\begin{definice}[Cost-complexity pruning \textnormal{(weakest-link pruning, CART)}]
Zavedeme \emph{penalizaci za velikost} stromu. Pro podstrom $T$ definujme cenu
\[ R_\alpha(T) = R(T) + \alpha\,|{\sim}T|, \]
kde $R(T)$ je trénovací chyba (např.\ počet či podíl chybně klasifikovaných příkladů), $|{\sim}T|$ počet listů a~$\alpha\ge 0$ \emph{cena za list}. Pro uzel $t$ a~jeho podstrom $T_t$ je \emph{efektivní} $\alpha$
\[ \alpha_{\mathrm{eff}}(t) = \frac{R(t) - R(T_t)}{|{\sim}T_t| - 1}, \]
kde $R(t)$ je chyba, kdyby se $t$ stal listem. Iterativně ořezáváme uzel s~\emph{nejmenším} $\alpha_{\mathrm{eff}}$ (nejlevnější zjednodušení na list), čímž vzniká \emph{posloupnost} stromů pro rostoucí $\alpha$; nejlepší $\alpha$ (kompromis accuracy vs.\ velikost) se vybere křížovou validací.
\end{definice}

\begin{intuice}
Cost-complexity pruning je regularizace stromu: $\alpha$ je „pokuta`` za každý list. Při $\alpha=0$ zůstává plný strom; jak $\alpha$ roste, vyplácí se obětovat listy, které zlepšují trénovací chybu jen málo. $\alpha_{\mathrm{eff}}$ je „cena za list ušetřený sloučením`` -- ořezáváme tam, kde je nejlevnější. Reduced-error pruning je intuitivnější (přímo se ptá validačních dat), cost-complexity dává hladkou škálu složitosti laditelnou křížovou validací. Obě řeší totéž: \emph{vyměnit kus trénovací accuracy za jednodušší, lépe generalizující strom} (Ockhamova břitva).
\end{intuice}

\begin{poznamka}[Stromy jako stavební kámen ansámblů]
Jeden strom má vysoký \emph{rozptyl} (drobná změna dat může dát úplně jiný strom). Proto se kombinují do ansámblů: \emph{random forest} (bagging stromů na bootstrapových vzorcích s~náhodnou podmnožinou příznaků v~každém uzlu) a~\emph{gradient boosting} (postupné stromy opravující chybu). To je téma sekce~\ref{sec:kombinace}; zde stromy bereme jako samostatný model.
\end{poznamka}

\subsection{Řešené příklady}

Následují všechny čtyři reálné SZZ otázky k~rozhodovacím stromům (tři klasifikační, jedna regresní), krokované. Postavené stromy jsou překresleny do TikZ; číselné hodnoty (entropie, zisky, struktura) jsme ověřili nezávisle v~Pythonu (\texttt{scikit-learn} \texttt{DecisionTreeClassifier} i~ruční výpočet).

\begin{priklad}[SZZ léto 2023 č.~7: Výběr kořene a~mez $100\,\%$ accuracy]
\szzotazka{léto 2023}{7}{Rozhodovací stromy}\par
\textbf{Zadání.} Trénujeme strom na datech níže se dvěma kategorickými atributy $x_1,x_2$ a~cílem $y$. (1)~Popište algoritmus a~kritéria větvení. (2)~Podle kterého atributu se větví \emph{kořen}? Proč? (3)~Lze data klasifikovat se $100\,\%$ accuracy (přesností)? Vysvětlete.

\smallskip
\begin{center}\small
\begin{tabular}{ccc}
\toprule
$x_1$ & $x_2$ & $y$ \\
\midrule
0 & 0 & 0 \\ 1 & 0 & 1 \\ 0 & 1 & 1 \\ 2 & 1 & 0 \\
0 & 0 & 1 \\ 1 & 1 & 1 \\ 2 & 0 & 0 \\ 0 & 1 & 0 \\
\bottomrule
\end{tabular}
\end{center}

\smallskip
\textbf{Řešení.} (1)~Hladově shora dolů (rozděl a~panuj): v~každém uzlu zvol atribut s~\emph{největším informačním ziskem} (resp.\ nejmenším váženým Gini), rozděl příklady podle jeho hodnot a~rekurzivně pokračuj, dokud nejsou listy čisté nebo nedojdou atributy (pak většinová třída).

\smallskip
(2)~Z~$8$ příkladů je $4\times\,y{=}0$ a~$4\times\,y{=}1$, takže entropie kořene je $B(\tfrac48)=B(\tfrac12)=\boxed{1}$ bit. Spočteme zisk obou atributů.

\emph{Atribut $x_1$} (hodnoty $0,1,2$) rozdělí příklady:
\[ x_1{=}0:\ (2,2)\ \to B(\tfrac12)=1,\quad x_1{=}1:\ (0,2)\to 0,\quad x_1{=}2:\ (2,0)\to 0. \]
(Zde $(a,b)$ značí $a$ příkladů s~$y{=}0$ a~$b$ s~$y{=}1$.) Zbytek a~zisk:
\[ \mathrm{Remainder}(x_1)=\tfrac48\cdot 1 + \tfrac28\cdot 0 + \tfrac28\cdot 0 = \tfrac12,\qquad \mathrm{Gain}(x_1)=1-\tfrac12=\boxed{0{,}5}. \]
\emph{Atribut $x_2$} (hodnoty $0,1$): obě skupiny $x_2{=}0$ i~$x_2{=}1$ mají poměr $(2,2)$, tedy $B(\tfrac12)=1$:
\[ \mathrm{Remainder}(x_2)=\tfrac48\cdot 1 + \tfrac48\cdot 1 = 1,\qquad \mathrm{Gain}(x_2)=1-1=\boxed{0}. \]
Kořen se větví podle $\mathbf{x_1}$ (vyšší zisk; $x_2$ samo o~sobě nerozliší vůbec nic). Větve $x_1{=}1$ a~$x_1{=}2$ jsou rovnou čisté listy.

\smallskip
(3)~\textbf{Ne.} V~datech jsou \emph{shodné vstupy s~různým cílem}: řádky $(x_1,x_2)=(0,0)$ se objeví s~$y{=}0$ i~$y{=}1$, stejně $(0,1)$ s~$y{=}1$ i~$y{=}0$. Žádný model rozhodující jen podle $(x_1,x_2)$ tyto dva páry nerozliší -- ať strom roste jakkoli hluboko, ve větvi $x_1{=}0$ zbudou čtyři příklady $(0,0,0),(0,0,1),(0,1,1),(0,1,0)$, které mají jen dvě různé kombinace vstupů, ale obě třídy. Maximální dosažitelná accuracy je tedy $\boxed{<100\,\%}$ (konkrétně $6/8$: ve větvi $x_1{=}0$ predikuje list většinu, dvě protichůdná pozorování zůstanou špatně).

\begin{center}
\begin{tikzpicture}[scale=1.15]
  \node[treenode] (r) at (0,3) {$x_1?$};
  \node[treenode] (z) at (-2.6,1.4) {$x_2?$};
  \node[treeleaf] (j) at (0,1.4) {$y{=}1$};
  \node[treeleaf] (d) at (2.4,1.4) {$y{=}0$};
  \draw[vetev] (r) -- node[popis,above left]{$0$} (z);
  \draw[vetev] (r) -- node[popis,left]{$1$} (j);
  \draw[vetev] (r) -- node[popis,above right]{$2$} (d);
  \node[treeleaf,fill=black!6] (z0) at (-3.6,-0.1) {$(2,2)$};
  \node[treeleaf,fill=black!6] (z1) at (-1.6,-0.1) {$(2,2)$};
  \draw[vetev] (z) -- node[popis,above left]{$0$} (z0);
  \draw[vetev] (z) -- node[popis,above right]{$1$} (z1);
\end{tikzpicture}
\obrpopis{Strom pro léto 2023: kořen $x_1$. Větve $x_1{=}1$ a~$x_1{=}2$ jsou čisté listy. Větev $x_1{=}0$ se sice dá ještě dělit podle $x_2$, ale obě podskupiny zůstanou poměr $(2,2)$ -- nerozdělitelný šum se shodnými vstupy a~různými cíli. Proto $100\,\%$ accuracy není dosažitelná.}
\end{center}
\end{priklad}

\begin{priklad}[SZZ podzim 2024 č.~31: Stavba stromu, kořen $A$, prořezávání validačními daty]
\szzotazka{podzim 2024}{31}{Rozhodovací strom}\par
\textbf{Zadání.} Data se třemi atributy $A,B,C$ a~cílem $Y\in\{\text{n},\text{y}\}$. (1)~Postavte plný strom (bez omezení). (2)~Pro atribut v~kořeni popište výpočet kritéria. (3)~Proveďte prořezávání validačními daty.

\smallskip
\begin{center}\small
\begin{minipage}{0.4\linewidth}\centering Trénovací:\par\medskip
\begin{tabular}{ccccc}\toprule & $A$ & $B$ & $C$ & $Y$ \\\midrule
0 & 1 & 2 & 6 & n \\ 1 & 1 & 2 & 6 & n \\ 2 & 3 & 2 & 8 & y \\ 3 & 3 & 4 & 6 & y \\
4 & 3 & 4 & 8 & y \\ 5 & 3 & 4 & 8 & n \\ 6 & 3 & 4 & 8 & n \\ 7 & 3 & 4 & 6 & y \\\bottomrule
\end{tabular}\end{minipage}\hfill
\begin{minipage}{0.4\linewidth}\centering Validační:\par\medskip
\begin{tabular}{ccccc}\toprule & $A$ & $B$ & $C$ & $Y$ \\\midrule
8  & 1 & 2 & 8 & n \\ 9  & 3 & 4 & 6 & y \\ 10 & 3 & 4 & 6 & y \\ 11 & 3 & 4 & 8 & y \\\bottomrule
\end{tabular}\end{minipage}
\end{center}

\smallskip
\textbf{Řešení.} (1)+(2)~\textbf{Kořen.} Z~$8$ příkladů jsou $4\times$ n a~$4\times$ y, entropie kořene $B(\tfrac12)=1$. Spočteme zisk řezů (kontrola jednotlivých prahů):
\[
\begin{aligned}
A\le 2\ (\text{tj.\ }A{=}1\text{ vs.\ }A{=}3): &\quad (2,0)\,|\,(2,4)\Rightarrow \mathrm{Remainder}=\tfrac28\cdot 0 + \tfrac68 B(\tfrac26)=\tfrac34\cdot 0{,}918\approx 0{,}689,\\
 &\quad \mathrm{Gain}\approx 1-0{,}689 = \boxed{0{,}311};\\
B\le 2: &\quad (2,1)\,|\,(2,3)\Rightarrow \mathrm{Gain}\approx 0{,}049;\\
C\le 6: &\quad (2,2)\,|\,(2,2)\Rightarrow \mathrm{Gain}=0.
\end{aligned}
\]
Největší zisk má $\boxed{A}$ (oficiální náčrt uvádí $\approx 0{,}32$ kvůli zaokrouhlení tabulky logaritmů). Větev $A\le 2$ je čistý list (n). V~pravé větvi $A>2$ ($6$ příkladů, $2$ n, $4$ y, $B(\tfrac26)\approx 0{,}918$) dělíme dál podle $C\le 7$: vlevo $(0,2)$ čistý list y, vpravo $C>7$ má $(2,2)$, kde rozdělí $B\le 3$ na čistý list y $(0,1)$ a~list $(2,1)$. Strom vyjde \emph{stejný pro entropii i~Gini} (ověřeno).

\begin{center}
\begin{tikzpicture}[scale=1.15]
  \node[treenode] (a) at (0,4.2) {$A \le 2$};
  \node[treeleaf] (ln) at (-2.8,2.7) {n \scriptsize$(2,0)$};
  \node[treenode] (c) at (1.4,2.7) {$C \le 7$};
  \draw[vetev] (a) -- node[popis,above left]{ano} (ln);
  \draw[vetev] (a) -- node[popis,above right]{ne} (c);
  \node[treeleaf] (ly) at (-0.4,1.2) {y \scriptsize$(0,2)$};
  \node[treenode] (b) at (2.8,1.2) {$B \le 3$};
  \draw[vetev] (c) -- node[popis,above left]{ano} (ly);
  \draw[vetev] (c) -- node[popis,above right]{ne} (b);
  \node[treeleaf] (ly2) at (1.6,-0.3) {y \scriptsize$(0,1)$};
  \node[treeleaf,fill=black!6] (ln2) at (4.0,-0.3) {n \scriptsize$(2,1)$};
  \draw[vetev] (b) -- node[popis,above left]{ano} (ly2);
  \draw[vetev] (b) -- node[popis,above right]{ne} (ln2);
\end{tikzpicture}
\obrpopis{Plný strom (podzim 2024), totožný pro entropii i~Gini. Dvojice $(n,y)$ udává počty tříd v~uzlu. Kořen $A$, pak v~pravé větvi $C$, pak $B$. List $(2,1)$ není čistý (zbylé příklady mají shodné $A,B,C$, ale různé $Y$), predikuje většinu n.}
\end{center}

\smallskip
(3)~\textbf{Prořezávání validačními daty (reduced-error pruning).} Spustíme $4$ validační příklady plným stromem:
\[
\begin{array}{c|cccc|cc}
\text{id} & A & B & C & Y & \text{plný strom} & \text{po ořezu } A{>}2\!\to\!\text{y} \\\hline
8  & 1 & 2 & 8 & \text{n} & \text{n \faCheck} & \text{n \faCheck} \\
9  & 3 & 4 & 6 & \text{y} & \text{y \faCheck} & \text{y \faCheck} \\
10 & 3 & 4 & 6 & \text{y} & \text{y \faCheck} & \text{y \faCheck} \\
11 & 3 & 4 & 8 & \text{y} & \text{n \faTimes} & \text{y \faCheck}
\end{array}
\]
Plný strom dělá na validaci $1$ chybu (id 11: prochází $A{>}2,C{>}7,B{>}3$ do listu n, ale skutečnost je y). Nahradíme-li celý \emph{pravý podstrom} (uzel $C$, tj.\ vše pod $A>2$) \emph{listem} s~většinovou trénovací třídou y ($4$ y vs.\ $2$ n), validační chyba klesne na $\boxed{0}$. Prořezání tedy chybu \emph{zlepšilo} -- ořez přijmeme. (Postupem zdola nahoru dojdeme ke stejnému stromu: nejdřív uzel $B\le 3$, jehož trénovací příklady mají poměr $(2,2)$ -- ať remízu rozhodneme jakkoli, validační chyba se nezhorší, takže se ořízne; pak uzel $C\le 7$ nahrazený listem y dá validační chybu $0$, takže se ořízne také.) Strom se zmenší na $A\le 2 \to$ n, jinak y. (Alternativně lze použít cost-complexity pruning.)
\end{priklad}

\begin{priklad}[SZZ jaro 2026 č.~30: Zadaný kořen $D$, kritérium pravého podstromu, prořezávání]
\szzotazka{jaro 2026}{30}{Rozhodovací strom}\par
\textbf{Zadání.} (a)~Postavte strom, do kořene zvolte atribut $D$; volby zdůvodněte, pro \emph{pravý podstrom} $D$ spočtěte kritérium pro všechny atributy. (b)~Vysvětlete prořezávání, uveďte neprořezaný a~prořezaný strom a~kritérium prořezávání. Sloupec \emph{id} se k~učení nepoužívá.

\smallskip
\begin{center}\small
\begin{tabular}{cccccc}\toprule
id & $A$ & $B$ & $C$ & $D$ & $Y$ \\\midrule
0 & 3 & 4 & 8 & 6 & 0 \\ 1 & 1 & 2 & 8 & 6 & 0 \\ 2 & 1 & 2 & 8 & 3 & 1 \\ 3 & 3 & 2 & 8 & 6 & 1 \\
4 & 1 & 4 & 6 & 6 & 1 \\ 5 & 3 & 2 & 8 & 6 & 0 \\ 6 & 3 & 4 & 6 & 3 & 0 \\ 7 & 3 & 4 & 8 & 3 & 1 \\\bottomrule
\end{tabular}
\end{center}

\smallskip
\textbf{Řešení.} (a)~Kořen je zadán: $D\le 4{,}5$ (tj.\ $D{=}3$ vlevo, $D{=}6$ vpravo). Kořen má $4\times 0$, $4\times 1$, entropie $1$.

\emph{Levý podstrom} $D{=}3$ (id $2,6,7$, cíle $1,0,1$, tj.\ poměr $(1,2)$): dělení podle $C\le 7$ dá vlevo $C{=}6$ jediný příklad id 6 (list $0$), vpravo $C{=}8$ příklady id $2,7$ (oba $Y{=}1$, čistý list $1$). Levý podstrom je tedy hotový.

\emph{Pravý podstrom} $D{=}6$ (id $0,1,3,4,5$, cíle $0,0,1,1,0$, poměr $(3,2)$, $B(\tfrac25)\approx 0{,}971$) -- spočítáme \emph{vážené entropie po dělení} pro všechny atributy:
\[
\begin{aligned}
A\le 2:\ & (1,1)_{n2}\,|\,(2,1)_{n3} \Rightarrow \tfrac25\cdot 1 + \tfrac35 B(\tfrac13) \approx 0{,}4 + 0{,}6\cdot 0{,}918 \approx \boxed{0{,}951};\\
B\le 3:\ & (2,1)_{n3}\,|\,(1,1)_{n2} \Rightarrow \tfrac35 B(\tfrac13)+\tfrac25\cdot 1 \approx \boxed{0{,}951};\\
C\le 7:\ & (0,1)_{n1}\,|\,(3,1)_{n4} \Rightarrow \tfrac15\cdot 0 + \tfrac45 B(\tfrac14) \approx \tfrac45\cdot 0{,}811 \approx \boxed{0{,}649}.
\end{aligned}
\]
Nejnižší vážená entropie (a~tedy \emph{největší zisk}, $\approx 0{,}971-0{,}649=0{,}322$) má $\mathbf{C}$. Vlevo $C{=}6$ čistý list $1$ (id 4), vpravo $C{=}8$ má id $0,1,3,5$, poměr $(3,1)$; dělí se podle $A\le 2$ (id 1 list $0$; stejný zisk by dal i~$B\le 3$) a~$A>2$ (id $0,3,5$, poměr $(2,1)$), kde řez $B\le 3$ dá id $3,5$ (poměr $(1,1)$, nerozdělitelný šum) a~$B>3$ čistý list $0$ (id $0$). Strom dává stejné predikce pro entropii i~Gini (ověřeno); nakreslené pořadí uzlů odpovídá entropii.

\begin{center}
\begin{tikzpicture}[scale=1.0]
  \node[treenode] (d) at (0,5.6) {$D \le 4{,}5$};
  % leva vetev
  \node[treenode] (cl) at (-3.4,4.0) {$C \le 7$};
  \draw[vetev] (d) -- node[popis,above left]{ano} (cl);
  \node[treeleaf] (cl6) at (-4.8,2.4) {$0$ \scriptsize$(1,0)$};
  \node[treeleaf] (cl8) at (-2.6,2.4) {$1$ \scriptsize$(0,2)$};
  \draw[vetev] (cl) -- node[popis,above left]{ano} (cl6);
  \draw[vetev] (cl) -- node[popis,above right]{ne} (cl8);
  % prava vetev
  \node[treenode] (cr) at (2.6,4.0) {$C \le 7$};
  \draw[vetev] (d) -- node[popis,above right]{ne} (cr);
  \node[treeleaf] (cr6) at (1.0,2.4) {$1$ \scriptsize$(0,1)$};
  \node[treenode] (ar) at (3.8,2.4) {$A \le 2$};
  \draw[vetev] (cr) -- node[popis,above left]{ano} (cr6);
  \draw[vetev] (cr) -- node[popis,above right]{ne} (ar);
  \node[treeleaf] (ar2) at (2.4,0.8) {$0$ \scriptsize$(1,0)$};
  \node[treenode] (br) at (5.2,0.8) {$B \le 3$};
  \draw[vetev] (ar) -- node[popis,above left]{ano} (ar2);
  \draw[vetev] (ar) -- node[popis,above right]{ne} (br);
  \node[treeleaf,fill=black!6] (br3) at (4.2,-0.8) {\scriptsize$(1,1)$};
  \node[treeleaf] (br4) at (6.2,-0.8) {$0$ \scriptsize$(1,0)$};
  \draw[vetev] (br) -- node[popis,above left]{ano} (br3);
  \draw[vetev] (br) -- node[popis,above right]{ne} (br4);
\end{tikzpicture}
\obrpopis{Neprořezaný strom do maximální hloubky (jaro 2026), shodný pro entropii i~Gini. Dvojice $(n_0,n_1)$ udává počty tříd. Pravá větev $D{=}6$ se dělí podle $C$, pak $A$, pak $B$; list $(1,1)$ je nerozdělitelný (shodné vstupy id 3 a~5 s~různým $Y$).}
\end{center}

(b)~\textbf{Prořezávání} odstraňuje overfitting: na trénovacích datech je chyba malá, ale na nových větší než u~jednoduššího stromu. Nahradíme uzel $\mathbf{A\le 2}$ (celý pravý spodní podstrom) \emph{listem}: jeho trénovací příklady id $0,1,3,5$ mají poměr $(3,1)$, většina je $0$. Plný podstrom stejně jeden příklad (id 3 vs.\ 5) klasifikuje špatně (list $(1,1)$), takže náhrada listem $0$ chybu \emph{nezhorší} a~strom zjednoduší (méně listů). Prořezaný strom: $D\le 4{,}5 \to (C\le 7?\,0:1)$; jinak $C\le 7 \to 1$, $C>7 \to 0$.
\[ \text{Kritérium ořezu (cost-complexity):}\quad \alpha_{\mathrm{eff}}(t) = \frac{R(t) - R(T_t)}{|{\sim}T_t| - 1}, \]
ořezává se uzel s~nejmenším $\alpha_{\mathrm{eff}}$, dokud je pod prahem; $R(t)$ je chyba uzlu jako listu, $R(T_t)$ chyba podstromu, $|{\sim}T_t|$ počet jeho listů. (Alternativně reduced-error pruning podle validačních dat.)
\end{priklad}

\begin{priklad}[SZZ podzim 2022 č.~19: Regresní strom -- výběr atributu]
\szzotazka{podzim 2022}{19}{Rozhodovací stromy}\par
\textbf{Zadání.} (1)~Definujte rozhodovací strom a~popište algoritmus jeho trénování. (2)~Uveďte kritérium pro výběr atributu pro \emph{klasifikaci} a~pro \emph{regresi}. (3)~Pro data níže (kategorické atributy $A_1,A_2$, číselný cíl $C$) určete, podle kterého atributu se dělí kořen pro \emph{regresi}, a~proč.

\smallskip
\begin{center}\small
\begin{tabular}{ccc}
\toprule
$A_1$ & $A_2$ & $C$ \\
\midrule
0 & 0 & 5 \\ 0 & 0 & 3 \\ 0 & 1 & 4 \\
1 & 0 & 6 \\ 1 & 1 & 2 \\ 1 & 1 & 4 \\
\bottomrule
\end{tabular}
\end{center}

\smallskip
\textbf{Řešení.} \emph{(1)} Rozhodovací strom je stromový model: vnitřní uzly testují atribut, větve odpovídají jeho hodnotám a~listy nesou predikci (třídu nebo číslo). Trénuje se hladově shora dolů (rozděl a~panuj): v~každém uzlu se vybere atribut podle zvoleného kritéria, data se podle něj rozdělí a~rekurzivně se pokračuje, dokud nejsou listy čisté, nedojdou atributy nebo není splněna zastavovací podmínka.

\emph{(2)} \emph{Klasifikace}: maximalizace informačního zisku (pokles entropie) nebo minimalizace Gini indexu. \emph{Regrese}: minimalizace vážené sumy čtvercových odchylek (SSE) v~potomcích, ekvivalentně maximalizace redukce rozptylu.

\emph{(3)} Celkový průměr je $\bar C = 4$, SSE rodiče $=\sum (C_i-4)^2 = 1+1+0+4+4+0 = \boxed{10}$.
\begin{itemize}
  \item \emph{Dělení podle $A_1$}: $A_1{=}0 \to \{5,3,4\}$, průměr $4$, SSE $2$; $A_1{=}1 \to \{6,2,4\}$, průměr $4$, SSE $8$. Celkem $2+8=10$, redukce $\boxed{0}$.
  \item \emph{Dělení podle $A_2$}: $A_2{=}0 \to \{5,3,6\}$, průměr $\approx 4{,}67$, SSE $\approx 4{,}67$; $A_2{=}1 \to \{4,2,4\}$, průměr $\approx 3{,}33$, SSE $\approx 2{,}67$. Celkem $\approx 7{,}33$, redukce $\boxed{\approx 2{,}67}$.
\end{itemize}
V~kořeni proto dělíme podle $\mathbf{A_2}$ (větší redukce rozptylu). Atribut $A_1$ nedává žádnou redukci, protože obě jeho větve mají stejný průměr $4$ jako rodič.
\end{priklad}

\begin{poznamka}[Co si odnést -- rozhodovací stromy]
\begin{itemize}
  \item \emph{Inference} = sleduj testy do listu. \emph{Učení} = hladově shora dolů (rozděl a~panuj): v~každém uzlu řez (příznak + práh) minimalizující kritérium potomků.
  \item \emph{Kritéria}: klasifikace -- \textbf{entropie}/informační zisk $\mathrm{Gain}=B(\tfrac{p}{p+n})-\mathrm{Remainder}(A)$ nebo \textbf{Gini} $=1-\sum_k p_k^2$ (levnější, prakticky stejný strom); regrese -- \textbf{rozptyl/MSE} $c_{\mathrm{SE}}=\sum(t_i-\hat t)^2$. Kritérium plyne z~minimální dosažitelné ztráty listu.
  \item \emph{Prořezávání}: \textbf{pre-pruning} (zastav brzy: hloubka, min.\ příkladů, $\chi^2$) vs.\ \textbf{post-pruning} (přerost a~zařízni): \textbf{reduced-error} (ořež, co nezhorší validační chybu) a~\textbf{cost-complexity} (penalizace $\alpha$ za list, ořez podle nejmenšího $\alpha_{\mathrm{eff}}$, $\alpha$ ladí křížová validace).
  \item \emph{Mez stromu}: shodné vstupy s~různým cílem nelze rozdělit (léto 2023: $100\,\%$ accuracy nedosažitelná). Jeden strom má vysoký rozptyl $\Rightarrow$ ansámbly (sekce~\ref{sec:kombinace}).
\end{itemize}
\end{poznamka}

\section{Metoda podpůrných vektorů (SVM)}\label{sec:svm}

\textbf{Metoda podpůrných vektorů} (support vector machines, SVM) je lineární klasifikátor, který mezi všemi oddělujícími nadrovinami vybere tu s~\emph{největším marginem} (margin, odstup od dat). Ve spojení s~\emph{kernel trikem} zvládne i~nelineárně oddělitelná data a~patří k~nejúspěšnějším klasifikátorům „z~krabice``. Tato sekce jde do hloubky, kterou S1 (sekce \emph{Strojové učení}) výslovně přenechává S3.

\begin{poznamka}[Vztah k~S1]
Okruh S1 zavádí SVM \emph{přehledově}: maximální margin, podpůrné vektory a~jádrovou transformaci jako hlavní myšlenku, s~dvoupanelovým obrázkem. Zde to nezdvojujeme. Stavíme na tom a~jdeme do detailu: jak přesně vypadá optimalizační úloha (QP), proč závisí jen na skalárních součinech, soft margin a~hinge loss, konkrétní jádra a~klasifikace do více tříd.
\end{poznamka}

\subsection{Lineárně separabilní třídy: maximální margin}

Mějme binární klasifikaci s~daty $(\mathbf{x}_i,t_i)$, kde cíle značíme \textbf{symetricky} $t_i\in\{-1,+1\}$ (ne $\{0,1\}$ jako u~logistické regrese -- u~SVM se to vyplatí). Lineární model je
\[ y(\mathbf{x}) = \mathbf{w}^{\!\top}\mathbf{x} + b, \qquad \text{predikce} = \sign\,y(\mathbf{x}). \]
Rovnice $\mathbf{w}^{\!\top}\mathbf{x}+b=0$ určuje \emph{oddělující nadrovinu} (v~rovině přímku); vektor $\mathbf{w}$ je na ni kolmý. Předpokládejme, že data jsou \textbf{lineárně separabilní}, tj.\ existuje nadrovina, která bezchybně oddělí obě třídy. Takových nadrovin je obecně nekonečně mnoho; SVM mezi nimi volí jednu konkrétní.

\begin{definice}[Margin, podpůrné vektory]
\textbf{Margin} (margin, šíře odstupu) oddělující nadroviny je vzdálenost od nadroviny k~nejbližšímu trénovacímu bodu (na obě strany). \textbf{Maximálně marginový (max-margin) klasifikátor} je nadrovina s~\emph{největším možným marginem}. Trénovací body ležící \emph{nejblíž} k~nadrovině (přesně na marginu) se nazývají \textbf{podpůrné vektory} (support vectors); právě ony nadrovinu určují -- ostatní body lze odebrat, aniž se hranice změní.
\end{definice}

\begin{intuice}
Mezi mnoha přímkami, které data oddělí, vybere SVM tu „nejodvážnější`` -- s~co největší rezervou na obě strany. Široký margin zlepšuje \emph{generalizaci}: malé posuny dat hranici nepřekročí, takže model je robustní vůči šumu. Tomuto principu se říká \emph{maximalizace marginu} a~lze ho chápat jako jistou formu regularizace.
\end{intuice}

\subsubsection*{Od marginu k~optimalizační úloze}

Vzdálenost bodu $\mathbf{x}_i$ od nadroviny je $\dfrac{|y(\mathbf{x}_i)|}{\|\mathbf{w}\|}$. Klasifikuje-li model všechny body správně, je $t_i\,y(\mathbf{x}_i)>0$, takže $|y(\mathbf{x}_i)| = t_i\,y(\mathbf{x}_i)$ a~vzdálenost je $\dfrac{t_i\,y(\mathbf{x}_i)}{\|\mathbf{w}\|}$. Chceme maximalizovat margin, tj.\ vzdálenost \emph{nejbližšího} bodu:
\[ \argmax_{\mathbf{w},b}\ \frac{1}{\|\mathbf{w}\|}\min_i\big[t_i(\mathbf{w}^{\!\top}\mathbf{x}_i+b)\big]. \]
Takto se úloha optimalizuje špatně. Pomůže pozorování, že model je \emph{invariantní vůči přeškálování}: vynásobíme-li $\mathbf{w}$ i~$b$ stejnou konstantou, nadrovina (ani znaménko $y$) se nezmění. Můžeme proto škálu \uv{ukotvit} tak, aby pro nejbližší body platilo $t_i\,y(\mathbf{x}_i)=1$. Pak pro \emph{všechny} body platí $t_i\,y(\mathbf{x}_i)\ge 1$ a~maximalizace $1/\|\mathbf{w}\|$ je totéž co minimalizace $\tfrac12\|\mathbf{w}\|^2$.

\begin{veta}[Primární úloha hard-margin SVM (kvadratické programování)]
Pro lineárně separabilní data je max-margin nadrovina řešením úlohy
\[ \boxed{\ \min_{\mathbf{w},b}\ \tfrac12\|\mathbf{w}\|^2 \quad\text{za podmínek}\quad t_i(\mathbf{w}^{\!\top}\mathbf{x}_i+b)\ge 1\ \text{pro všechna }i.\ } \]
Účelová funkce je \emph{kvadratická} a~podmínky \emph{lineární} (afinní) v~$\mathbf{w},b$, jde tedy o~úlohu \textbf{kvadratického programování} (QP). Účelová funkce je navíc \emph{konvexní}, takže optimum je jednoznačné (globální). Při této normalizaci je margin roven
\[ \text{margin} = \frac{1}{\|\mathbf{w}\|},\qquad \text{tj.\ celková šíře pásu mezi marginy} = \frac{2}{\|\mathbf{w}\|}. \]
\end{veta}

\begin{intuice}[Proč $2/\|\mathbf{w}\|$]
Body s~$t_i\,y(\mathbf{x}_i)=1$ leží na dvou rovnoběžných nadrovinách $\mathbf{w}^{\!\top}\mathbf{x}+b=+1$ a~$\mathbf{w}^{\!\top}\mathbf{x}+b=-1$ (kladný a~záporný margin). Vzdálenost mezi dvěma rovnoběžkami $\mathbf{w}^{\!\top}\mathbf{x}+b=\pm1$ je $\dfrac{2}{\|\mathbf{w}\|}$. Minimalizovat $\|\mathbf{w}\|$ tedy znamená \emph{rozšiřovat} pás mezi marginy.
\end{intuice}

\begin{center}
\begin{tikzpicture}[scale=1.3]
  % decision boundary: x1 + x2 = 3  (w = (1,1))
  % positive margin: x1+x2 = 4 ; negative margin: x1+x2 = 2
  \draw[osa] (0,0) -- (5,0) node[right,font=\scriptsize]{$x_1$};
  \draw[osa] (0,0) -- (0,5) node[above,font=\scriptsize]{$x_2$};
  % margin band fill (between x1+x2=2 and x1+x2=4)
  \fill[blue!7] (0.0,4.0) -- (4.0,0.0) -- (2.0,0.0) -- (0.0,2.0) -- cycle;
  % margin lines (dashed): x1+x2=4 and x1+x2=2
  \draw[margin] (0.0,4.0) -- (4.0,0.0);
  \draw[margin] (0.0,2.0) -- (2.0,0.0);
  % decision boundary (solid): x1+x2=3
  \draw[hranice] (0.0,3.0) -- (3.0,0.0);
  \node[font=\scriptsize,black!75,anchor=west] at (3.15,0.35) {$\mathbf{w}^{\!\top}\mathbf{x}+b=0$};
  % positive class (cls2 = red squares) above the band
  \node[cls2] (p1) at (3.4,2.2) {};
  \node[cls2] (p2) at (4.2,1.4) {};
  \node[cls2,podpora] (sp) at (2.6,1.4) {}; % support vector on positive margin (x1+x2=4)
  \node[font=\scriptsize,red!65!black,anchor=south west] at (3.5,2.25) {třída $+1$};
  % negative class (cls1 = blue circles) below the band
  \node[cls1] (n1) at (0.5,0.5) {};
  \node[cls1] (n2) at (1.3,0.2) {};
  \node[cls1,podpora] (sn) at (1.0,1.0) {}; % support vector on negative margin (x1+x2=2)
  \node[font=\scriptsize,blue!60!black,anchor=north east] at (1.25,0.1) {třída $-1$};
  % normal vector w (perpendicular to boundary, direction (1,1)) -- v prazdne casti vlevo nahore
  \draw[->,>=Stealth,thick,black] (0.9,2.1) -- (1.45,2.65) node[midway,above left,font=\scriptsize]{$\mathbf{w}$};
  % margin arrow (perpendicular distance between boundary and pos margin) -- vpravo dole, mimo body
  \draw[<->,>=Stealth,blue!55!black,thick] (2.5,0.5) -- (3.0,1.0);
  \node[font=\scriptsize,blue!55!black,anchor=west] at (3.05,0.9) {margin $=\tfrac{1}{\|\mathbf{w}\|}$};
  % label support vectors (umisteno tak, aby nekolidovalo s w ani se sebou)
  \node[font=\scriptsize,anchor=north] at (2.6,1.2) {p.\,v.};
  \node[font=\scriptsize,anchor=west] at (1.2,1.05) {p.\,v.};
\end{tikzpicture}
\obrpopis{Max-margin separátor: plná čára je rozhodovací hranice $\mathbf{w}^{\!\top}\mathbf{x}+b=0$, čárkované čáry jsou dva marginy $\mathbf{w}^{\!\top}\mathbf{x}+b=\pm1$. Zakroužkované body leží přesně na marginu -- to jsou \emph{podpůrné vektory} (p.\,v.), které hranici určují. Vektor $\mathbf{w}$ je na hranici kolmý; šíře pásu je $2/\|\mathbf{w}\|$.}
\end{center}

\subsubsection*{Duální tvar a~role skalárních součinů}

Úloha s~nerovnostními podmínkami se řeší přes \emph{Lagrangeovy multiplikátory} $a_i\ge 0$ (jeden na každou podmínku; podmínky optimality KKT -- Karush--Kuhn--Tucker). Plnou derivaci vynecháváme; klíčové je, jak vypadá výsledek. Po dosazení optimálních $\mathbf{w}=\sum_i a_i t_i\,\mathbf{x}_i$ vznikne ekvivalentní \textbf{duální úloha}
\[ \max_{\mathbf{a}}\ \sum_i a_i - \tfrac12\sum_i\sum_j a_i a_j\,t_i t_j\,\underbrace{\mathbf{x}_i^{\!\top}\mathbf{x}_j}_{\text{skalární součin}}\quad\text{za podmínek}\quad a_i\ge 0,\ \sum_i a_i t_i = 0. \]

\begin{intuice}[Dvě zásadní pozorování]
\begin{itemize}
  \item \textbf{Závislost jen na skalárních součinech.} V~duální úloze se data objevují \emph{výhradně} přes skalární součiny $\mathbf{x}_i^{\!\top}\mathbf{x}_j$. To je brána ke kernel triku (níže): stačí umět počítat skalární součiny, samotné vektory nepotřebujeme.
  \item \textbf{Řídkost: jen podpůrné vektory.} Z~podmínek KKT plyne $a_i\,\big(t_i y(\mathbf{x}_i)-1\big)=0$ pro každé $i$. Tedy buď je bod přesně na marginu ($t_i y(\mathbf{x}_i)=1$), nebo $a_i=0$. Body s~$a_i=0$ se v~predikci $y(\mathbf{x})=\sum_i a_i t_i\,\mathbf{x}_i^{\!\top}\mathbf{x}+b$ neuplatní. Stačí tedy uchovat jen \emph{podpůrné vektory} (body s~$a_i>0$). SVM je sice neparametrický model (počet podpůrných vektorů roste s~daty), ale paměťově úsporný.
\end{itemize}
\end{intuice}

\subsection{Lineárně neseparabilní třídy: soft margin}

Reálná data bývají \emph{zašuměná} a~oddělit je beze zbytku přímkou nelze (hard-margin úloha pak nemá řešení). \textbf{Soft margin} (měkký margin) povolí, aby některé body byly \emph{uvnitř} marginu, nebo i~na \emph{špatné straně} hranice, ale za pokutu.

\begin{definice}[Slack proměnné, soft-margin SVM]
Pro každý bod zavedeme \textbf{slack proměnnou} $\xi_i\ge 0$ měřící porušení marginu. Podmínku $t_i y(\mathbf{x}_i)\ge 1$ uvolníme na $t_i y(\mathbf{x}_i)\ge 1-\xi_i$. Optimalizujeme
\[ \boxed{\ \min_{\mathbf{w},b,\boldsymbol{\xi}}\ \tfrac12\|\mathbf{w}\|^2 + C\sum_i \xi_i \quad\text{za podmínek}\quad t_i(\mathbf{w}^{\!\top}\mathbf{x}_i+b)\ge 1-\xi_i,\ \ \xi_i\ge 0.\ } \]
Hodnota slacku rozlišuje polohu bodu: $\xi_i=0$ bod \emph{vně} marginu (správně, s~rezervou), $0<\xi_i<1$ bod \emph{uvnitř} marginu (správně, ale bez rezervy), $\xi_i=1$ bod \emph{na} hranici a~$\xi_i>1$ bod na \emph{špatné straně} hranice (chyba klasifikace).
\end{definice}

\begin{intuice}[Význam konstanty $C$]
Konstanta $C>0$ je \emph{kompromis} mezi šíří marginu a~chybami: penalizuje součet slacků $\sum_i\xi_i$.
\begin{itemize}
  \item \textbf{Velké $C$} -- chyby jsou drahé, model se snaží klasifikovat vše správně (úzký margin, blíž k~hard-margin; riziko \emph{overfittingu}, tj.\ přeučení).
  \item \textbf{Malé $C$} -- chyby jsou levné, model upřednostní široký margin a~toleruje porušení (víc \emph{regularizace}, hladší hranice).
\end{itemize}
$C$ je hyperparametr; ladí se křížovou validací (viz sekce~\ref{sec:evaluace}).
\end{intuice}

\begin{center}
\begin{tikzpicture}[scale=1.4]
  % geometrie: hranice x1+x2=3, kladny margin x1+x2=4, zaporny margin x1+x2=2
  % smer kolmy na primky = (1,1)/sqrt2; krok (0.5,0.5) posune x1+x2 o 1
  \draw[osa] (0,0) -- (5.2,0) node[right,font=\scriptsize]{$x_1$};
  \draw[osa] (0,0) -- (0,5.2) node[above,font=\scriptsize]{$x_2$};
  % vyplň pásu mezi marginy x1+x2=2 a x1+x2=4
  \fill[blue!6] (0.0,4.0) -- (4.0,0.0) -- (2.0,0.0) -- (0.0,2.0) -- cycle;
  % marginy (čárkované)
  \draw[margin] (0.4,3.6) -- (3.6,0.4);
  \draw[margin] (0.2,1.8) -- (1.8,0.2);
  \node[font=\scriptsize,black!55,anchor=west] at (3.55,0.5) {margin $+1$};
  \node[font=\scriptsize,black!55,anchor=west] at (1.75,0.3) {margin $-1$};
  % rozhodovaci hranice x1+x2=3 (plna)
  \draw[hranice] (0.3,2.7) -- (3.0,0.0);
  \node[font=\scriptsize,black!75,anchor=west] at (2.65,0.55) {hranice};
  % --- spravne klasifikovane kladne (cervene ctverce), xi=0 ---
  \node[cls2] at (3.3,1.5) {};   % x1+x2=4.8
  \node[cls2] at (4.2,1.0) {};   % x1+x2=5.2
  \node[cls2,podpora] at (2.2,1.8) {}; % podpurny vektor na kladnem marginu (x1+x2=4)
  % --- kladny bod UVNITR marginu (spravna strana, 0<xi<1): x1+x2=3.5 ---
  \node[cls2] (pin) at (2.6,0.9) {};
  % slack: kolma usecka od kladneho marginu k bodu, foot na x1+x2=4 -> (2.85,1.15)
  \draw[<->,>=Stealth,orange!80!black,thick] (2.85,1.15) -- (2.6,0.9);
  \node[font=\scriptsize,orange!70!black,anchor=west] at (2.9,1.0) {$0<\xi_i<1$};
  % --- spravne klasifikovane zaporne (modra kolecka), xi=0 ---
  \node[cls1] at (0.5,0.5) {};   % x1+x2=1.0
  \node[cls1] at (0.9,0.4) {};   % x1+x2=1.3
  \node[cls1,podpora] at (0.8,1.2) {}; % podpurny vektor na zapornem marginu (x1+x2=2)
  % --- zaporny bod na SPATNE strane (xi>1, chyba): x1+x2=3.4 ---
  \node[cls1] (nwr) at (1.2,2.2) {};
  % slack: kolma usecka od zaporneho marginu k bodu, foot na x1+x2=2 -> (0.5,1.5)
  \draw[<->,>=Stealth,red!75!black,thick] (0.5,1.5) -- (1.2,2.2);
  \node[font=\scriptsize,red!70!black,anchor=south east] at (1.15,2.25) {$\xi_j>1$};
  % legenda
  \node[font=\scriptsize,red!65!black,anchor=west] at (3.4,2.2) {třída $+1$ (čtverce)};
  \node[font=\scriptsize,blue!60!black,anchor=west] at (1.9,2.9) {třída $-1$ (kolečka)};
\end{tikzpicture}
\obrpopis{Soft margin: většina bodů leží správně vně marginu ($\xi=0$, zakroužkované jsou podpůrné vektory přesně na marginu). Oranžová úsečka -- kladný bod uvnitř pásu na správné straně hranice ($0<\xi_i<1$). Červená úsečka -- bod záporné třídy zabloudil na kladnou stranu hranice ($\xi_j>1$, chyba). Slack $\xi_i$ je (škálovaná) kolmá vzdálenost od příslušného marginu k bodu; jeho součet $\sum_i\xi_i$ pokutuje konstanta $C$.}
\end{center}

\subsubsection*{Hinge loss: SVM jako regularizovaná ztráta}

Z~podmínek lze slack vyjádřit explicitně: $\xi_i=\max\big(0,\,1-t_i y(\mathbf{x}_i)\big)$ (je-li bod správně s~rezervou, je $1-t_iy\le 0$ a~slack je $0$; jinak je kladný). Dosazením dostaneme \emph{bezpodmínkovou} formulaci přes \textbf{hinge loss} (závěsovou ztrátu):
\[ \mathcal{L}_{\text{hinge}}(t,y) = \max(0,\,1-t\,y),\qquad \min_{\mathbf{w},b}\ \underbrace{C\sum_i \mathcal{L}_{\text{hinge}}\big(t_i,y(\mathbf{x}_i)\big)}_{\text{datová ztráta}} + \underbrace{\tfrac12\|\mathbf{w}\|^2}_{\text{regularizace}}. \]

\begin{intuice}
Hinge loss penalizuje bod, dokud není správně klasifikován \emph{s~rezervou} ($t_i y\ge 1$); za touto hranicí je ztráta nulová (proto „závěs``). To je rozdíl oproti logistické regresi (sekce~\ref{sec:logisticka}), jejíž ztráta $-\log\sigma(ty)$ klesá k~nule jen asymptoticky a~odměňuje i~body daleko za hranicí. Formulace ukazuje, že \textbf{$C$ je \emph{inverzní} síla regularizace}: $C\to\infty$ znamená žádnou regularizaci (hard-margin), $C\to 0$ ignoruje data. SVM tak zapadá do téhož rámce „datová ztráta $+$ regularizace`` jako lineární a~logistická regrese; liší se jen volbou ztrátové funkce.
\end{intuice}

\subsection{Jádrové funkce: kernel trick}

Lineární SVM zvládne jen lineárně (téměř) oddělitelná data. Trik, jak ho zmocnit pro nelineární hranice: data nejprve \emph{transformovat} zobrazením $\boldsymbol{\varphi}\colon\mathbb{R}^D\to\mathbb{R}^{F}$ do prostoru \emph{vyšší dimenze}, kde se stanou lineárně separabilní. Lineární hranice v~novém prostoru odpovídá \emph{zakřivené} hranici v~původním.

\begin{intuice}[Proč vyšší dimenze pomáhá]
Body, které v~rovině nejdou oddělit přímkou (např.\ jedna třída uvnitř kruhu, druhá kolem), se přidáním vhodných příznaků (např.\ $x_1^2,\,x_2^2,\,x_1x_2$) v~prostoru vyšší dimenze \uv{rozestoupí} tak, že je oddělí nadrovina. Vzdálenost od počátku ($x_1^2+x_2^2$) je nový příznak, podle nějž se vnitřní a~vnější kruh liší.
\end{intuice}

\begin{center}
\begin{tikzpicture}[scale=1.05]
  % --- levy panel: v R^2 nelze oddelit primkou (vnitrni shluk vs prstenec) ---
  \begin{scope}
    \draw[osa] (-2.3,0) -- (2.3,0) node[right,font=\scriptsize]{$x_1$};
    \draw[osa] (0,-2.3) -- (0,2.5) node[above,font=\scriptsize]{$x_2$};
    % vnitrni trida (cervene ctverce) v malem shluku kolem pocatku
    \foreach \x/\y in {0/0, 0.45/0.25, {0-0.4}/0.3, 0.3/{0-0.4}, {0-0.35}/{0-0.3}}{
      \node[cls2] at (\x,\y) {};}
    % vnejsi trida (modra kolecka) rovnomerne na kruznici (8 bodu po 45 stupnich)
    \foreach \a in {0,45,...,315}{
      \node[cls1] at ({1.8*cos(\a)},{1.8*sin(\a)}) {};}
    \node[font=\scriptsize] at (0,-2.85) {(a) v~$\mathbb{R}^2$: nejde oddělit přímkou};
  \end{scope}
  % sipka phi
  \draw[->,>=Stealth,thick,black!70] (2.7,0.1) -- (3.9,0.1) node[midway,above,font=\scriptsize]{$\boldsymbol{\varphi}$};
  % --- pravy panel: po transformaci oddelitelne (osa y = x1^2+x2^2) ---
  \begin{scope}[xshift=6.7cm]
    \draw[osa] (-2.3,0) -- (2.3,0) node[right,font=\scriptsize]{$x_1$};
    \draw[osa] (0,0) -- (0,3.6) node[above,font=\scriptsize]{$x_1^2+x_2^2$};
    % vnitrni trida: maly polomer -> mala vyska (kolem 0.2-0.3)
    \foreach \x/\h in {0/0.0, 0.45/0.27, {0-0.4}/0.25, 0.3/0.25, {0-0.35}/0.21}{
      \node[cls2] at (\x,\h) {};}
    % vnejsi trida: polomer 1.8 -> stejna vyska; x rozprostreno, aby se markery neprekryvaly
    \foreach \x in {-1.9,-1.35,-0.8,-0.25,0.3,0.85,1.4,1.95}{
      \node[cls1] at (\x,3.05) {};}
    % oddelujici nadrovina (vodorovna primka mezi shluky)
    \draw[hranice] (-2.1,1.6) -- (2.1,1.6);
    \node[font=\scriptsize,black!75,anchor=south] at (0,1.65) {nadrovina};
    \node[font=\scriptsize] at (0,-0.85) {(b) po transformaci: oddělí vodorovná nadrovina};
  \end{scope}
\end{tikzpicture}
\obrpopis{Kernel trik. (a)~Vnitřní třída (červené čtverce, malý shluk kolem počátku) a~vnější třída (modrá kolečka rovnoměrně na kružnici) nejdou v~rovině oddělit přímkou. (b)~Po transformaci do prostoru vyšší dimenze, kde je novým příznakem vzdálenost od počátku $x_1^2+x_2^2$, se třídy rozdělí a~oddělí obyčejnou (vodorovnou) nadrovinou; ta v~původní rovině odpovídá kružnici.}
\end{center}

Transformace do vysoké (i~nekonečné) dimenze by byla výpočetně drahá. Tady se uplatní, že duální SVM závisí na datech \emph{jen přes skalární součiny}: stačí umět spočítat $\boldsymbol{\varphi}(\mathbf{x})^{\!\top}\boldsymbol{\varphi}(\mathbf{x}')$, transformaci samotnou nikdy nepotřebujeme.

\begin{definice}[Jádro, kernel trick]
\textbf{Jádro} (kernel) odpovídající transformaci $\boldsymbol{\varphi}$ je funkce
\[ K(\mathbf{x},\mathbf{x}') \;\stackrel{\text{def}}{=}\; \boldsymbol{\varphi}(\mathbf{x})^{\!\top}\boldsymbol{\varphi}(\mathbf{x}'). \]
\textbf{Kernel trick}: v~duální úloze i~v~predikci $y(\mathbf{x})=\sum_i a_i t_i\,K(\mathbf{x},\mathbf{x}_i)+b$ všude nahradíme skalární součin $\mathbf{x}_i^{\!\top}\mathbf{x}_j$ jádrem $K(\mathbf{x}_i,\mathbf{x}_j)$. Pracujeme tak ve vysokodimenzionálním prostoru, \emph{aniž bychom $\boldsymbol{\varphi}$ kdy explicitně vyčíslili}.
\end{definice}

\begin{definice}[Nejčastější jádra]
\begin{itemize}
  \item \textbf{Polynomiální jádro} stupně (nejvýše) $d$:
    \[ K(\mathbf{x},\mathbf{x}') = (\gamma\,\mathbf{x}^{\!\top}\mathbf{x}' + 1)^d. \]
    Odpovídá transformaci na všechny součiny až $d$ vstupních příznaků (interakce do stupně $d$). Pro $d=2$ a~$\gamma=1$ je $\boldsymbol{\varphi}(x_1,x_2)=(1,\sqrt2\,x_1,\sqrt2\,x_2,x_1^2,\sqrt2\,x_1x_2,x_2^2)$.
  \item \textbf{Gaussovo (RBF) jádro} (radial basis function):
    \[ K(\mathbf{x},\mathbf{x}') = \exp\!\big(-\gamma\,\|\mathbf{x}-\mathbf{x}'\|^2\big),\qquad \gamma>0. \]
    Odpovídá skalárnímu součinu v~prostoru \emph{nekonečné} dimenze (kombinaci polynomiálních jader všech stupňů).
\end{itemize}
\end{definice}

\begin{intuice}[Kdy jádra pomohou a~jak se chová RBF]
Jádra pomohou, když hranice mezi třídami \emph{není lineární}. RBF jádro je funkcí \emph{vzdálenosti}: $K(\mathbf{x},\mathbf{x}')$ je velké pro blízké body a~klesá k~nule pro vzdálené. SVM s~RBF tak vlastně váží „svědectví`` trénovacích bodů podle vzdálenosti -- je to chytřejší varianta $k$-NN (sekce~\ref{sec:knn}), která bere v~úvahu všechny body, každý vážený podobností. Parametr $\gamma$ řídí \emph{dosah}: velké $\gamma$ = úzké špičky kolem bodů (riziko \emph{overfittingu}), malé $\gamma$ = široký vliv (hladší hranice). Na \emph{lineárně oddělitelných} datech jádro nepomůže -- naopak může vést k~\emph{overfittingu}; tam stačí lineární SVM.
\end{intuice}

\subsection{Klasifikace do více tříd}

SVM je ze své podstaty \emph{binární}. Pro $K>2$ tříd se skládá z~více binárních klasifikátorů dvěma standardními způsoby.

\begin{definice}[One-vs-rest a~one-vs-one]
\begin{itemize}
  \item \textbf{One-versus-rest} (jeden proti zbytku, OvR): natrénuje se $K$ binárních klasifikátorů; $i$-tý odděluje třídu $i$ od \emph{všech ostatních}. Při predikci vybereme třídu, jejíž klasifikátor dá nejvyšší skóre (vyžaduje \emph{srovnatelná} skóre, ideálně kalibrované pravděpodobnosti -- což SVM přímo nedává).
  \item \textbf{One-versus-one} (každý s~každým, OvO): natrénuje se $\binom{K}{2}$ klasifikátorů, jeden pro každou dvojici tříd $(i,j)$. Při predikci každý hlasuje pro jednu ze své dvojice a~vítězí třída s~\emph{nejvíce hlasy}. (Tuto variantu používá pro SVM např.\ knihovna \texttt{libsvm}/scikit-learn.)
\end{itemize}
\end{definice}

\begin{intuice}[Kompromis OvR vs.\ OvO]
OvR trénuje jen $K$ klasifikátorů (méně modelů), ale každý na \emph{nevyvážených} datech (jedna třída proti $K-1$) a~vyžaduje porovnatelná skóre. OvO trénuje $\binom{K}{2}=\tfrac{K(K-1)}{2}$ klasifikátorů (kvadraticky mnoho), zato každý na malé, vyvážené dvojici tříd; proto se hodí pro SVM, jehož trénink škáluje nadlineárně s~počtem bodů. Nevýhoda hlasování: vznikají \emph{nerozhodné} oblasti, kde víc tříd dostane stejně hlasů.
\end{intuice}

\begin{poznamka}[Nad rámec požadavků]
Jen pro úplnost a~bez detailu: duální úlohu SVM v~praxi řeší algoritmus \textbf{SMO} (Sequential Minimal Optimization), což je souřadnicový sestup optimalizující vždy dvojici multiplikátorů $a_i,a_j$ najednou (kvůli podmínce $\sum_i a_i t_i=0$). SVM lze upravit i~na \emph{regresi} (SVR) s~$\varepsilon$-necitlivou ztrátou. Tyto detaily požadavky nevyžadují.
\end{poznamka}

\subsection{Řešený příklad}

\begin{priklad}[SZZ léto 2024 č.~30: Metoda podpůrných vektorů (SVM)]
\szzotazka{léto 2024}{30}{Metoda podpůrných vektorů (SVM)}\par
\textbf{Zadání.} Data se dvěma příznaky, dvě třídy. Třída A (křížky, $t=+1$): $(1,1),(2,2),(2,0)$. Třída B (kolečka, $t=-1$): $(0,0),(1,0),(0,1)$. (1)~Popište, jak fungují a~trénují se lineární SVM pro lineárně separabilní data. (2)~Zakreslete přibližnou rozhodovací hranici lineárního SVM. (3)~Jak ovlivní hranici přidání bodu $(3,3)$ do třídy A? (4)~Pomohlo by na těchto datech RBF jádro? Nakreslete dataset, kde RBF pomohou.

\smallskip
\textbf{Řešení.}

\emph{(1)~Princip.} Viz výklad výše: SVM hledá oddělující přímku s~\emph{maximálním marginem}, tj.\ řeší QP úlohu $\min\tfrac12\|\mathbf{w}\|^2$ za podmínek $t_i(\mathbf{w}^{\!\top}\mathbf{x}_i+b)\ge 1$. Hranici určují \emph{podpůrné vektory} (body s~$a_i>0$); ty leží přesně na marginu, obráceně to platit nemusí.

\emph{(2)~Hranice.} Geometrie dat: třída B je v~„levém dolním`` rohu, třída A nad ní. Nejbližší body třídy A leží na přímce $x_1+x_2=2$, nejbližší body třídy B na $x_1+x_2=1$; separátor kolmý na směr $(1,1)$ uprostřed mezi nimi je zároveň maximálně marginový. Lepší být nemůže: bod $(1,1)$ třídy A má od úsečky $(1,0)$--$(0,1)$, která leží v~konvexním obalu třídy B, vzdálenost $1/\sqrt2$, a~každá oddělující přímka musí ležet mezi nimi, takže žádná nemá margin větší než $1/(2\sqrt2)$. Tento separátor ho dosahuje (ověřeno i~řešičem, viz níže). Je to přímka
\[ \boxed{x_1+x_2 = 1{,}5},\qquad \text{tj.\ } \mathbf{w}=(1,1),\ b=-1{,}5 \]
(v~kanonické normalizaci $\mathbf{w}=(2,2),\ b=-3$, aby nejbližší body měly $|y|=1$). Na marginu leží $(1,1)$ a~$(2,0)$ z~A a~$(1,0),(0,1)$ z~B. Jako \textbf{podpůrné vektory} (s~$a_i>0$) vrátí řešič $(1,1)$, $(1,0)$ a~$(0,1)$; bod $(2,0)$ je na marginu také, ale k~určení optima už není potřeba (degenerovaný případ). Při kanonické normalizaci $\mathbf{w}=(2,2)$ je $\|\mathbf{w}\|=2\sqrt2$, takže margin (vzdálenost separátoru k~nejbližšímu bodu) je $\dfrac{1}{\|\mathbf{w}\|}=\dfrac{1}{2\sqrt2}\approx 0{,}354$ a~celková šíře pásu mezi marginy $\dfrac{2}{\|\mathbf{w}\|}=\dfrac{1}{\sqrt2}\approx 0{,}707$.

\begin{center}
\begin{tikzpicture}[scale=1.25]
  \draw[osa] (-0.4,0) -- (3.1,0) node[right,font=\scriptsize]{$x_1$};
  \draw[osa] (0,-0.4) -- (0,3.1) node[above,font=\scriptsize]{$x_2$};
  \foreach \t in {1,2,3}{\draw[black!25] (\t,-0.05)--(\t,0.05); \node[font=\tiny,anchor=north] at (\t,-0.05){\t};
                          \draw[black!25] (-0.05,\t)--(0.05,\t); \node[font=\tiny,anchor=east] at (-0.05,\t){\t};}
  % decision boundary x1+x2=1.5
  \draw[hranice] (-0.2,1.7) -- (1.7,-0.2);
  % popisek hranice v otevrene casti nad pasem, mimo body i osove popisky
  \node[font=\scriptsize,black!75,anchor=south west] at (1.15,1.55) {$x_1{+}x_2{=}1{,}5$};
  % margins: x1+x2=2 (A side) and x1+x2=1 (B side)
  \draw[margin] (0.0,2.0) -- (2.0,0.0);
  \draw[margin] (0.0,1.0) -- (1.0,0.0);
  % class A: crosses (use cls2 = squares to match preamble; mark with x via node text)
  \node[cls2,podpora] at (1,1) {}; % support vector (on margin x1+x2=2)
  \node[cls2] at (2,2) {};
  \node[cls2] at (2,0) {};         % lies ON margin x1+x2=2 but not chosen as SV (degenerate)
  % class B: circles
  \node[cls1] at (0,0) {};
  \node[cls1,podpora] at (1,0) {}; % support vector (on margin x1+x2=1)
  \node[cls1,podpora] at (0,1) {}; % support vector
  % labels
  \node[font=\scriptsize,red!65!black,anchor=south west] at (2.05,2.0) {A ($t{=}{+}1$)};
  \node[font=\scriptsize,blue!60!black,anchor=north east] at (-0.05,-0.05) {B ($t{=}{-}1$)};
\end{tikzpicture}
\obrpopis{Skutečný max-margin separátor pro data ze zadání (ověřeno \texttt{sklearn}): plná čára $x_1+x_2=1{,}5$, čárkované marginy $x_1+x_2=2$ (strana A) a~$x_1+x_2=1$ (strana B). Zakroužkované podpůrné vektory: $(1,1)$ z~A a~$(1,0),(0,1)$ z~B. Body $(2,2)$ a~$(0,0)$ leží daleko a~hranici neovlivňují; $(2,0)$ leží sice na marginu, ale optimum je určeno už třemi vektory.}
\end{center}

\emph{(3)~Přidání bodu $(3,3)$ do A.} Bod $(3,3)$ leží \emph{hluboko} na správné straně (pro něj $x_1+x_2=6\gg 2$), daleko od marginu. Není to podpůrný vektor, jeho odebrání či přidání hranici nezmění: \textbf{separátor i~margin zůstanou stejné} ($\mathbf{w}=(1,1),\ b=-1{,}5$). To je klíčová vlastnost SVM -- hranici určují jen body na marginu, ne vzdálené body (na rozdíl třeba od lineární regrese, kterou by odlehlý bod „přitáhl``). \emph{(Ověřeno numericky: po přidání $(3,3)$ vrací \texttt{sklearn} týž separátor a~$(3,3)$ mezi podpůrnými vektory není.)}

\emph{(4)~RBF na těchto datech.} Data \emph{jsou} lineárně separabilní (právě jsme našli přímku), takže RBF jádro by zde \textbf{nepomohlo} -- spíš by hrozil \emph{overfitting} (přeučení). RBF se vyplatí, až když lineární hranice nestačí. Příklad takových dat -- \emph{soustředné kruhy}:

\begin{center}
\begin{tikzpicture}[scale=1.0]
  \draw[osa] (-2.4,0) -- (2.4,0) node[right,font=\scriptsize]{$x_1$};
  \draw[osa] (0,-2.4) -- (0,2.6) node[above,font=\scriptsize]{$x_2$};
  % vnitrni trida (cervene ctverce) -- maly shluk kolem pocatku (stejny vzor jako u kernel triku)
  \foreach \x/\y in {0/0, 0.45/0.25, {0-0.4}/0.3, 0.3/{0-0.4}, {0-0.35}/{0-0.3}}{
    \node[cls2] at (\x,\y) {};}
  % vnejsi trida (modra kolecka) -- rovnomerne na kruznici (8 bodu po 45 stupnich)
  \foreach \a in {0,45,...,315}{
    \node[cls1] at ({1.8*cos(\a)},{1.8*sin(\a)}) {};}
  % carkovana RBF hranice (kruznice mezi shlukem a prstencem)
  \draw[margin] (0,0) circle (1.15);
  \node[font=\scriptsize,black!70,anchor=south west] at (0.82,0.82) {hranice (RBF)};
\end{tikzpicture}
\obrpopis{Dataset, kde RBF jádro pomáhá: jedna třída uvnitř (čtverce, malý shluk kolem počátku), druhá v~prstenci kolem (kolečka rovnoměrně na kružnici). Žádná přímka je neoddělí, RBF jádro ano, kruhovou hranicí. (Na vygenerovaných datech tohoto typu dosáhne lineární SVM jen $\approx 63\,\%$ správnosti, RBF $100\,\%$; ověřeno \texttt{sklearn}.)}
\end{center}
\end{priklad}

\begin{poznamka}[Numerické ověření]
Skutečný separátor, podpůrné vektory, neměnnost při přidání $(3,3)$ i~rozdíl lineární vs.\ RBF na kruzích jsou ověřeny skriptem \texttt{src/fig/gen\_s7\_svm\_verify.py} (\texttt{sklearn.svm.SVC}). Výsledek: $\mathbf{w}=(2,2),\ b=-3$ (kanonicky), margin $1/(2\sqrt2)\approx0{,}354$ (šíře pásu $1/\sqrt2$), podpůrné vektory $(1,1),(1,0),(0,1)$; lineární SVM na soustředných kruzích $63\,\%$, RBF $100\,\%$.
\end{poznamka}

\section{Kombinace více modelů}\label{sec:kombinace}

Jeden model má svá omezení: lineární modely nezachytí nelinearity, jeden hluboký rozhodovací strom má vysoký \emph{rozptyl} (drobná změna dat dá úplně jiný strom; viz sekce~\ref{sec:stromy}). \textbf{Kombinace více modelů} (ensembling) tyto slabiny obchází tím, že více modelů necháme rozhodovat společně. Tato sekce pokrývá tři klasické přístupy: \textbf{bagging} (paralelní modely na bootstrapových vzorcích, snižuje rozptyl), \textbf{boosting} (sekvenční modely opravující chyby předchůdců) a~jejich nejznámější zástupce -- \textbf{random forest} (náhodný les), tj.\ bagging rozhodovacích stromů s~náhodným výběrem příznaků.

\begin{intuice}
Proč vůbec kombinovat? Představme si výbor expertů: pokud každý dělá chyby \emph{jinde} (jejich chyby nejsou korelované), pak hlasováním nebo průměrem se chyby navzájem vyruší a~výbor je lepší než kterýkoli jednotlivý expert. Naopak když všichni dělají \emph{tytéž} chyby, kombinace nepomůže. Celé umění ensemblingu je vyrobit modely, které jsou jednotlivě slušné, ale jejich chyby spolu \emph{nesouvisí}.
\end{intuice}

\begin{definice}[Kombinace modelů: hlasování a~průměrování]
Mějme $M$ natrénovaných modelů $y_1,\dots,y_M$. Jejich predikce zkombinujeme do jediné:
\begin{itemize}
  \item \textbf{hard voting} (tvrdé hlasování) -- u~klasifikace vrátíme třídu, kterou predikuje \emph{nejvíce} modelů (většina);
  \item \textbf{soft voting} (měkké hlasování / averaging) -- zprůměrujeme \emph{rozdělení pravděpodobností} tříd vracená modely a~vybereme třídu s~nejvyšší průměrnou pravděpodobností; u~regrese prostě zprůměrujeme číselné predikce.
\end{itemize}
\end{definice}

\subsection{Bagging (bootstrap agregace)}

\begin{definice}[Bootstrap, bagging]
\textbf{Bootstrapový vzorek} trénovací množiny o~$N$ příkladech vznikne tak, že z~ní vylosujeme $N$ příkladů \emph{s~opakováním} (s~vracením). Vzniklá množina má stejnou velikost $N$, ale některé příklady obsahuje vícekrát a~jiné vůbec. \textbf{Bagging} (\emph{bootstrap aggregating}) natrénuje $M$ kopií téhož modelu, každou na vlastním nezávislém bootstrapovém vzorku, a~jejich predikce \emph{agreguje} -- u~regrese průměrem, u~klasifikace hlasováním.
\end{definice}

\begin{intuice}
Losování s~opakováním vyrobí $M$ \emph{podobných, ale ne identických} trénovacích množin, takže vzniklé modely jsou různé. Bagging hlavně \textbf{snižuje rozptyl} (variance) odhadu, aniž zvýší vychýlení (bias): je proto ideální pro modely s~\emph{nízkým vychýlením a~vysokým rozptylem} -- typicky hluboké rozhodovací stromy. Zhruba $1/e\approx 37\,\%$ příkladů se do daného bootstrapového vzorku nedostane (tzv.\ \emph{out-of-bag} příklady); na nich lze model „zdarma`` validovat.
\end{intuice}

\begin{poznamka}[Proč právě bootstrap]
Pro neuronové sítě často stačí trénovat modely s~různou náhodnou inicializací: ztráta má mnoho lokálních minim, takže modely vyjdou dost odlišné i~na stejných datech. U~modelů s~\emph{konvexní} ztrátovou funkcí (lineární či logistická regrese) ale učení na týchž datech konverguje k~témuž optimu bez ohledu na náhodnost (inicializaci, pořadí příkladů). Různé modely tu dostaneme jen změnou samotných trénovacích dat -- a~to dělá bagging.
\end{poznamka}

\begin{center}
\begin{tikzpicture}[font=\small,
  ds/.style={draw,thick,rounded corners=2pt,fill=blue!8,minimum width=20mm,minimum height=8mm,align=center},
  bs/.style={draw,thick,rounded corners=2pt,fill=yellow!18,minimum width=20mm,minimum height=7mm,align=center},
  md/.style={draw,thick,rounded corners=2pt,fill=green!12,minimum width=20mm,minimum height=7mm,align=center},
  ag/.style={draw,thick,rounded corners=2pt,fill=red!10,minimum width=22mm,minimum height=9mm,align=center},
  fl/.style={->,>=Stealth,thick}]
  \node[ds] (D) at (0,0) {trénovací\\množina $\mathcal{D}$};
  \node[bs] (B1) at (3.6,2.0) {bootstrap $\mathcal{D}_1$};
  \node[bs] (B2) at (3.6,0.0) {bootstrap $\mathcal{D}_2$};
  \node[bs] (B3) at (3.6,-2.0) {bootstrap $\mathcal{D}_M$};
  \node[md] (M1) at (7.0,2.0) {model $y_1$};
  \node[md] (M2) at (7.0,0.0) {model $y_2$};
  \node[md] (M3) at (7.0,-2.0) {model $y_M$};
  \node[ag] (A) at (10.6,0.0) {agregace\\(hlasování\\/ průměr)};
  \node[font=\scriptsize] at (3.6,-1.0) {\vdots};
  \node[font=\scriptsize] at (7.0,-1.0) {\vdots};
  \foreach \b in {B1,B2,B3}{\draw[fl] (D) -- (\b);}
  \draw[fl] (B1) -- (M1); \draw[fl] (B2) -- (M2); \draw[fl] (B3) -- (M3);
  \foreach \m in {M1,M2,M3}{\draw[fl] (\m) -- (A);}
  \node[font=\scriptsize,below=1pt of D] {\textit{losování s~opakováním}};
\end{tikzpicture}
\obrpopis{Schéma baggingu. Z~jediné trénovací množiny $\mathcal{D}$ vzniká $M$ bootstrapových vzorků (losování $N$ příkladů \emph{s~opakováním}). Na každém se natrénuje samostatný model; jejich predikce se na konci agregují (hlasováním u~klasifikace, průměrem u~regrese). Modely jsou trénované \emph{paralelně} a~nezávisle.}
\end{center}

\subsubsection*{Proč průměrování snižuje chybu}

\begin{veta}[Redukce chyby průměrováním (nekorelované chyby)]
Označme chybu $i$-tého modelu na příkladu $\mathbf{x}$ jako $\varepsilon_i(\mathbf{x})=y_i(\mathbf{x})-t$. Mají-li chyby \emph{nulovou střední hodnotu} a~jsou \emph{nekorelované} ($\mathbb{E}[\varepsilon_i\varepsilon_j]=0$ pro $i\ne j$), pak střední kvadratická chyba průměru $M$ modelů je
\[ \mathbb{E}\!\left[\Big(\tfrac{1}{M}\textstyle\sum_i \varepsilon_i(\mathbf{x})\Big)^2\right] = \frac{1}{M}\,\mathbb{E}\!\left[\tfrac{1}{M}\textstyle\sum_i \varepsilon_i^2(\mathbf{x})\right], \]
tj.\ $\tfrac{1}{M}$-násobek průměrné chyby jednotlivých modelů.
\end{veta}

\begin{intuice}
Roznásobíme-li $\big(\sum_i\varepsilon_i\big)^2=\sum_{i,j}\varepsilon_i\varepsilon_j$ a~vezmeme střední hodnotu, smíšené členy $\mathbb{E}[\varepsilon_i\varepsilon_j]$ ($i\ne j$) vypadnou (nekorelovanost) a~zbude jen $M$ diagonálních členů $\mathbb{E}[\varepsilon_i^2]$; faktor $\tfrac{1}{M^2}$ před sumou pak dá výsledný $\tfrac{1}{M}$. Toto je \emph{ideální} mez: chyba klesá až $M$-krát. V~praxi jsou chyby \emph{korelované} (modely se učí z~překrývajících se dat), takže zlepšení je menší. Předpoklad nekorelovanosti je tedy klíčový -- proto se bagging snaží modely co nejvíc „rozhodit``.
\end{intuice}

\subsection{Boosting a~algoritmus AdaBoost}

\begin{definice}[Boosting]
\textbf{Boosting} buduje silný klasifikátor jako \emph{vážený součet} mnoha \emph{slabých} klasifikátorů (např.\ mělkých stromů), trénovaných \textbf{sekvenčně}. Po každém kole se zvýší \emph{váhy} těch trénovacích příkladů, které dosavadní kombinace klasifikuje \emph{špatně}, takže další slabý klasifikátor se zaměří právě na ně.
\end{definice}

\begin{intuice}
Zatímco bagging trénuje modely \emph{paralelně a~nezávisle} (cíl: snížit rozptyl), boosting je trénuje \emph{postupně a~závisle} -- každý nový model „dohání`` chyby předchozích (cíl: snížit vychýlení). Slabý klasifikátor stačí lepší než náhoda (chybovost $<\tfrac12$); jejich vážená kombinace může být libovolně přesná. Metafora: tým, kde se každý nováček soustředí na úlohy, které tým dosud nezvládal.
\end{intuice}

\begin{algo}[AdaBoost (binární klasifikace, třídy $t_i\in\{-1,+1\}$)]
Vstup: trénovací data $(\mathbf{x}_i,t_i)$, $i=1,\dots,N$; počet kol $T$.\\[2pt]
Inicializuj váhy příkladů rovnoměrně: $w_i^{(1)}=\tfrac1N$. \cmt{na začátku stejně důležité}\\
\textbf{pro} $m=1,\dots,T$:
\begin{enumerate}[leftmargin=2.4em,itemsep=1pt]
  \item Natrénuj slabý klasifikátor $h_m(\mathbf{x})\in\{-1,+1\}$ na datech s~váhami $w^{(m)}$.
  \item Spočti váženou chybu \;$\varepsilon_m=\sum_{i:\,h_m(\mathbf{x}_i)\ne t_i} w_i^{(m)}$. \cmt{podíl chybně klasifikované váhy}
  \item Spočti váhu klasifikátoru \;$\alpha_m=\tfrac12\ln\!\dfrac{1-\varepsilon_m}{\varepsilon_m}$. \cmt{čím menší chyba, tím větší vliv}
  \item Převáž příklady: \;$w_i^{(m+1)}\propto w_i^{(m)}\,e^{-\alpha_m t_i h_m(\mathbf{x}_i)}$ a~normalizuj na součet $1$. \cmt{špatně klasifikované získají větší váhu}
\end{enumerate}
Výstup: \;$H(\mathbf{x})=\sign\!\Big(\sum_{m=1}^{T}\alpha_m\,h_m(\mathbf{x})\Big)$. \cmt{vážené hlasování slabých klasifikátorů}
\end{algo}

\begin{intuice}
V~kroku~4 je $t_i h_m(\mathbf{x}_i)=+1$ pro správně a~$-1$ pro špatně klasifikovaný příklad, takže $e^{-\alpha_m t_i h_m}$ je $<1$ pro správné (váha klesne) a~$>1$ pro chybné (váha vzroste). Příklady, na kterých model selhává, tak v~dalším kole „křičí`` hlasitěji. Klasifikátor s~malou chybou $\varepsilon_m$ dostane velké $\alpha_m$ (velký hlas); klasifikátor s~$\varepsilon_m=\tfrac12$ (náhoda) dostane $\alpha_m=0$ (žádný hlas).
\end{intuice}

\begin{center}
\begin{tikzpicture}[font=\small,
  box/.style={draw,thick,rounded corners=2pt,minimum width=22mm,minimum height=15mm,align=center,fill=blue!7},
  fl/.style={->,>=Stealth,thick}]
  \node[box] (h1) at (0,0) {slabý\\klasif.\ $h_1$};
  \node[box] (h2) at (3.5,0) {slabý\\klasif.\ $h_2$};
  \node[box] (h3) at (7.0,0) {slabý\\klasif.\ $h_3$};
  \node (dots) at (9.3,0) {\dots};
  \draw[fl] (h1) -- node[above,font=\scriptsize,align=center]{převáž\\chyby} (h2);
  \draw[fl] (h2) -- node[above,font=\scriptsize,align=center]{převáž\\chyby} (h3);
  \draw[fl] (h3) -- (dots);
  % body s rostouci vahou pod kazdym klasifikatorem
  \def\yb{-1.5}
  \foreach \cx/\sa/\sb in {0/1.0/1.0, 3.5/0.6/1.8, 7.0/0.5/2.4}{
    \fill[green!55!black] (\cx-0.7,\yb) circle (\sa pt);
    \fill[green!55!black] (\cx-0.35,\yb) circle (\sa pt);
    \fill[red!70!black]   (\cx,\yb)      circle (\sb pt);
    \fill[red!70!black]   (\cx+0.35,\yb) circle (\sb pt);
    \fill[green!55!black] (\cx+0.7,\yb)  circle (\sa pt);
  }
  \node[font=\scriptsize,anchor=west] at (8.0,\yb) {velikost $=$ váha};
  \node[font=\scriptsize,anchor=north] at (0,\yb-0.25) {rovnoměrné};
  \node[font=\scriptsize,anchor=north] at (3.5,\yb-0.25) {chybné $\uparrow$};
  \node[font=\scriptsize,anchor=north] at (7.0,\yb-0.25) {chybné $\uparrow\uparrow$};
\end{tikzpicture}
\obrpopis{Schéma boostingu. Slabé klasifikátory se trénují \emph{sekvenčně}; po každém kole vzroste váha (velikost bodu) příkladů, které dosud klasifikujeme špatně (červeně), takže další klasifikátor se zaměří na ně. Finální predikce je \emph{vážené hlasování} všech slabých klasifikátorů.}
\end{center}

\begin{poznamka}[Gradientní boosting a~XGBoost (jen koncept)]
Modernější varianta, \textbf{gradientní boosting} (GBDT), místo převažování příkladů trénuje každý nový strom tak, aby predikoval \emph{záporný gradient ztráty} dosavadní kombinace (tj.\ „opravoval rezidua`` předchozích stromů); přidává se s~malým koeficientem (learning rate). Efektivní implementace \textbf{XGBoost} a~\textbf{LightGBM} (využívají i~druhý řád ztráty a~chytré dělení uzlů) patří k~nejsilnějším modelům pro \emph{tabulární data}. Detailní matematiku (Newtonův krok, Hessián) zde nerozebíráme; podstatná je myšlenka „sekvenční stromy opravující chybu předchůdců``.
\end{poznamka}

\subsection{Random forest (náhodný les)}

\begin{definice}[Random forest]
\textbf{Random forest} (náhodný les) je \emph{bagging rozhodovacích stromů} s~jednou přidanou náhodností: při hledání nejlepšího řezu v~\emph{každém uzlu} se uvažuje jen \textbf{náhodná podmnožina příznaků} (ne všechny). Predikce random forest vznikne \emph{hlasováním} stromů (klasifikace), resp.\ \emph{průměrem} (regrese).
\end{definice}

\begin{algo}[Trénink random forest]
Vstup: data $\mathcal{D}$ o~$N$ příkladech s~$D$ příznaky; počet stromů $M$; velikost podmnožiny příznaků $k$ (typicky $k\approx\sqrt{D}$ u~klasifikace).\\[2pt]
\textbf{pro} $m=1,\dots,M$:
\begin{enumerate}[leftmargin=2.4em,itemsep=1pt]
  \item Vytvoř bootstrapový vzorek $\mathcal{D}_m$ (losuj $N$ příkladů s~opakováním). \cmt{bagging}
  \item Natrénuj rozhodovací strom na $\mathcal{D}_m$, ale v~\emph{každém uzlu} vyber nejlepší řez jen z~\emph{náhodných $k$ příznaků} (čerstvá podmnožina pro každý uzel). \cmt{náhodné příznaky}
  \item Strom \emph{neprořezávej} (necháváme přerůst -- rozptyl srovná až agregace).
\end{enumerate}
Predikce: \;hlasování (klasifikace) / průměr (regrese) přes všech $M$ stromů.
\end{algo}

\begin{intuice}
Proč navíc \emph{náhodné příznaky}? Samotný bagging stromů by dal stromy stále dost podobné -- pokud existuje jeden velmi silný příznak, vybraly by ho téměř všechny stromy v~kořeni a~jejich chyby by zůstaly korelované. Omezení na náhodnou podmnožinu příznaků nutí různé stromy „dívat se jinam``, čímž \emph{dekoreluje} jejich chyby; podle věty výše to je přesně to, co dělá agregaci účinnou. (Ještě radikálnější jsou \emph{extra trees}: nehledají ani nejlepší hodnotu řezu, volí ji náhodně z~rozsahu příznaku.)
\end{intuice}

\begin{poznamka}[Random forest vs.\ jeden strom -- výhody a~nevýhody]
\textbf{Jeden strom:} rychlý trénink, snadná \emph{interpretace} (cestu k~rozhodnutí lze přečíst), ale vysoký \emph{rozptyl} a~sklon k~overfittingu. \textbf{Random forest:} výrazně \emph{přesnější a~stabilnější} (lépe generalizuje, odolný vůči overfittingu), zvládá vysoký počet příznaků, dává „zdarma`` odhad chyby (out-of-bag) i~důležitost příznaků. Cena: \emph{ztráta interpretovatelnosti} (stovky stromů už člověk nepřečte), pomalejší a~paměťově náročnější trénink i~inference. Random forest a~gradientní boosting stromů patří k~nejlepším modelům pro tabulární data (viz příklad v~sekci~\ref{sec:evaluace}).
\end{poznamka}

\subsection{Řešené příklady}

\begin{priklad}[SZZ jaro 2023 č.~6: Učení pomocí kolektivních prediktorů]
\szzotazka{jaro 2023}{6}{Učení pomocí kolektivních prediktorů}\par
\textbf{Zadání.} (1)~Vysvětlete princip trénování metodami \emph{bagging} a~\emph{boosting}. (2)~Popište trénovací algoritmus \emph{random forest} (náhodných lesů). (3)~Uveďte konkrétní příklad algoritmu typu \emph{boosting} a~stručně vysvětlete jeho hlavní myšlenku.

\smallskip
\textbf{Řešení.}
\emph{(1)} \textbf{Bagging} trénuje $M$ kopií modelu \emph{paralelně}, každou na vlastním \emph{bootstrapovém} vzorku (losování $N$ příkladů s~opakováním); predikce agreguje hlasováním/průměrem. Snižuje \emph{rozptyl} (vhodný pro nestabilní modely, např.\ hluboké stromy). \textbf{Boosting} trénuje slabé klasifikátory \emph{sekvenčně}; po každém kole zvýší \emph{váhu} chybně klasifikovaných příkladů, takže další model se zaměří na ně. Finální model je vážené hlasování. Snižuje hlavně \emph{vychýlení}.

\emph{(2)} \textbf{Random forest} $=$ bagging rozhodovacích stromů $+$ náhodná podmnožina příznaků v~každém uzlu. Pro $m=1,\dots,M$: vytvoř bootstrapový vzorek, natrénuj na něm (nezprořezaný) strom, ale v~každém uzlu hledej nejlepší řez jen z~náhodně zvolené podmnožiny příznaků (typicky $\sqrt{D}$). Predikce: hlasování stromů. Dvojí náhodnost (data $+$ příznaky) \emph{dekoreluje} stromy.

\emph{(3)} \textbf{AdaBoost.} Slabé klasifikátory $h_m\in\{-1,+1\}$ se trénují postupně na vážených datech; klasifikátor s~váženou chybou $\varepsilon_m$ dostane hlas $\alpha_m=\tfrac12\ln\tfrac{1-\varepsilon_m}{\varepsilon_m}$ a~váhy chybně klasifikovaných příkladů se vynásobí $e^{\alpha_m}$ (vzrostou). Výsledek $H(\mathbf{x})=\sign\big(\sum_m \alpha_m h_m(\mathbf{x})\big)$. Hlavní myšlenka: \emph{každé další kolo se soustředí na to, co dosavadní kombinace ještě nezvládá}.
\end{priklad}

\begin{priklad}[SZZ podzim 2023 č.~30: Kombinace více modelů]
\szzotazka{podzim 2023}{30}{Kombinace více modelů}\par
\textbf{Zadání.} (1)~Popište myšlenky baggingu a~boostingu; jeden boosting algoritmus podrobněji. (2)~Jak funguje a~jak se trénuje random forest? (3)~Porovnejte rozhodovací stromy a~random forest (výhody/nevýhody).

\smallskip
\textbf{Řešení.}
\emph{(1)} Viz definice baggingu a~boostingu výše a~algoritmus \textbf{AdaBoost} v~předchozím příkladu (slabé klasifikátory, převažování chyb, vážené hlasování $\alpha_m$).

\emph{(2)} Viz algoritmus tréninku random forest výše: každý strom na bootstrapovém vzorku, v~každém uzlu řez jen z~náhodné podmnožiny příznaků; výsledek hlasováním stromů.

\emph{(3)} \textbf{Stromy:} interpretovatelné, rychlé, ale vysoký rozptyl a~overfitting. \textbf{Random forest:} přesnější, stabilnější, lépe generalizuje, OOB odhad chyby a~důležitost příznaků zdarma; cena je ztráta interpretovatelnosti a~vyšší výpočetní nároky. (Viz poznámka „random forest vs.\ strom`` výše.)
\end{priklad}

\begin{priklad}[SZZ jaro 2026 č.~29: Meze accuracy ensemblu při hlasování]
\szzotazka{jaro 2026}{29}{Kombinace modelů}\par
\textbf{Zadání.} (1)~Stručně popište myšlenky baggingu a~boostingu a~uveďte modely, které je používají. (2)~Tři modely $M_1,M_2,M_3$ pro binární klasifikaci mají accuracy $0{,}70$, $0{,}75$, $0{,}80$. Vytvoříme ensemble \emph{většinovým hlasováním} (finální odpověď je ta častější ze tří). \emph{Za jakých podmínek} bude mít ensemble vyšší accuracy než nejlepší model? \emph{Jaká je nejlepší dosažitelná} accuracy ensemblu? \emph{Kdy} bude naopak horší?

\smallskip
\textbf{Řešení.}
\emph{(1)} \textbf{Bagging} -- paralelní modely na bootstrapových vzorcích, agregace hlasováním/průměrem; příklad: \emph{random forest}. \textbf{Boosting} -- sekvenční modely, převažování chyb; příklad: \emph{AdaBoost}, \emph{XGBoost} (gradientní boosting).

\emph{(2)} Klíč je uvažovat \textbf{chybové množiny} $E_i=\{$vzorky, které $M_i$ klasifikuje špatně$\}$. Z~přesností: $|E_1|=30\,\%$, $|E_2|=25\,\%$, $|E_3|=20\,\%$ z~testovací množiny. Při \emph{většinovém hlasování tří} modelů je vzorek klasifikován \textbf{špatně}, právě když chybují \emph{alespoň dva} modely, tj.\ vzorek leží \emph{alespoň ve dvou} z~množin $E_1,E_2,E_3$. Chyba ensemblu je tedy
\[ \mathrm{err}_{\mathrm{ens}} = \big|\,\text{vzorky pokryté alespoň dvěma z~}E_1,E_2,E_3\,\big|. \]
Vše dál plyne z~toho, \emph{jak se chybové množiny překrývají}.

\smallskip
\begin{center}
\begin{tikzpicture}[scale=1.15,font=\small]
  \def\r{1.6}
  \coordinate (c1) at (-0.95,0.55);
  \coordinate (c2) at (0.95,0.55);
  \coordinate (c3) at (0,-1.05);
  % výplně kruhů (poloprůhledné)
  \fill[blue!20,opacity=0.5]  (c1) circle (\r);
  \fill[red!20,opacity=0.5]   (c2) circle (\r);
  \fill[green!22,opacity=0.5] (c3) circle (\r);
  % vyšrafování oblasti "chybují >= 2 modely" = sjednocení tří párových průniků
  \begin{scope}
    \clip (c1) circle (\r); \fill[pattern=north east lines,pattern color=black!55] (c2) circle (\r);
  \end{scope}
  \begin{scope}
    \clip (c1) circle (\r); \fill[pattern=north east lines,pattern color=black!55] (c3) circle (\r);
  \end{scope}
  \begin{scope}
    \clip (c2) circle (\r); \fill[pattern=north east lines,pattern color=black!55] (c3) circle (\r);
  \end{scope}
  % obrysy kruhů
  \draw[blue!55!black,thick]  (c1) circle (\r);
  \draw[red!55!black,thick]   (c2) circle (\r);
  \draw[green!45!black,thick] (c3) circle (\r);
  % popisky množin (vně kruhů)
  \node[blue!55!black]  at (-2.35,2.05) {$E_1$ ($30\%$)};
  \node[red!55!black]   at (2.35,2.05)  {$E_2$ ($25\%$)};
  \node[green!45!black] at (0,-2.95)    {$E_3$ ($20\%$)};
  % anotace vyšrafované oblasti s vodicí šipkou na trojný střed
  \node[font=\scriptsize,align=center,anchor=west] (anot) at (2.6,-0.2)
    {šrafování: chybují\\$\ge 2$ modely\\\textit{($=$ chyba ensemblu)}};
  \draw[->,>=Stealth,black!60] (anot.west) to[bend right=12] (0,-0.05);
\end{tikzpicture}
\obrpopis{Chybové množiny tří modelů. Vzorek je ensemblem klasifikován špatně, právě když leží alespoň ve dvou množinách (vyšrafovaná oblast tvořená třemi párovými průniky). Disjunktní množiny (žádný překryv) $\Rightarrow$ ensemble se nikdy nemýlí; velké překryvy $\Rightarrow$ ensemble se mýlí často.}
\end{center}

\textbf{Nejlepší dosažitelná accuracy.} Mají-li modely \emph{disjunktní} chybové množiny (každý chybuje jinde), žádný vzorek neleží ve dvou množinách, takže $\mathrm{err}_{\mathrm{ens}}=0$, tj.\ \emph{accuracy $100\,\%$}. To je možné, protože $30\%+25\%+20\%=75\%\le 100\%$ se do testovací množiny vejde disjunktně.
\[ \boxed{\text{nejlepší dosažitelná accuracy ensemblu} = 100\,\%\quad(\text{disjunktní chyby}).} \]

\textbf{Kdy je ensemble lepší než nejlepší model.} Nejlepší model $M_3$ má chybu $20\,\%$. Ensemble je lepší, právě když $\mathrm{err}_{\mathrm{ens}} < 20\,\%$, tj.\ když se chybové množiny překrývají jen málo (chyby modelů jsou málo korelované).
Stačí tedy, aby modely chybovaly převážně na \emph{různých} vzorcích.

\textbf{Kdy je ensemble horší.} Při \emph{silně korelovaných} (překrývajících se) chybách:
\begin{itemize}
  \item je-li $E_2\subseteq E_1$ (kdykoli chybuje $M_2$, chybuje i~$M_1$), je na celém $E_2$ chyba $\ge 2$, takže $\mathrm{err}_{\mathrm{ens}}\ge 25\%$, tj.\ accuracy \emph{nejvýše} $0{,}75$;
  \item je-li navíc i~zbytek $E_1\setminus E_2$ pokryt $E_3$, je celé $E_1$ pokryté dvěma množinami, $\mathrm{err}_{\mathrm{ens}}\ge 30\%$, tj.\ accuracy \emph{nejvýše} $0{,}70$ (ensemble není lepší ani než nejhorší model).
\end{itemize}
Vůbec \emph{nejhorší} dosažitelná hodnota (maximální překryv při daných velikostech) je $1-\tfrac{0{,}30+0{,}25+0{,}20}{2}=\mathbf{0{,}625}$: každý vzorek, na kterém ensemble chybuje, spotřebuje aspoň dvě chyby jednotlivých modelů, a~těch je dohromady $75\,\%$. Mez se dosáhne např.\ při $|E_1\cap E_2|=17{,}5\%$, $|E_1\cap E_3|=12{,}5\%$, $|E_2\cap E_3|=7{,}5\%$ (trojný průnik prázdný): aspoň dvěma množinami je pokryto $37{,}5\%$ vzorků a~accuracy klesne na $\boxed{0{,}625}$. (Mají-li chyby být celá procenta, tj.\ $N=100$, je nejhorší $0{,}63$, např.\ při $17/12/8\,\%$.) Ensemble tedy může být horší \emph{než i~nejhorší jednotlivý model}, pokud modely dělají tytéž chyby.

\smallskip
\noindent\textbf{Závěr.} Většinové hlasování pomáhá jen tehdy, jsou-li chyby modelů málo korelované: při \emph{nezávislých} chybách chybují právě dva modely s~pravděpodobností $0{,}3\cdot0{,}25\cdot0{,}8+0{,}3\cdot0{,}75\cdot0{,}2+0{,}7\cdot0{,}25\cdot0{,}2=0{,}14$ a~všechny tři s~$0{,}015$, takže accuracy ensemblu je $0{,}845>0{,}80$; při \emph{disjunktních} chybových množinách až $100\,\%$. Korelované chyby kombinaci uškodí. (Přesnosti $0{,}70/0{,}75/0{,}80$ samy o~sobě výsledek \emph{neurčují} -- rozhoduje překryv chybových množin. Ověřeno vyčerpávajícím prohledáním rozložení chyb, \texttt{src/fig/gen\_s8\_ensemble.py}.)
\end{priklad}

\section{Statistické testy}\label{sec:testy}

Statistický test rozhoduje \emph{ano/ne} otázku o~datech: liší se průměrná hmotnost sušenek od deklarované? Liší se počty zákazníků mezi dny? Tato látka patří do \textbf{pravděpodobnosti a~statistiky} (okruh \textbf{M5}, předmět NMAI059), ne do strojového učení v~užším smyslu; požadavky S3 ji ale jmenují, protože se ve strojovém učení používá k~ověření, zda je rozdíl mezi dvěma modely (nebo mezi pozorováním a~předpokladem) \emph{statisticky významný}, nebo jen náhodný šum. Zde proto zopakujeme obecný rámec testování hypotéz (podrobně viz M5) a~rozebereme tři testy z~požadavků: \emph{jednovýběrový} a~\emph{dvouvýběrový Studentův t-test} a~\emph{chí-kvadrát test dobré shody}.

\begin{poznamka}[Rozsah]
Držíme se přesně požadavků: u~chí-kvadrát testu jen \emph{test dobré shody} (ne test nezávislosti v~kontingenční tabulce), u~t-testu jednovýběrovou a~dvouvýběrovou variantu. Resamplingové testy (bootstrap, permutační test) z~NPFL129 (přednáška~12) ani ANOVA do požadavků nepatří, neuvádíme je.
\end{poznamka}

\subsection{Rámec testování hypotéz}

Připomeňme z~M5 stavební kameny každého testu. Cíl je jediný: rozhodnout, zda data dostatečně silně \emph{odporují} výchozímu předpokladu.

\begin{definice}[Hypotézy, testová statistika, kritický obor]
\begin{itemize}
  \item \textbf{Nulová hypotéza} $H_0$: výchozí, konzervativní tvrzení („nic zajímavého se neděje`` -- sušenky váží přesně tolik, kolik mají; zákazníci přicházejí každý den stejně).
  \item \textbf{Alternativní hypotéza} $H_1$: tvrzení, které přijmeme, pokud $H_0$ zamítneme (hmotnost se liší; rozdělení neodpovídá).
  \item \textbf{Testová statistika} $T$: číslo spočtené z~dat, jehož rozdělení \emph{za platnosti $H_0$} známe (zde $t$-rozdělení nebo $\chi^2$-rozdělení).
  \item \textbf{Hladina významnosti} $\alpha$ (typicky $0{,}05$): maximální přípustná pravděpodobnost \emph{chyby 1.~druhu}, tj.\ zamítnutí $H_0$, ačkoli platí.
  \item \textbf{Kritická hodnota} $c$ a~\textbf{kritický obor} $W$: hranice a~množina hodnot $T$, pro které $H_0$ zamítáme. $W$ se volí tak, aby $P(T\in W\mid H_0)=\alpha$.
\end{itemize}
\end{definice}

\begin{intuice}
Postup je vždy stejný: (1)~formuluj $H_0$ a~$H_1$; (2)~zvol $\alpha$; (3)~spočti testovou statistiku $T$ z~dat; (4)~urči kritickou hodnotu z~příslušného rozdělení (z~tabulek nebo softwaru); (5)~je-li $T$ v~kritickém oboru (zde $T$ \uv{příliš velké}), \textbf{zamítni} $H_0$, jinak ji \emph{nezamítej}. Logika je důkazem sporem: předpokládáme $H_0$ a~ptáme se, jak nepravděpodobná jsou pak pozorovaná data. Pokud velmi (statistika padne do ocasu), $H_0$ zavrhneme.
\end{intuice}

\begin{poznamka}[p-hodnota]
Ekvivalentní pohled: \textbf{p-hodnota} je pravděpodobnost (za~$H_0$), že bychom dostali výsledek \emph{aspoň tak extrémní} jako pozorovaný. Rozhodovací pravidlo $T\in W$ je totéž co $p\le\alpha$. P-hodnota navíc říká, \emph{jak silně} data $H_0$ porušují. Pozor: p-hodnota \emph{není} pravděpodobnost, že $H_0$ platí. Rámec p-hodnoty viz M5.
\end{poznamka}

\subsection{Studentův t-test}

T-test porovnává \emph{střední hodnoty} u~normálně rozdělených dat, když \textbf{neznáme rozptyl} a~musíme ho odhadnout z~výběru. Právě nahrazení neznámého $\sigma$ výběrovou směrodatnou odchylkou $s$ vede místo normálního rozdělení na \emph{$t$-rozdělení} (Studentovo) s~$\nu$ \emph{stupni volnosti}: je symetrické kolem nuly, ale má těžší ocasy než $N(0,1)$; s~rostoucím $\nu$ se k~$N(0,1)$ blíží (pro $\nu\gtrsim 30$ jsou prakticky k~nerozeznání).

\subsubsection*{Jednovýběrový t-test}

\begin{definice}[Jednovýběrový t-test]
Máme výběr $X_1,\dots,X_n$ z~normálního rozdělení s~neznámou střední hodnotou $\mu$ a~neznámým rozptylem. Testujeme, zda se $\mu$ rovná dané konstantě $\mu_0$:
\[ H_0\colon \mu=\mu_0, \qquad H_1\colon \mu\neq\mu_0 \quad(\text{oboustranně}). \]
Testová statistika je
\[ t=\frac{\bar X-\mu_0}{s/\sqrt{n}}, \qquad \bar X=\frac1n\sum_{i=1}^n X_i,\quad s^2=\frac{1}{n-1}\sum_{i=1}^n (X_i-\bar X)^2, \]
kde $\bar X$ je výběrový průměr a~$s$ výběrová směrodatná odchylka. Za~$H_0$ má $t$ \textbf{Studentovo $t$-rozdělení s~$n-1$ stupni volnosti}. $H_0$ na hladině $\alpha$ zamítáme, pokud $|t|\ge t_{n-1,\,1-\alpha/2}$ (kritická hodnota oboustranného testu).
\end{definice}

\begin{intuice}
Čitatel $\bar X-\mu_0$ měří, jak daleko je pozorovaný průměr od testované hodnoty; jmenovatel $s/\sqrt{n}$ je \emph{směrodatná chyba průměru}, tedy typický rozptyl $\bar X$. Statistika $t$ tak vyjadřuje rozdíl \emph{v~jednotkách jeho vlastní nejistoty}: $t=10$ znamená, že pozorovaný průměr je deset směrodatných chyb od $\mu_0$ -- to je za~$H_0$ extrémně nepravděpodobné. Jednostrannou variantu ($H_1\colon \mu>\mu_0$, resp.\ $<$) použijeme, je-li předem jasný směr odchylky; pak je kritická hodnota $t_{n-1,\,1-\alpha}$ a~kritický obor jen jeden ocas.
\end{intuice}

\begin{center}
\begin{tikzpicture}
\begin{axis}[
    width=11cm, height=5.4cm,
    axis lines=middle,
    xlabel={$t$}, ylabel={hustota},
    xtick={-1.98,0,1.98}, xticklabels={$-1{,}98$,$0$,$1{,}98$},
    ytick=\empty,
    xmin=-4.6, xmax=4.6, ymin=0, ymax=0.46,
    enlargelimits=false, clip=false,
    every axis x label/.style={at={(current axis.right of origin)},anchor=west},
    every axis y label/.style={at={(current axis.above origin)},anchor=south west},
]
  % hustota t_99 (z .dat, scipy)
  \addplot[blue!70!black, thick] table[x=x,y=p]{src/fig/s9_t_df99.dat};
  % levy ocas (kriticky obor) -- vykreslime znovu jen cast a uzavreme k ose
  \addplot[draw=none, fill=red!55, fill opacity=0.32, restrict x to domain=-4.6:-1.98]
    table[x=x,y=p]{src/fig/s9_t_df99.dat} \closedcycle;
  % pravy ocas (kriticky obor)
  \addplot[draw=none, fill=red!55, fill opacity=0.32, restrict x to domain=1.98:4.6]
    table[x=x,y=p]{src/fig/s9_t_df99.dat} \closedcycle;
  % kriticke hodnoty (hranice kritickeho oboru) -- jen po krivku, ne az do vrcholu
  \addplot[black, dashed] coordinates {(-1.98,0) (-1.98,0.30)};
  \addplot[black, dashed] coordinates {(1.98,0) (1.98,0.30)};
  \node[red!60!black, anchor=north, font=\small] at (axis cs:-2.75,0.085) {$\alpha/2$};
  \node[red!60!black, anchor=north, font=\small] at (axis cs:2.75,0.085) {$\alpha/2$};
  \node[anchor=center, font=\small] at (axis cs:0,0.24) {nezamítáme};
\end{axis}
\end{tikzpicture}
\obrpopis{Hustota $t$-rozdělení s~$99$ stupni volnosti (oboustranný t-test, $n=100$). Při $\alpha=0{,}05$ leží kritický obor v~obou ocasech za~$\pm t_{99,\,0{,}975}\approx\pm1{,}98$ (\textcolor{red!70!black}{červeně}, dohromady míra $0{,}05$). Pro $99$ stupňů volnosti je rozdělení prakticky nerozeznatelné od $N(0,1)$ (kritická hodnota $1{,}98$ vs.\ $z_{0{,}975}=1{,}96$). Hustota generována scipy (\texttt{src/fig/gen\_s9\_dist.py}).}
\end{center}

\subsubsection*{Dvouvýběrový t-test}

\begin{definice}[Dvouvýběrový t-test]
Máme \emph{dva nezávislé} výběry $X_1,\dots,X_{n_1}$ a~$Y_1,\dots,Y_{n_2}$ z~normálních rozdělení se středními hodnotami $\mu_1,\mu_2$. Testujeme rovnost středních hodnot:
\[ H_0\colon \mu_1=\mu_2,\qquad H_1\colon \mu_1\neq\mu_2. \]
Při \emph{shodných rozptylech} obou skupin se testová statistika spočte ze \emph{sdruženého} odhadu rozptylu $s_p^2$:
\[ t=\frac{\bar X-\bar Y}{s_p\sqrt{\tfrac{1}{n_1}+\tfrac{1}{n_2}}},\qquad
   s_p^2=\frac{(n_1-1)s_1^2+(n_2-1)s_2^2}{n_1+n_2-2}, \]
kde $\bar X,\bar Y$ jsou výběrové průměry a~$s_1^2,s_2^2$ výběrové rozptyly. Za~$H_0$ má $t$ Studentovo rozdělení s~$n_1+n_2-2$ stupni volnosti; $H_0$ zamítáme při $|t|\ge t_{n_1+n_2-2,\,1-\alpha/2}$.
\end{definice}

\begin{intuice}
Myšlenka je stejná jako u~jednovýběrového testu, jen místo „průměr vs.\ konstanta`` porovnáváme „průměr vs.\ průměr``: čitatel je rozdíl průměrů $\bar X-\bar Y$, jmenovatel jeho směrodatná chyba. Sdružený rozptyl $s_p^2$ je vážený průměr rozptylů obou skupin (váhami jsou stupně volnosti) -- nejlepší společný odhad společného $\sigma^2$. Nelze-li předpokládat shodné rozptyly, používá se \emph{Welchova} varianta (jiný jmenovatel a~stupně volnosti); princip rozhodování zůstává.
\end{intuice}

\begin{poznamka}[Párový vs.\ nepárový]
Jsou-li data \emph{párovaná} (např.\ stejní pacienti před a~po léčbě), počítáme rozdíly $D_i=X_i-Y_i$ a~aplikujeme \emph{jednovýběrový} t-test na $H_0\colon \mu_D=0$. Výše uvedený dvouvýběrový test předpokládá \emph{nezávislé} skupiny.
\end{poznamka}

\begin{priklad}[Dvouvýběrový t-test -- ilustrace]
\textbf{Zadání.} Dvě skupiny studentů (každá $n=8$) psaly stejný test. Skupina A měla průměr $80$ a~výběrovou směrodatnou odchylku $4$, skupina B průměr $74$ a~odchylku $4$. Liší se průměry statisticky významně na hladině $\alpha=0{,}05$ (předpokládáme normalitu a~shodné rozptyly)?

\smallskip
\textbf{Řešení.} Sdružený rozptyl $s_p^2=\frac{7\cdot 16+7\cdot 16}{14}=16$, tedy $s_p=4$. Testová statistika
\[ t=\frac{80-74}{4\sqrt{\tfrac18+\tfrac18}}=\frac{6}{4\cdot 0{,}5}=\boxed{3{,}0},\qquad \text{df}=n_1+n_2-2=14. \]
Kritická hodnota $t_{14,\,0{,}975}\approx 2{,}14$. Protože $|3{,}0|\ge 2{,}14$, $H_0$ \textbf{zamítáme} -- skupiny se liší.
\end{priklad}

\begin{priklad}[SZZ jaro 2024 č.~26: Jednovýběrový t-test (sušenky)]
\szzotazka{jaro 2024}{26}{Statistické testy -- t-test}\par
\textbf{Zadání.} Výrobce deklaruje hmotnost sušenek $100\,$g. Z~kontroly: $n=100$ sušenek, průměrná hmotnost $\bar X=102\,$g, výběrová směrodatná odchylka $s=2\,$g. T-testem rozhodněte, zda se průměrná hmotnost liší od deklarovaných $100\,$g. (1)~Hypotézy. (2)~Testová statistika a~její rozdělení. (3)~Rozhodnutí, je-li kritická hodnota na $\alpha=0{,}05$ přibližně $2$.

\smallskip
\textbf{Řešení.}

\emph{(1)~Hypotézy.} Jednovýběrový t-test proti deklarované hodnotě $\mu_0=100$:
\[ H_0\colon \mu=100,\qquad H_1\colon \mu\neq 100 \quad(\text{oboustranný}). \]

\emph{(2)~Testová statistika.} Dosadíme do vzorce $t=(\bar X-\mu_0)/(s/\sqrt n)$:
\[ t=\frac{102-100}{2/\sqrt{100}}=\frac{2}{2/10}=\frac{2}{0{,}2}=\boxed{10}. \]
Za~$H_0$ má $t$-rozdělení s~$n-1=\boxed{99}$ stupni volnosti.

\emph{(3)~Rozhodnutí.} Kritická hodnota je $t_{99,\,0{,}975}\approx 2$ (přesně $1{,}984$). Protože $|t|=10\gg 2$, statistika leží hluboko v~kritickém oboru, \textbf{$H_0$ zamítáme}. Průměrná hmotnost sušenek se na hladině $\alpha=0{,}05$ \emph{statisticky významně liší} od deklarovaných $100\,$g (p-hodnota je prakticky nulová).
\end{priklad}

\subsection{Chí-kvadrát test dobré shody}

Zatímco t-test pracuje se \emph{spojitými} daty a~jejich průměrem, \textbf{chí-kvadrát test dobré shody} (goodness-of-fit) pracuje s~\emph{kategoriálními} daty: máme pozorované počty (četnosti) v~několika kategoriích a~ptáme se, zda odpovídají nějakému teoretickému rozdělení.

\begin{definice}[Chí-kvadrát test dobré shody]
Pozorování je rozděleno do $k$ kategorií s~\emph{pozorovanými četnostmi} $O_1,\dots,O_k$ (celkem $n=\sum_i O_i$). Za~$H_0$ jsou teoretické pravděpodobnosti kategorií $p_1,\dots,p_k$ ($\sum p_i=1$), což dává \emph{očekávané četnosti} $E_i=n\,p_i$:
\[ H_0\colon \text{pozorované četnosti odpovídají očekávaným } (E_i=n p_i), \qquad H_1\colon \text{neodpovídají}. \]
Testová statistika je
\[ \chi^2=\sum_{i=1}^k \frac{(O_i-E_i)^2}{E_i}. \]
Za~$H_0$ má pro dostatečně velká $n$ přibližně \textbf{$\chi^2$-rozdělení s~$k-1$ stupni volnosti}. $H_0$ zamítáme, pokud $\chi^2\ge \chi^2_{k-1,\,1-\alpha}$ (test je vždy \emph{jednostranný} -- velká hodnota svědčí proti shodě).
\end{definice}

\begin{intuice}
Každý sčítanec $(O_i-E_i)^2/E_i$ měří, jak moc se pozorovaná četnost v~kategorii odchýlila od očekávané, \emph{relativně} k~očekávané (dělení $E_i$ normuje měřítko). Když data sedí, jsou rozdíly malé a~$\chi^2$ je malé; velký nesoulad v~kterékoli kategorii statistiku nafoukne. Stupňů volnosti je $k-1$, ne $k$, protože jedna vazba $\sum_i O_i=n$ ubírá jeden stupeň volnosti (poslední četnost je dopočitatelná z~ostatních). Pravidlo: aproximace $\chi^2$-rozdělením je spolehlivá, jen když jsou očekávané četnosti dost velké (obvykle $E_i\ge 5$).
\end{intuice}

\begin{center}
\begin{tikzpicture}
\begin{axis}[
    width=11cm, height=5.4cm,
    axis lines=middle,
    xlabel={$\chi^2$}, ylabel={hustota},
    xtick={0,5,9.49,15,20}, xticklabels={$0$,$5$,$9{,}49$,$15$,$20$},
    ytick=\empty,
    xmin=0, xmax=20.5, ymin=0, ymax=0.20,
    enlargelimits=false, clip=false,
    every axis x label/.style={at={(current axis.right of origin)},anchor=west},
    every axis y label/.style={at={(current axis.above origin)},anchor=south west},
]
  % hustota chi^2_4 (z .dat, scipy)
  \addplot[blue!70!black, thick] table[x=x,y=p]{src/fig/s9_chi2_df4.dat};
  % kriticky obor: pravy ocas x >= 9.49
  \addplot[draw=none, fill=red!55, fill opacity=0.32, restrict x to domain=9.49:20.5]
    table[x=x,y=p]{src/fig/s9_chi2_df4.dat} \closedcycle;
  % kriticka hodnota (hranice kritickeho oboru)
  \addplot[black, dashed] coordinates {(9.49,0) (9.49,0.105)};
  % pozorovana statistika (pozorovana hodnota chi^2)
  \addplot[black!55!red, thick] coordinates {(10.1,0) (10.1,0.155)};
  % popisek kriticke hodnoty -- vlevo dole, jasne oddeleny
  \node[anchor=south east, font=\small] at (axis cs:9.35,0.108) {$c=9{,}49$};
  % popisek pozorovane statistiky -- vyse a vpravo, bez kolize s c
  \node[black!55!red, anchor=south west, font=\small] at (axis cs:10.25,0.155) {$\chi^2=10{,}1$};
  \node[red!60!black, anchor=west, font=\small] at (axis cs:12.0,0.05) {kritický obor ($\alpha=0{,}05$)};
\end{axis}
\end{tikzpicture}
\obrpopis{Hustota $\chi^2$-rozdělení se~$4$ stupni volnosti (pět kategorií, $k-1=4$). Kritický obor na hladině $\alpha=0{,}05$ je pravý ocas za~kritickou hodnotou $c=\chi^2_{4,\,0{,}95}\approx 9{,}49$ (\textcolor{red!70!black}{červeně}, míra $0{,}05$). Pozorovaná statistika $\chi^2=10{,}1$ leží v~kritickém oboru, takže $H_0$ \emph{zamítáme}. Rozdělení je nesymetrické a~kladné (statistika je součet čtverců). Hustota generována scipy (\texttt{src/fig/gen\_s9\_dist.py}).}
\end{center}

\begin{priklad}[SZZ podzim 2023 č.~29: Chí-kvadrát test dobré shody (zákazníci)]
\szzotazka{podzim 2023}{29}{chí-kvadrát test (test dobré shody)}\par
\textbf{Zadání.} (1)~Popište $\chi^2$ test dobré shody: hypotézy, testová statistika, její rozdělení. (2)~Obchod měl za pět pracovních dní celkem $100$ zákazníků; po dnech (po--pá) postupně $25$, $30$, $15$, $14$, $16$. Rozhodněte $\chi^2$ testem, zda se počty mezi dny statisticky významně liší. \textit{Nápověda:} kritické hodnoty na $\alpha=0{,}05$ jsou $7{,}8$ (df${}=3$), $9{,}5$ (df${}=4$), $11{,}1$ (df${}=5$), $12{,}6$ (df${}=6$).

\smallskip
\textbf{Řešení.}

\emph{(1)~Popis testu.} $H_0$: pozorované četnosti odpovídají očekávaným; $H_1$: neodpovídají. Statistika $\chi^2=\sum_i (O_i-E_i)^2/E_i$ má za~$H_0$ rozdělení $\chi^2$ s~$k-1$ stupni volnosti, kde $k$ je počet kategorií. Velká hodnota svědčí proti $H_0$.

\emph{(2)~Výpočet.} \uv{Zákazníků je každý den stejně}: $5$ kategorií, $n=100$, takže za~$H_0$ je $p_i=\tfrac15$ a~očekávaná četnost $E_i=n p_i=100\cdot\tfrac15=20$ pro každý den. Dosadíme pozorované četnosti:
\[ \chi^2=\frac{(25-20)^2}{20}+\frac{(30-20)^2}{20}+\frac{(15-20)^2}{20}+\frac{(14-20)^2}{20}+\frac{(16-20)^2}{20}. \]
Čitatelé jsou $5^2,10^2,(-5)^2,(-6)^2,(-4)^2=25,100,25,36,16$, součet $202$, tedy
\[ \chi^2=\frac{25+100+25+36+16}{20}=\frac{202}{20}=\boxed{10{,}1}. \]
Stupně volnosti: $k-1=5-1=4$, kritická hodnota z~nápovědy $c=\chi^2_{4,\,0{,}95}=9{,}5$.

\emph{Rozhodnutí.} $\chi^2=10{,}1>9{,}5=c$, statistika padla do kritického oboru, \textbf{$H_0$ zamítáme}. Počty zákazníků v~jednotlivých dnech se na hladině $\alpha=0{,}05$ \emph{statisticky významně liší} (rozdíl není jen náhodný šum; p-hodnota $\approx 0{,}039<0{,}05$). Největší příspěvek ke statistice dává úterý ($30$ zákazníků, příspěvek $100/20=5$ z~celkových $10{,}1$), ostatní dny přispívají nejvýše $1{,}8$ (čtvrtek).
\end{priklad}

\begin{poznamka}[Když test nezamítá]
Pozor na interpretaci: \emph{nezamítnutí} $H_0$ neznamená, že $H_0$ \uv{platí} -- jen že data nemají dost síly ji vyvrátit. U~obou příkladů by stačil výrazně menší rozdíl (sušenky $100{,}3\,$g, vyrovnanější počty zákazníků) k~tomu, aby test nezamítl, ač $H_0$ doslova neplatí. Statistická významnost také není totéž co \emph{praktická} významnost: u~obrovského vzorku zamítneme $H_0$ i~při zanedbatelném rozdílu.
\end{poznamka}

\section{Učení bez učitele: shlukování a redukce dimenze}\label{sec:bezucitele}

Dosud měla každá data \emph{cíl} $t$ (label) -- učili jsme se s~učitelem. Při \textbf{učení bez učitele} máme jen vstupy $\mathbf{x}_1,\dots,\mathbf{x}_N$ \emph{bez} cílů a~hledáme v~nich \emph{strukturu}: přirozené skupiny (\emph{shlukování}) nebo nízkodimenzionální popis (\emph{redukce dimenze}). Není „správná odpověď`` k~porovnání, takže evaluace je vnitřní (jak kompaktní jsou shluky, kolik rozptylu zachováme).

\begin{poznamka}[Vztah k~S1 a~M2]
S1 (sekce \emph{Strojové učení}) zmiňuje EM jako Bayesovské učení se skrytými proměnnými a~poznamenává, že $k$-means je jeho speciální případ s~„tvrdým`` přiřazením. Zde shlukování rozvádíme do hloubky (algoritmus, kritérium, konvergence) a~EM dořešíme na konkrétní SZZ úloze. \emph{Redukce dimenze} přes PCA stojí na vlastních číslech a~vlastních vektorech symetrické matice (okruh \textbf{M2}); odtud bereme spektrální rozklad jako hotový nástroj.
\end{poznamka}

% =====================================================================
\subsection{Shlukování algoritmem \texorpdfstring{$k$}{k}-means}\label{ssec:kmeans}

\begin{definice}[Úloha shlukování a~kritérium $k$-means]
Mějme $N$ bodů $\mathbf{x}_1,\dots,\mathbf{x}_N\in\mathbb{R}^D$ a~zadaný počet shluků $K$. Hledáme \emph{centroidy} (středy) $\mu_1,\dots,\mu_K\in\mathbb{R}^D$ a~přiřazení každého bodu k~jednomu shluku, popsané indikátory $z_{i,k}\in\{0,1\}$ ($z_{i,k}=1$, právě když bod $\mathbf{x}_i$ patří do shluku $k$; pro každé $i$ je právě jedno $z_{i,k}=1$). Minimalizujeme \textbf{kritérium} (součet čtverců vzdáleností bodů od jejich centroidu)
\[
  J(\{z_{i,k}\},\{\mu_k\}) \;=\; \sum_{i=1}^N\sum_{k=1}^K z_{i,k}\,\|\mathbf{x}_i-\mu_k\|^2 .
\]
\end{definice}

\begin{intuice}
$J$ je celková „rozptýlenost`` shluků: čím menší, tím kompaktnější jsou shluky kolem svých středů. Najít globální minimum přes \emph{všechna} přiřazení je NP-těžké (exponenciálně mnoho rozdělení), proto $k$-means minimalizuje $J$ \emph{střídavě}: drží středy a~optimalizuje přiřazení, pak drží přiřazení a~optimalizuje středy. To je přesně struktura EM (sekce~\ref{ssec:em}).
\end{intuice}

\begin{algo}[$k$-means (Lloydův algoritmus)]
\textbf{Vstup:} body $\mathbf{x}_1,\dots,\mathbf{x}_N$, počet shluků $K$.\\[2pt]
\textbf{1. Inicializace.} Zvol $K$ počátečních centroidů $\mu_1,\dots,\mu_K$ \cmt{náhodně z~dat / k-means++}\\[2pt]
\textbf{2.} \textbf{opakuj} až do konvergence (přiřazení se nemění) nebo do limitu iterací:\\
\hspace*{1.5em}\textbf{(a) Krok přiřazení.} Každý bod přiřaď k~\emph{nejbližšímu} centroidu:\hfill\cmt{drží středy, mění $z$}
\[
  z_{i,k}=\begin{cases}1 & k=\argmin_{j}\|\mathbf{x}_i-\mu_j\|^2,\\[1pt]0 & \text{jinak.}\end{cases}
\]
\hspace*{1.5em}\textbf{(b) Krok přepočtu.} Každý centroid přepočti jako \emph{průměr} bodů svého shluku:\hfill\cmt{drží $z$, mění středy}
\[
  \mu_k=\frac{\sum_i z_{i,k}\,\mathbf{x}_i}{\sum_i z_{i,k}} .
\]
\textbf{Výstup:} přiřazení $z_{i,k}$ a~centroidy $\mu_k$.
\end{algo}

\begin{intuice}
Oba kroky jsou \emph{optimální pro svoji polovinu}. (a)~Při pevných středech je $J$ minimalizováno tím, že každý bod jde k~nejbližšímu středu (jiné přiřazení by jen zvětšilo jeho člen $\|\mathbf{x}_i-\mu_k\|^2$). (b)~Při pevném přiřazení je $\sum_i z_{i,k}\|\mathbf{x}_i-\mu\|^2$ jako funkce $\mu$ minimalizována \emph{průměrem} bodů shluku -- derivace podle $\mu$ je $-2\sum_i z_{i,k}(\mathbf{x}_i-\mu)=0$, odkud $\mu=\bar{\mathbf{x}}$ shluku. Proto se každý krok počítá tak, jak se počítá.
\end{intuice}

\subsubsection*{Inicializace: náhodně vs.\ k-means++}

Náhodná inicializace volí $K$ počátečních středů jako náhodné body z~dat. Riziko: dva středy padnou do téhož hustého shluku a~jeden přirozený shluk zůstane bez středu. \textbf{k-means++} volí středy „rozprostřeně``: první náhodně, každý další s~pravděpodobností \emph{úměrnou čtverci vzdálenosti} od nejbližšího už zvoleného středu. Vzdálené (dosud nepokryté) oblasti tak mají větší šanci dostat střed. Dokazatelně dává v~očekávání $\mathcal{O}(\log K)$-aproximaci optimálního $J$.

\begin{poznamka}[Praxe]
Protože $k$-means konverguje jen k~\emph{lokálnímu} minimu (viz níže), spouští se obvykle \emph{vícekrát} s~různými inicializacemi a~vybere se běh s~\emph{nejmenším} $J$ (scikit-learn má ve výchozím stavu \texttt{n\_init=10}).
\end{poznamka}

\subsubsection*{Konvergence k~lokálnímu optimu}

\begin{veta}[$k$-means konverguje (k~lokálnímu minimu)]
Algoritmus $k$-means se po konečně mnoha iteracích zastaví; kritérium $J$ při tom neroste a~výsledek je \emph{lokální} (ne nutně globální) minimum.
\end{veta}

\begin{intuice}[Argument konvergence -- dotažený]
Tvrzení stojí na dvou pozorováních a~jedné konečnosti:
\begin{enumerate}
  \item \textbf{Žádný krok $J$ nezvýší.} Krok přiřazení (a) přiřadí každý bod k~nejbližšímu středu, takže jeho příspěvek $\|\mathbf{x}_i-\mu_k\|^2$ klesne nebo zůstane -- $J$ neroste. Krok přepočtu (b) nahradí každý střed průměrem svého shluku, který \emph{minimalizuje} $\sum_i z_{i,k}\|\mathbf{x}_i-\mu\|^2$ (viz intuice výše) -- $J$ opět neroste. Posloupnost hodnot $J$ je tedy \emph{nerostoucí} a~zdola omezená nulou.
  \item \textbf{Konečně mnoho stavů.} $J$ je plně určeno \emph{přiřazením} $\{z_{i,k}\}$ (středy jsou pak dané průměry). Různých přiřazení $N$ bodů do $K$ shluků je konečně mnoho ($\le K^N$).
  \item \textbf{Závěr.} Algoritmus prochází přiřazeními tak, že $J$ neroste; kdyby se přiřazení změnilo, $J$ \emph{ostře} klesne (jinak jsme v~minimu a~iterace se zastaví). Ostře klesající hodnota na konečné množině stavů nemůže pokračovat donekonečna -- žádný stav se nezopakuje. Po konečně mnoha krocích se přiřazení přestane měnit, tehdy se zastavíme.
\end{enumerate}
Že jde jen o~\emph{lokální} minimum, plyne z~toho, že střídavá minimalizace optimalizuje vždy jen jednu skupinu proměnných; výsledek je „nejlepší v~okolí``, ne globálně. \emph{Proto} je výsledek citlivý na inicializaci a~spouští se vícekrát (viz výše).
\end{intuice}

\begin{poznamka}[Pozor na prázdnou frázi]
Říct „$k$-means konverguje, protože $J$ klesá`` je \emph{neúplné}: klesání samo ještě nezaručuje \emph{zastavení} (klesat lze i~nekonečně, $1/n\to 0$). Klíč je \emph{konečnost} počtu přiřazení -- proto se posloupnost monotónně klesajících hodnot na konečné množině stavů musí ustálit.
\end{poznamka}

\begin{center}
\begin{tikzpicture}[scale=0.84]
% --- Panel 1: krok prirazeni ---
\begin{scope}
  \draw[osa] (0,0) -- (4.2,0) node[right,font=\scriptsize]{$x_1$};
  \draw[osa] (0,0) -- (0,3.9) node[above,font=\scriptsize]{$x_2$};
  % body shluku 1 (vlevo dole), 2 (vpravo nahore)
  \node[cls1] at (0.8,0.7) {}; \node[cls1] at (1.3,1.1) {}; \node[cls1] at (0.6,1.4) {};
  \node[cls1] at (1.1,0.5) {};
  \node[cls2] at (3.0,3.0) {}; \node[cls2] at (3.4,2.5) {}; \node[cls2] at (2.6,3.3) {};
  \node[cls2] at (3.2,3.5) {};
  % stare centroidy (jeste neprepoctene -- mimo teziste shluku)
  \node[centroid] (m1) at (1.7,1.6) {};
  \node[centroid] (m2) at (2.4,2.5) {};
  \node[font=\scriptsize,anchor=south] at (1.7,1.80) {$\mu_1$};
  \node[font=\scriptsize,anchor=south] at (2.4,2.70) {$\mu_2$};
  \node[font=\scriptsize] at (2.1,3.78) {\textbf{(a) krok přiřazení}};
\end{scope}
% --- Panel 2: krok prepoctu ---
\begin{scope}[xshift=5.6cm]
  \draw[osa] (0,0) -- (4.2,0) node[right,font=\scriptsize]{$x_1$};
  \draw[osa] (0,0) -- (0,3.9) node[above,font=\scriptsize]{$x_2$};
  \node[cls1] at (0.8,0.7) {}; \node[cls1] at (1.3,1.1) {}; \node[cls1] at (0.6,1.4) {};
  \node[cls1] at (1.1,0.5) {};
  \node[cls2] at (3.0,3.0) {}; \node[cls2] at (3.4,2.5) {}; \node[cls2] at (2.6,3.3) {};
  \node[cls2] at (3.2,3.5) {};
  % stare centroidy (svetle) a sipka na nove (prumer = teziste shluku)
  \node[star,star points=5,draw=black!35,fill=yellow!30,inner sep=0pt,minimum size=9pt] (o1) at (1.7,1.6) {};
  \node[star,star points=5,draw=black!35,fill=yellow!30,inner sep=0pt,minimum size=9pt] (o2) at (2.4,2.5) {};
  \draw[->,>=Stealth,black!65,thick] (o1) -- (0.98,0.95);
  \draw[->,>=Stealth,black!65,thick] (o2) -- (3.05,3.075);
  \node[centroid] (n1) at (0.95,0.925) {};
  \node[centroid] (n2) at (3.05,3.075) {};
  \node[font=\scriptsize,anchor=west] at (1.18,0.88) {$\mu_1'$};
  \node[font=\scriptsize,anchor=west] at (3.22,3.02) {$\mu_2'$};
  \node[font=\scriptsize] at (2.1,3.78) {\textbf{(b) krok přepočtu}};
\end{scope}
\end{tikzpicture}
\obrpopis{Jedna iterace $k$-means. (a)~Krok přiřazení: každý bod (modrý kruh / červený čtverec) jde k~bližšímu z~centroidů (žluté hvězdy $\mu_1,\mu_2$). (b)~Krok přepočtu: každý centroid se posune do těžiště (průměru) bodů svého shluku -- šipka vede od staré polohy k~nové $\mu_k'$. Kroky se střídají, dokud se přiřazení nemění.}
\end{center}

\subsubsection*{Řešené příklady: $k$-means}

\begin{priklad}[SZZ léto 2022 č.~12: K-means -- první iterace]
\szzotazka{léto 2022}{12}{Shlukování (K-means)}\par
\textbf{Zadání.} Šest bodů s~příznaky $A_1,A_2$: $1(2,4)$, $2(1,3)$, $3(0,4)$, $4(5,2)$, $5(6,2)$, $6(4,0)$; $K=2$, počáteční náhodné přiřazení $C_1=\{2,3,4\}$, $C_2=\{1,5,6\}$. (1)~Spočítejte centroidy obou shluků. (2)~Přiřaďte každý bod k~nejbližšímu centroidu (euklidovsky) a~uveďte nové shluky.

\smallskip
\textbf{Řešení.} (1)~Centroid je průměr bodů shluku:
\[ c_1=\Big(\tfrac{1+0+5}{3},\tfrac{3+4+2}{3}\Big)=(2,3),\qquad c_2=\Big(\tfrac{2+6+4}{3},\tfrac{4+2+0}{3}\Big)=(4,2). \]
(2)~Pro každý bod porovnáme čtverce vzdáleností k~$c_1=(2,3)$ a~$c_2=(4,2)$:
\[
\begin{array}{c|cc|c}
\text{bod} & \|{\cdot}-c_1\|^2 & \|{\cdot}-c_2\|^2 & \text{shluk}\\\hline
1\,(2,4) & 1  & 8  & C_1\\
2\,(1,3) & 1  & 10 & C_1\\
3\,(0,4) & 5  & 20 & C_1\\
4\,(5,2) & 10 & 1  & C_2\\
5\,(6,2) & 17 & 4  & C_2\\
6\,(4,0) & 13 & 4  & C_2
\end{array}
\]
Nové shluky: $C_1=\{1,2,3\}$, $C_2=\{4,5,6\}$.
\end{priklad}

\begin{priklad}[SZZ léto 2023 č.~8: $k$-means iterace + hierarchické shlukování]
\szzotazka{léto 2023}{8}{Shlukování}\par
\textbf{Zadání.} Data $(1,1),(2,1),(3,2),(4,3),(4,2)$, $k=2$, počáteční středy $\mu_1=(1,1)$, $\mu_2=(3,2)$. (1)~Nové středy po jedné iteraci? (2)~Co v~dalších iteracích? (3)~Popište hierarchické shlukování a~jeho výsledek na těchto datech.

\smallskip
\textbf{Řešení.} (1) \emph{Krok přiřazení} -- ke každému bodu spočteme vzdálenost k~oběma středům (stačí porovnat čtverce):
\[
\begin{array}{c|cc|c}
\text{bod} & \|{\cdot}-\mu_1\| & \|{\cdot}-\mu_2\| & \text{shluk}\\\hline
(1,1) & 0 & \sqrt5 & 1\\
(2,1) & 1 & \sqrt2 & 1\\
(3,2) & \sqrt5 & 0 & 2\\
(4,3) & \sqrt{13} & \sqrt2 & 2\\
(4,2) & \sqrt{10} & 1 & 2
\end{array}
\]
Shluk~1 $=\{(1,1),(2,1)\}$, shluk~2 $=\{(3,2),(4,3),(4,2)\}$. \emph{Krok přepočtu} (průměry):
\[
  \mu_1=\Big(\tfrac{1+2}{2},\tfrac{1+1}{2}\Big)=\boxed{(1{,}5,\;1)},\qquad
  \mu_2=\Big(\tfrac{3+4+4}{3},\tfrac{2+3+2}{3}\Big)=\boxed{\big(\tfrac{11}{3},\,\tfrac{7}{3}\big)}.
\]
(2)~Opakované přiřazení s~novými středy dá \emph{totéž} rozdělení (body $(1,1),(2,1)$ zůstanou u~$\mu_1$, zbytek u~$\mu_2$), takže se středy už \textbf{nemění -- algoritmus zkonvergoval} (po jedné iteraci).

(3)~\emph{Aglomerativní} shlukování (viz~\ref{ssec:hier}): začni s~pěti jednobodovými shluky a~opakovaně slučuj dva nejbližší. Při \emph{single linkage} (vzdálenost nejbližších bodů) mají tři dvojice vzdálenost $1$: $(1,1)$--$(2,1)$, $(3,2)$--$(4,2)$ a~$(4,2)$--$(4,3)$. Ve výšce $1$ tak vzniknou skupiny $\{(1,1),(2,1)\}$ a~$\{(3,2),(4,2),(4,3)\}$ (při postupném slučování po jedné dvojici v~libovolném pořadí), které se nakonec spojí ve výšce $\sqrt2$ (vzdálenost $(2,1)$--$(3,2)$). Řez dendrogramu na dvě skupiny dá \emph{stejné dva shluky jako $k$-means}: $\{(1,1),(2,1)\}$ a~$\{(3,2),(4,2),(4,3)\}$.
\begin{center}
\begin{tikzpicture}[scale=0.85]
  \draw[osa] (0,0) -- (5.0,0) node[right,font=\scriptsize]{$x_1$};
  \draw[osa] (0,0) -- (0,3.8) node[above,font=\scriptsize]{$x_2$};
  \foreach \t in {1,2,3,4}{\draw (\t,0.05)--(\t,-0.05) node[below,font=\tiny]{\t};}
  \foreach \t in {1,2,3}{\draw (0.05,\t)--(-0.05,\t) node[left,font=\tiny]{\t};}
  % shluk 1
  \node[cls1] at (1,1) {}; \node[cls1] at (2,1) {};
  % shluk 2
  \node[cls2] at (3,2) {}; \node[cls2] at (4,3) {}; \node[cls2] at (4,2) {};
  % centroidy po konvergenci
  \node[centroid] at (1.5,1) {}; \node[centroid] at (3.6667,2.3333) {};
  \node[font=\tiny,anchor=south] at (1.5,1.15) {$\mu_1$};
  \node[font=\tiny,anchor=west] at (3.78,2.33) {$\mu_2$};
\end{tikzpicture}
\obrpopis{Data úlohy léto~2023. $k$-means s~$k=2$ konverguje na shluky $\{(1,1),(2,1)\}$ a~$\{(3,2),(4,2),(4,3)\}$ s~centroidy $(1{,}5,1)$ a~$(\tfrac{11}{3},\tfrac73)$ (žluté hvězdy). Single-linkage aglomerativní shlukování dává řezem na dvě skupiny totéž rozdělení.}
\end{center}
\end{priklad}

\begin{priklad}[SZZ jaro 2024 č.~28: $k$-means se třemi shluky]
\szzotazka{jaro 2024}{28}{Shlukování}\par
\textbf{Zadání.} Body $A(-2,1)$, $B(-1,-1)$, $C(3,-2)$, $D(4,-1)$, $E(3,3)$, $F(4,4)$, $G(-3,0)$, $H(4,2)$, $I(4,-2)$, $J(6,1)$; $k=3$, počáteční středy $A$, $F$, $J$. (1)~Hlavní kroky $k$-means. (2)~Rozdělení a~nové středy po první iteraci (euklidovská vzdálenost). (3)~Změní se středy dál?

\smallskip
\textbf{Řešení.} (1)~Inicializuj $K$ středů; pak opakuj \emph{přiřazení} (každý bod k~nejbližšímu středu) a~\emph{přepočet} (střed $=$ průměr bodů shluku), dokud se přiřazení mění.

(2)~\emph{Krok přiřazení} ke středům $A(-2,1)$, $F(4,4)$, $J(6,1)$. Vzdálenosti levých bodů jsou nejmenší k~$A$, horních pravých k~$F$, dolních/krajně pravých k~$J$. Např.\ $C(3,-2)$: k~$A$ je $\sqrt{25+9}=\sqrt{34}$, k~$F$ je $\sqrt{1+36}=\sqrt{37}$, k~$J$ je $\sqrt{9+9}=\sqrt{18}$ -- nejblíž $J$; $H(4,2)$: k~$F$ je $\sqrt{0+4}=2$, k~$J$ je $\sqrt{4+1}=\sqrt5$ -- nejblíž $F$. Výsledek:
\[
  \text{shluk }A:\{A,B,G\},\quad \text{shluk }F:\{E,F,H\},\quad \text{shluk }J:\{C,D,I,J\}.
\]
\emph{Krok přepočtu} (průměry):
\[
  \mu_A=\Big(\tfrac{-2-1-3}{3},\tfrac{1-1+0}{3}\Big)=\boxed{(-2,\,0)},\quad
  \mu_F=\Big(\tfrac{3+4+4}{3},\tfrac{3+4+2}{3}\Big)=\boxed{\big(\tfrac{11}{3},\,3\big)},
\]
\[
  \mu_J=\Big(\tfrac{3+4+4+6}{4},\tfrac{-2-1-2+1}{4}\Big)=\boxed{\big(\tfrac{17}{4},\,-1\big)}.
\]
(3)~Opakované přiřazení k~novým středům dá \emph{totéž} rozdělení, takže se středy už \textbf{nemění} -- algoritmus zkonvergoval.
\end{priklad}

\begin{priklad}[SZZ léto 2025 č.~31: kritérium a~lokální optimum]
\szzotazka{léto 2025}{31}{Shlukování ($k$-means)}\par
\textbf{Zadání.} (1)~Popište $k$-means (volba a~úprava středů, ukončení). (2)~Jaké kritérium optimalizuje? (3)~Na datech se třemi zjevnými shluky $k$-means ($k=3$) nenašel optimální shluky -- proč a~jak zvýšit šanci na lepší výsledek?

\smallskip
\textbf{Řešení.} (1)~Středy se volí (náhodně z~dat / k-means++); pak se střídá přiřazení bodů k~nejbližšímu středu a~přepočet středu na průměr shluku; končí se, když se přiřazení přestane měnit (nebo po daném počtu iterací). (2)~Minimalizuje $J=\sum_{i}\sum_{k}z_{i,k}\|\mathbf{x}_i-\mu_k\|^2$ (součet čtverců vzdáleností bodů od jejich centroidu).

(3)~$k$-means konverguje jen k~\emph{lokálnímu} minimu $J$, závislému na inicializaci. Při nešťastné inicializaci se některý přirozený shluk rozdělí mezi dva středy, zatímco jiný přirozený shluk vlastní střed nedostane. Na obrázku ze zadání si velký levý dolní shluk rozdělily značky „čtvereček`` a~„kolečko``, přičemž kolečka zároveň pokrývá malý horní prostřední shluk; oficiální náčrt to popisuje tak, že prostřední shluk (\emph{mnohem menší} než zbylé dva) byl přiřazen k~části levého. Takové rozdělení je lokálně stabilní (žádný bod nechce přejít k~jinému středu), ale má vyšší $J$ než pravé tři shluky. \textbf{Náprava:} spustit $k$-means \emph{vícekrát} z~různých inicializací a~vzít běh s~nejmenším $J$; použít \emph{k-means++} inicializaci, která rozprostře počáteční středy.
\end{priklad}

\begin{priklad}[SZZ jaro 2026 č.~28: $k$-means a~argument konvergence]
\szzotazka{jaro 2026}{28}{Shlukování pomocí $K$-means}\par
\textbf{Zadání.} (1)~Popište $k$-means, k~čemu slouží, jeho kroky a~inicializaci. (2)~Jakou ztrátovou funkci optimalizuje? (3)~Konverguje? Zformulujte argument a~uveďte praktické důsledky.

\smallskip
\textbf{Řešení.} (1)~$k$-means je metoda \emph{učení bez učitele} pro rozdělení dat do $K$ shluků podle geometrické podobnosti. Inicializuje $K$ centroidů (náhodně z~dat nebo k-means++: první náhodně, další úměrně čtverci vzdálenosti od nejbližšího už zvoleného) a~opakuje přiřazení (bod $\to$ nejbližší centroid) a~přepočet (centroid $=$ průměr shluku), dokud se přiřazení mění. (2)~$J=\sum_i\sum_k z_{i,k}\|\mathbf{x}_i-\mu_k\|^2$.

(3)~\textbf{Ano, konverguje k~lokálnímu minimu.} Argument: oba kroky $J$ \emph{nezvýší} (přiřazení k~nejbližšímu středu i~přepočet na průměr jsou každý optimální pro svoji proměnnou); $J$ je zdola omezené nulou; a~hodnota $J$ je určena \emph{přiřazením}, kterých je jen \emph{konečně mnoho} ($\le K^N$). Nerostoucí posloupnost, která při každé skutečné změně přiřazení ostře klesá, se na konečné množině stavů musí po konečně mnoha krocích ustálit -- tehdy se přiřazení nemění a~zastavíme. \textbf{Praktické důsledky:} výsledek je jen \emph{lokální} optimum a~závisí na inicializaci, proto se algoritmus spouští vícekrát a~bere se nejlepší (nejmenší $J$); inicializace k-means++ kvalitu zlepšuje. (Konvergence je obvykle rychlá -- pár iterací.)
\end{priklad}

% =====================================================================
\subsection{Hierarchické shlukování}\label{ssec:hier}

$k$-means vyžaduje předem zvolené $K$ a~dává „placaté`` rozdělení. \textbf{Aglomerativní (zdola-nahoru) hierarchické shlukování} žádné $K$ předem nepotřebuje: vybuduje celou \emph{hierarchii} vnořených shluků a~rozhodnutí o~počtu skupin nechá až na konec (řez dendrogramu).

\begin{algo}[Aglomerativní hierarchické shlukování]
\textbf{Vstup:} body $\mathbf{x}_1,\dots,\mathbf{x}_N$, míra vzdálenosti shluků (\emph{linkage}).\\[2pt]
\textbf{1.} Začni s~$N$ shluky, každý bod sám jeden shluk. \cmt{singletony}\\
\textbf{2.} \textbf{opakuj}, dokud nezbude jeden shluk:\\
\hspace*{1.5em}(a) najdi dvojici shluků s~\emph{nejmenší} vzdáleností (dle zvoleného linkage);\\
\hspace*{1.5em}(b) sluč je do jednoho a~zaznamenej výšku spojení $=$ jejich vzdálenost.\\
\textbf{Výstup:} \emph{dendrogram} -- strom slučování s~výškami.
\end{algo}

\begin{definice}[Linkage: vzdálenost dvou shluků $\mathcal{A},\mathcal{B}$]
\[
\begin{aligned}
&\textbf{single}\ (\text{nejbližší soused}): && d(\mathcal A,\mathcal B)=\min_{a\in\mathcal A,b\in\mathcal B}\|a-b\|;\\
&\textbf{complete}\ (\text{nejvzdálenější}): && d(\mathcal A,\mathcal B)=\max_{a\in\mathcal A,b\in\mathcal B}\|a-b\|;\\
&\textbf{average}\ (\text{průměrná}): && d(\mathcal A,\mathcal B)=\tfrac{1}{|\mathcal A||\mathcal B|}\sum_{a,b}\|a-b\|;\\
&\textbf{Ward}: && \text{slučuj dvojici, která nejméně zvýší vnitroshlukový rozptyl.}
\end{aligned}
\]
\end{definice}

\begin{intuice}
Volba linkage mění \emph{tvar} shluků: \emph{single} umí natáhnout shluk podél řetězce blízkých bodů (riziko „řetězení`` -- dva oddělené shluky spojené můstkem splynou), \emph{complete} preferuje kompaktní kulaté shluky, \emph{average} je kompromis. \emph{Ward} minimalizuje růst rozptylu -- dává podobně kompaktní shluky jako $k$-means a~je výchozí volbou např.\ ve \texttt{scikit-learn}. \textbf{Dendrogram} je strom, jehož listy jsou body a~výška spojení je vzdálenost slučovaných shluků. \emph{Vodorovný řez} v~určité výšce rozpadne strom na shluky: čím níž řežeme, tím \emph{víc} shluků. Počet shluků tak volíme až dodatečně, čtením z~dendrogramu.
\end{intuice}

\begin{center}
\begin{tikzpicture}[scale=1.0]
% Dendrogram 5 bodu: a,b blizko; d,e blizko; c pripoji k {d,e}; pak obe skupiny.
% listy na ose x: a=0, b=1, c=2, d=3, e=4 (jednotka = 1cm). vyska = vzdalenost spojeni.
  \def\hab{0.8}   % spojeni a-b
  \def\hde{1.0}   % spojeni d-e
  \def\hcde{1.8}  % c k {d,e}
  \def\htop{3.0}  % obe skupiny
  % osa vysky
  \draw[osa] (-0.6,0) -- (-0.6,3.4) node[above,font=\scriptsize]{výška (vzdálenost)};
  \foreach \y in {1,2,3}{\draw (-0.65,\y)--(-0.55,\y) node[left,font=\tiny,xshift={0-2pt}]{\y};}
  % listy
  \foreach \x/\lab in {0/a,1/b,2/c,3/d,4/e}{\node[font=\scriptsize] at (\x,-0.25){$\lab$}; \draw[black!30,dotted] (\x,0)--(\x,3.0);}
  % spojeni a-b
  \draw[thick] (0,0)--(0,\hab); \draw[thick] (1,0)--(1,\hab);
  \draw[thick] (0,\hab)--(1,\hab);
  % spojeni d-e
  \draw[thick] (3,0)--(3,\hde); \draw[thick] (4,0)--(4,\hde); \draw[thick] (3,\hde)--(4,\hde);
  % c k {d,e}
  \draw[thick] (2,0)--(2,\hcde); \draw[thick] (3.5,\hde)--(3.5,\hcde); \draw[thick] (2,\hcde)--(3.5,\hcde);
  % top
  \draw[thick] (0.5,\hab)--(0.5,\htop); \draw[thick] (2.75,\hcde)--(2.75,\htop); \draw[thick] (0.5,\htop)--(2.75,\htop);
  % rez
  \draw[red!70!black,dashed,thick] (-0.4,2.4)--(4.4,2.4) node[right,font=\tiny,red!70!black]{řez ($2$ shluky)};
\end{tikzpicture}
\obrpopis{Dendrogram aglomerativního shlukování pěti bodů. Nejdřív se spojí blízké páry ($a,b$ a~$d,e$), pak se $c$ přidá ke~skupině $\{d,e\}$, nakonec splynou obě skupiny. Výška vodorovné spojnice $=$ vzdálenost slučovaných shluků. Vodorovný \emph{řez} (červeně) protne dvě svislé čáry, tedy dělí data na \textbf{dva} shluky $\{a,b\}$ a~$\{c,d,e\}$; nižší řez by dal víc shluků.}
\end{center}

\begin{poznamka}[Hustotní shlukování -- jen zmínka]
Existují i~\emph{hustotní} metody (DBSCAN): shluk $=$ souvislá oblast s~dostatečnou hustotou bodů, určenou parametry $\varepsilon$ (poloměr) a~\texttt{minPts}; body v~řídkých oblastech zůstanou jako šum. Umí najít shluky libovolného tvaru a~nepotřebuje předem $K$, ale v~požadavcích S3 jmenováno není (uvádíme jednou větou pro úplnost).
\end{poznamka}

% =====================================================================
\subsection{Redukce dimenze: analýza hlavních komponent (PCA)}\label{ssec:pca}

\textbf{Prokletí dimenzionality} (sekce~\ref{sec:preuceni}) i~potřeba vizualizace motivují \emph{redukci dimenze}: nahradit $D$ příznaků menším počtem $M<D$ nových příznaků, které zachovají co nejvíc informace. \textbf{PCA} (analýza hlavních komponent) je lineární metoda, která nové osy volí ve směrech \emph{největšího rozptylu} dat.

\begin{definice}[PCA: hlavní komponenty]
Mějme matici dat $\mathbf{X}\in\mathbb{R}^{N\times D}$ (příklad na řádku). PCA:
\begin{enumerate}
  \item \textbf{Centrování:} odečti průměr, $\tilde{\mathbf{X}}=\mathbf{X}-\bar{\mathbf{x}}$ (každý sloupec má střední hodnotu $0$).
  \item \textbf{Kovarianční matice:} $\mathbf{S}=\frac1N\tilde{\mathbf{X}}^{\!\top}\tilde{\mathbf{X}}\in\mathbb{R}^{D\times D}$ (symetrická, pozitivně semidefinitní).
  \item \textbf{Spektrální rozklad:} najdi vlastní čísla $\lambda_1\ge\lambda_2\ge\dots\ge\lambda_D\ge0$ a~ortonormální vlastní vektory $\mathbf{v}_1,\dots,\mathbf{v}_D$ matice $\mathbf{S}$. \emph{Hlavní komponenty} jsou vlastní vektory $\mathbf{v}_k$, seřazené sestupně podle $\lambda_k$.
  \item \textbf{Projekce:} nové souřadnice bodu jsou jeho průměty na prvních $M$ komponent, $\mathbf{z}_i=(\mathbf{v}_1^{\!\top}\tilde{\mathbf{x}}_i,\dots,\mathbf{v}_M^{\!\top}\tilde{\mathbf{x}}_i)\in\mathbb{R}^M$.
\end{enumerate}
\end{definice}

\begin{intuice}
Vlastní číslo $\lambda_k$ je přesně \emph{rozptyl dat ve směru} $\mathbf{v}_k$ (kovarianční matice měří rozptyl, její vlastní směry jsou hlavní osy „mraku`` dat). První hlavní komponenta $\mathbf{v}_1$ míří tam, kde jsou data nejvíc rozprostřená; $\mathbf{v}_2$ je na ni kolmá s~druhým největším rozptylem atd. Zahodit komponenty s~malým $\lambda$ znamená zahodit směry, v~nichž se data skoro nemění -- ztratíme nejmíň informace (v~smyslu rozptylu). \emph{Kolik komponent} vzít? Podle \emph{podílu zachyceného rozptylu} $\frac{\lambda_1+\dots+\lambda_M}{\lambda_1+\dots+\lambda_D}$ (např.\ tolik $M$, aby zachoval $95\,\%$), nebo z~„lomu`` v~grafu seřazených $\lambda$ (scree plot).
\end{intuice}

\begin{poznamka}[Souvislost se SVD a~M2]
PCA je vlastní rozklad symetrické pozitivně semidefinitní matice $\mathbf{S}=\mathbf{V}\Lambda\mathbf{V}^{\!\top}$ (M2: taková matice má reálná nezáporná vlastní čísla a~ortonormální vlastní vektory). Ekvivalentně lze hlavní komponenty získat z~\emph{singulárního rozkladu} $\tilde{\mathbf{X}}=\mathbf{U}\Sigma\mathbf{V}^{\!\top}$ (sloupce $\mathbf{V}$ jsou tytéž vlastní vektory, $\lambda_k=\sigma_k^2/N$); v~praxi se počítá přes SVD, numericky stabilnější. Pro jen pár prvních komponent stačí \emph{power iteration} (opakované násobení $\mathbf{S}$ s~normalizací konverguje k~dominantnímu vlastnímu vektoru). Detail SVD a~Eckart-Youngovu větu (nejlepší nízkohodnotová aproximace) zde nerozvádíme -- jen zmiňujeme.
\end{poznamka}

\begin{center}
\begin{tikzpicture}
\begin{axis}[width=9.8cm,height=7.4cm,axis equal image,
  xlabel={$x_1$},ylabel={$x_2$},
  xlabel style={at={(axis description cs:1.0,0.5)},anchor=north,yshift=-2pt,font=\scriptsize},
  ylabel style={at={(axis description cs:0.5,1.0)},anchor=east,xshift=-2pt,font=\scriptsize},
  xmin=-4.4,xmax=4.4,ymin=-3.4,ymax=3.4,
  xtick={-3,0,3},ytick={-2,0,2},tick label style={font=\scriptsize},
  axis lines=middle,
  every axis plot/.append style={thick},clip=false]
  \addplot[only marks,mark=*,mark size=1.1pt,blue!55!black,opacity=0.7]
    table[x=x,y=y]{fig/s10_pca_cloud.dat};
  % PC1 (dlouha osa, max rozptyl) a PC2 (kratsi, kolma) -- z gen_s10_pca.py
  \addplot[->,>=Stealth,very thick,red!70!black] coordinates {(0,0) (2.779,1.463)};
  \addplot[->,>=Stealth,very thick,red!70!black] coordinates {(0,0) (-2.779,-1.463)};
  \addplot[->,>=Stealth,very thick,green!45!black] coordinates {(0,0) (-0.404,0.768)};
  \addplot[->,>=Stealth,very thick,green!45!black] coordinates {(0,0) (0.404,-0.768)};
  \node[font=\scriptsize\bfseries,red!70!black,anchor=south west] at (axis cs:2.95,1.55) {$\mathbf{v}_1$\,(PC1)};
  \node[font=\scriptsize\bfseries,green!45!black,anchor=south east] at (axis cs:-0.55,0.95) {$\mathbf{v}_2$\,(PC2)};
\end{axis}
\end{tikzpicture}
\obrpopis{PCA na protáhlém 2D mraku (centrovaná data). První hlavní komponenta $\mathbf{v}_1$ (červeně) míří směrem \emph{největšího rozptylu}, druhá $\mathbf{v}_2$ (zeleně) je na ni kolmá s~menším rozptylem. Zde $\mathbf{v}_1$ zachycuje přes $90\,\%$ rozptylu, takže projekcí na samotnou $\mathbf{v}_1$ ($M=1$) ztratíme jen málo informace. Komponenty jsou vlastní vektory kovarianční matice, seřazené dle vlastních čísel $=$ rozptylů.}
\end{center}

\subsubsection*{Řešený příklad: PCA a~vliv na klasifikaci}

\begin{priklad}[SZZ léto 2024 č.~29: směry komponent a~PCA před $k$-NN]
\szzotazka{léto 2024}{29}{PCA analýza}\par
\textbf{Zadání.} Data se dvěma atributy ze tří tříd (obrázek -- tři protáhlé, vzájemně oddělené mraky podél „antidiagonály``, tj.\ směru klesající přímky). (1)~Jak funguje PCA, jak spočítat hlavní komponenty? (2)~Bez počítání odhadněte směry obou hlavních komponent. (3)~$k$-NN po redukci PCA na 1D -- jaká accuracy (přesnost) oproti původním datům? Vysvětlete.

\smallskip
\textbf{Řešení.} (1)~Viz výklad výše: centruj data, sestav kovarianční matici, její vlastní vektory (seřazené dle vlastních čísel) jsou hlavní komponenty; projekcí na prvních $M$ z~nich snížíš dimenzi.

(2)~Rozptyl dat je \emph{největší podél antidiagonály} (mraky se táhnou a~zároveň jsou rozmístěny ve směru klesající přímky, přes celý graf). \textbf{1.~hlavní komponenta $\mathbf{v}_1$ tedy míří podél této antidiagonály} (směr $\approx(1,-1)$, klesající), \textbf{2.~komponenta $\mathbf{v}_2$ je na ni kolmá} (směr $\approx(1,1)$, stoupající). Obě se zakreslí jako šipky z~těžiště dat.

(3)~\textbf{Accuracy klesne} (na původních 2D datech je $k$-NN téměř bezchybný, po projekci na 1D se třídy $+$ a~$\bullet$ z~velké části překryjí). Vysvětlení: PCA je \emph{bez učitele} -- maximalizuje rozptyl, \textbf{ne} oddělitelnost tříd. Projekce na $\mathbf{v}_1$ (antidiagonála) sice zachová nejvíc rozptylu, ale \emph{stlačí dohromady} třídy, které se liší hlavně v~\emph{kolmém} (málorozptylovém) směru $\mathbf{v}_2$ -- jejich průměty na $\mathbf{v}_1$ se překryjí a~$k$-NN je splete. Redukce na 1D vždy zahodí jednu souřadnici, takže může smazat právě tu informaci, která třídy dělila. (Projekce na málorozptylovou $\mathbf{v}_2$ by tu nebyla lepší: v~ní se zase překryjí třídy $\times$ a~$\bullet$, jejichž součty souřadnic $x_1+x_2$ jsou podobné -- ukázka, že „víc rozptylu`` $\ne$ „lepší pro klasifikaci``.)
\begin{center}
\begin{tikzpicture}[scale=0.78]
  \draw[osa] (0,0) -- (5.5,0) node[right,font=\scriptsize]{atribut~1};
  \draw[osa] (0,0) -- (0,5.5) node[above,font=\scriptsize]{atribut~2};
  % tri protahle mraky podel antidiagonaly (schematicky)
  % trida X vlevo nahore
  \foreach \i in {0,...,5}{\pgfmathsetmacro\x{0.5+0.18*\i}\pgfmathsetmacro\y{4.7-0.18*\i}\node[font=\scriptsize] at (\x,\y) {$\times$};}
  % trida + uprostred vpravo
  \foreach \i in {0,...,5}{\pgfmathsetmacro\x{3.5+0.18*\i}\pgfmathsetmacro\y{2.2-0.18*\i}\node[font=\scriptsize] at (\x,\y) {$+$};}
  % trida o vpravo dole
  \foreach \i in {0,...,5}{\pgfmathsetmacro\x{3.7+0.18*\i}\pgfmathsetmacro\y{1.1-0.18*\i}\node[font=\scriptsize] at (\x,\y) {$\bullet$};}
  % tezisko (priblizne stred vsech)
  \coordinate (ctr) at (2.6,2.2);
  % PC1 podel antidiagonaly (klesajici), PC2 kolma (stoupajici)
  \draw[->,>=Stealth,very thick,red!70!black] (ctr) -- ++(1.7,-1.7) node[anchor=north west,font=\scriptsize]{$\mathbf{v}_1$};
  \draw[->,>=Stealth,very thick,red!70!black] (ctr) -- ++(-1.7,1.7);
  \draw[->,>=Stealth,very thick,green!45!black] (ctr) -- ++(1.0,1.0) node[anchor=south west,font=\scriptsize]{$\mathbf{v}_2$};
  \draw[->,>=Stealth,very thick,green!45!black] (ctr) -- ++(-1.0,-1.0);
\end{tikzpicture}
\obrpopis{Schéma dat úlohy léto~2024: tři protáhlé třídy ($\times$, $+$, $\bullet$) podél antidiagonály. Odhadnuté hlavní komponenty: $\mathbf{v}_1$ (červeně) podél směru největšího rozptylu (klesající antidiagonála), $\mathbf{v}_2$ (zeleně) kolmá. Projekce na samotnou $\mathbf{v}_1$ překryje třídy $+$ a~$\bullet$ (liší se kolmo), proto $k$-NN po redukci na 1D ztrácí accuracy.}
\end{center}
\end{priklad}

% =====================================================================
\subsection{Měkké shlukování: EM a~směs Gaussovců}\label{ssec:em}

\begin{poznamka}[Zařazení]
EM a~směs Gaussovců nejsou v~požadavcích S3 jmenovány přímo, ale jde o~\emph{pravděpodobnostní zobecnění $k$-means} (měkké místo tvrdého přiřazení) a~okruh má na ně reálnou SZZ otázku; uvádíme je proto stručně a~vyřešíme otázku. Obecný princip EM (E-krok/M-krok pro modely se skrytými proměnnými) je vyložen v~S1.
\end{poznamka}

$k$-means přiřazuje bod \emph{tvrdě} k~jedinému shluku. \textbf{Směs Gaussovců} (Gaussian mixture) modeluje data jako směs $K$ normálních rozdělení a~přiřazuje \emph{měkce}: každý bod patří do shluku $k$ s~pravděpodobností (zodpovědností) $r_{i,k}\in[0,1]$, $\sum_k r_{i,k}=1$. Navíc dovolí shlukům \emph{různou velikost i~elipsovitý tvar} (vlastní kovarianční matici), kdežto $k$-means implicitně předpokládá kulové shluky podobné velikosti. Parametry (váhy, středy, kovariance) se učí \textbf{EM algoritmem}:
\begin{itemize}
  \item \textbf{E-krok:} při současných parametrech spočti zodpovědnosti $r_{i,k}=P(\text{shluk }k\mid\mathbf{x}_i)$ (Bayesova věta).
  \item \textbf{M-krok:} přepočti parametry jako \emph{vážené} (zodpovědnostmi) odhady -- vážené průměry a~kovariance.
\end{itemize}
$k$-means je limitní případ EM se sférickými Gaussovci nekonečně malého rozptylu, kde zodpovědnosti zkolabují na $0/1$ (tvrdé přiřazení).

\begin{priklad}[SZZ léto 2025 č.~32: EM na bayesovské síti]
\szzotazka{léto 2025}{32}{EM algoritmus}\par
\textbf{Zadání.} Bayesovská síť: skrytá $\mathit{Cluster}\in\{1,2\}$ je rodičem $\mathit{Color}\in\{\text{blue},\text{green}\}$ a~$\mathit{Holes}\in\{\text{yes},\text{no}\}$ (mezi $\mathit{Color}$ a~$\mathit{Holes}$ hrana není). (1)~Popište EM. (2)~Určete parametry, inicializujte z~$\{0{,}4,0{,}6\}$. (3)~Proveďte jednu EM aktualizaci jednoho parametru na datech (řádky $\mathit{Color},\mathit{Holes}$, $\mathit{Cluster}$ chybí).

\smallskip
\textbf{Řešení.} (1)~EM odhaduje parametry modelu se \emph{skrytou} proměnnou (zde $\mathit{Cluster}$) střídáním E-kroku (spočti očekávané hodnoty / zodpovědnosti skryté proměnné při daných parametrech) a~M-kroku (přepočti parametry maximalizací věrohodnosti na „doplněných`` datech), dokud parametry nezkonvergují.

(2)~Model má 5~nezávislých parametrů; inicializujeme z~$\{0{,}4,0{,}6\}$ tak, aby podmínky pro $\mathit{Cluster}{=}1$ a~$\mathit{Cluster}{=}2$ nebyly stejné:
\[
\begin{array}{llc}
\theta=P(\mathit{Cluster}{=}1)=0{,}6, &
\theta_{B1}=P(\text{blue}\mid 1)=0{,}6, & \theta_{B2}=P(\text{blue}\mid 2)=0{,}4,\\
& \theta_{H1}=P(\text{yes}\mid 1)=0{,}6, & \theta_{H2}=P(\text{yes}\mid 2)=0{,}4.
\end{array}
\]
(3)~Aktualizujeme nejjednodušší parametr $\theta=P(\mathit{Cluster}{=}1)$. \emph{E-krok} -- pro každý řádek spočti zodpovědnost $P(\mathit{Cluster}{=}1\mid\text{data})$ Bayesovou větou (faktorizace dle sítě: $P(C,\text{col},\text{hol})=P(C)P(\text{col}\mid C)P(\text{hol}\mid C)$):
\[
P(C{=}1\mid \text{col},\text{hol})=\frac{\theta\,P(\text{col}\mid1)P(\text{hol}\mid1)}{\theta\,P(\text{col}\mid1)P(\text{hol}\mid1)+(1-\theta)\,P(\text{col}\mid2)P(\text{hol}\mid2)} .
\]
Po řádcích (čitatel $=\theta\,P(\text{col}\mid1)P(\text{hol}\mid1)$, jmenovatel $+\,(1-\theta)\,P(\text{col}\mid2)P(\text{hol}\mid2)$) -- blue,no: $0{,}6\cdot0{,}6\cdot0{,}4=0{,}144$ vs.\ $0{,}4\cdot0{,}4\cdot0{,}6=0{,}096$; blue,yes: $0{,}6\cdot0{,}6\cdot0{,}6=0{,}216$ vs.\ $0{,}4\cdot0{,}4\cdot0{,}4=0{,}064$; green,yes: $0{,}6\cdot0{,}4\cdot0{,}6=0{,}144$ vs.\ $0{,}4\cdot0{,}6\cdot0{,}4=0{,}096$:
\[
\begin{array}{c|c|c}
\mathit{Color} & \mathit{Holes} & P(C{=}1\mid\text{data})\\\hline
\text{blue} & \text{no} & \dfrac{0{,}144}{0{,}144+0{,}096}=0{,}6\\[4pt]
\text{blue} & \text{yes} & \dfrac{0{,}216}{0{,}216+0{,}064}\approx0{,}77\\[4pt]
\text{blue} & \text{no} & 0{,}6\\[2pt]
\text{green} & \text{yes} & \dfrac{0{,}144}{0{,}144+0{,}096}=0{,}6\\[4pt]
\text{green} & \text{yes} & 0{,}6
\end{array}
\]
\emph{M-krok} -- nové $\theta$ je průměr zodpovědností (očekávaný podíl příkladů ve shluku~1):
\[
  \theta^{\text{nové}}=\frac{0{,}6+0{,}77+0{,}6+0{,}6+0{,}6}{5}=\frac{3{,}17}{5}=\boxed{0{,}634}.
\]
(Číslo přesně: $\tfrac{4\cdot0{,}6+0{,}7714}{5}=0{,}6343$.) Ostatní parametry by se aktualizovaly analogicky -- váženými četnostmi $\sum_j r_{j}\,[\dots]$ místo tvrdých počtů.
\end{priklad}

% =====================================================================
\subsection*{Shrnutí sekce}

\emph{Shlukování} ($k$-means) minimalizuje $J=\sum_i\sum_k z_{i,k}\|\mathbf{x}_i-\mu_k\|^2$ střídáním přiřazení (bod $\to$ nejbližší střed) a~přepočtu (střed $=$ průměr shluku); konverguje k~\emph{lokálnímu} minimu (oba kroky $J$ nezvyšují, přiřazení je konečně mnoho), proto se spouští vícekrát / s~k-means++. \emph{Hierarchické} (aglomerativní) shlukování buduje dendrogram slučováním nejbližších shluků (linkage single/complete/average/Ward); počet shluků určí řez dendrogramu. \emph{PCA} centruje data, vezme vlastní vektory kovarianční matice (hlavní komponenty seřazené dle vlastních čísel $=$ rozptylů) a~projektuje na prvních $M$ -- pozor, maximalizuje rozptyl, ne oddělitelnost tříd. \emph{EM / směs Gaussovců} je měkké, pravděpodobnostní zobecnění $k$-means.


% ============================================================
\section*{Shrnutí}
\addcontentsline{toc}{section}{Shrnutí}
\begin{poznamka}
\textbf{Klíčové vlastnosti a~vzorce (přehled).} Stručný výčet klíčových definic, vět a vzorců okruhu S3. Anglická terminologie tam, kde čeština není běžná. Zkratky: TP/TN/FP/FN~= true/false positive/negative; $D$~= počet příznaků; $N$~= počet příkladů; $\mathbf{w}$~= váhy; $\alpha$~= krok SGD; $\lambda$~= regularizační síla; $\hat y$~= predikce modelu.

\smallskip
\emph{Evaluace klasifikátoru.}
\begin{itemize}
  \item \textbf{Matice záměn:} řádky $=$ skutečná třída, sloupce $=$ predikovaná; TP (správně pozitivní), TN (správně negativní), FP (falešně pozitivní), FN (falešně negativní).
  \item $\mathrm{accuracy}=\dfrac{TP+TN}{TP+TN+FP+FN}$. Pro nevyvážené třídy \textbf{nevhodná}; používej precision/recall/F1 nebo AUC.
  \item $\mathrm{precision}=\dfrac{TP}{TP+FP}$;\quad $\mathrm{recall}=\mathrm{sensitivity}=\mathrm{TPR}=\dfrac{TP}{TP+FN}$;\quad $F_1=\dfrac{2\,TP}{2\,TP+FP+FN}=\dfrac{2\cdot\text{prec}\cdot\text{recall}}{\text{prec}+\text{recall}}$.
  \item $\mathrm{specificity}=\mathrm{TNR}=\dfrac{TN}{TN+FP}$;\quad $\mathrm{FPR}=1-\mathrm{specificity}$. \textbf{ROC:} TPR vs.\ FPR při plynulém posouvání prahu; \textbf{AUC}~= plocha pod ROC ($0{,}5$~= náhoda, $1{,}0$~= perfektní).
  \item \textbf{Regresní metriky:} $\mathrm{MSE}=\tfrac1N\sum_i(y_i-t_i)^2$; RMSE $=\sqrt{\mathrm{MSE}}$; MAE $=\tfrac1N\sum_i|y_i-t_i|$.
  \item \textbf{Křížová validace} ($k$-fold): data dělí na $k$ foldů; každý slouží jednou jako validační, ostatní jako trénovací; průměr $k$ chyb $=$ odhad generalizační chyby. Leave-one-out: $k=N$.
  \item \textbf{MLE:} $\mathbf{w}_\mathrm{MLE}=\argmin_{\mathbf{w}}\sum_i\bigl(-\log p(t_i\mid\mathbf{x}_i;\mathbf{w})\bigr)$; pro normální šum $\Rightarrow$ MSE, pro binární klasifikaci $\Rightarrow$ cross-entropy.
  \item \textbf{Naivní Bayes:} $\hat c=\argmax_c P(c)\prod_j P(x_j\mid c)$; předpoklad podmíněné nezávislosti příznaků dané třídou; $P(c)$ a~$P(x_j\mid c)$ odhadujeme z~dat četnostmi.
\end{itemize}

\smallskip
\emph{Overfitting, regularizace a~prokletí dimenzionality.}
\begin{itemize}
  \item \textbf{Overfitting:} model se naučí šum; trénovací chyba malá, generalizační velká. \textbf{Underfitting:} model slabý; obě chyby velké. \emph{U-křivka:} generalizační chyba má minimum pro optimální kapacitu.
  \item \textbf{L2 regularizace} (ridge, Tikhonov): $E(\mathbf{w})=E_0(\mathbf{w})+\frac\lambda2\|\mathbf{w}\|^2$; penalizuje velké váhy, dává hladké řešení; analyticky pro lin.\ regresi: $\mathbf{w}=(\mathbf{X}^\top\mathbf{X}+\lambda\mathbf{I})^{-1}\mathbf{X}^\top\mathbf{t}$.
  \item \textbf{L1 regularizace} (LASSO): $E(\mathbf{w})=E_0(\mathbf{w})+\lambda\|\mathbf{w}\|_1$; tlačí váhy přesně na nulu~$\Rightarrow$ řídká řešení (feature selection). Geometricky: L1 koule (v~2D kosočtverec) má rohy na osách, L2 koule je hladká~-- proto L1 řešení typicky padne do rohu (nulová souřadnice).
  \item \textbf{Early stopping:} sleduj validační chybu; zastav trénování ve chvíli, kdy začne na validaci růst (model se začíná přeučovat). Hyperparametr je počet trpělivostních epoch (\textit{patience}).
  \item \textbf{Prokletí dimenzionality:} objem prostoru roste exponenciálně s~$D$ -- pokrýt $[0,1]^D$ mřížkou s~krokem $r$ vyžaduje $N\propto r^{-D}$ příkladů; objem se koncentruje u~povrchu; vzdálenosti ztrácejí kontrast ($k$-NN degraduje); nutná regularizace nebo redukce dimenze.
\end{itemize}

\smallskip
\emph{$k$-NN ($k$ nejbližších sousedů).}
\begin{itemize}
  \item \textbf{Líný/neparametrický:} trénování $=$ uložení dat; inference $=$ najdi $k$ nejbližších sousedů. \emph{Klasifikace} $=$ nejčastější třída, \emph{regrese} $=$ průměr cílů $k$ sousedů.
  \item Metrika: $L^p$ ($L^1$ manhattanská, $L^2$ euklidovská). Nutná \textbf{normalizace příznaků} (různé škály deformují vzdálenosti). Volba $k$ křížovou validací: malé $k$ $\Rightarrow$ overfitting, velké $k$ $\Rightarrow$ underfitting.
\end{itemize}

\smallskip
\emph{Lineární regrese.}
\begin{itemize}
  \item Model: $y(\mathbf{x})=\mathbf{w}^\top\mathbf{x}+b$; minimalizujeme $\mathrm{MSE}=\tfrac1N\|\mathbf{X}\mathbf{w}-\mathbf{t}\|^2$ (ekvivalentně součet čtverců).
  \item \textbf{Normální rovnice} (analytické řešení): $\mathbf{w}=(\mathbf{X}^\top\mathbf{X})^{-1}\mathbf{X}^\top\mathbf{t}$; s~L2: $(\mathbf{X}^\top\mathbf{X}+\lambda\mathbf{I})^{-1}\mathbf{X}^\top\mathbf{t}$. Přesné, ale pomalé ($\mathcal O(ND^2+D^3)$), neumí early stopping.
  \item \textbf{SGD krok} (jeden příklad $(\mathbf{x},t)$): $\mathbf{w}\leftarrow\mathbf{w}-\alpha\bigl[(\mathbf{w}^\top\mathbf{x}-t)\,\mathbf{x}+\lambda\mathbf{w}\bigr]$; bias bez L2: $b\leftarrow b-\alpha(\hat y-t)$. Škálovatelné, umí early stopping.
\end{itemize}

\smallskip
\emph{Logistická regrese.}
\begin{itemize}
  \item \textbf{Binární:} $\hat y=\sigma(\mathbf{w}^\top\mathbf{x}+b)=P(C_1\mid\mathbf{x})$; $\sigma(z)=\frac{1}{1+e^{-z}}\in(0,1)$; prahování na $0{,}5$; $P(C_0)=1-\hat y$.
  \item \textbf{Cross-entropy ztráta} (NLL Bernoulliho): $E=-\sum_i\bigl[t_i\log\hat y_i+(1-t_i)\log(1-\hat y_i)\bigr]$.
  \item \textbf{Gradient/SGD krok:} $\nabla_{\mathbf{w}}E_i=(\hat y_i-t_i)\,\mathbf{x}_i$; $\mathbf{w}\leftarrow\mathbf{w}-\alpha(\hat y-t)\mathbf{x}$ (tvar identický s~lin.\ regresí, liší se jen $\hat y$).
  \item \textbf{Multiclass (softmax):} $K$ tříd, $K$ vektorů vah $\mathbf{w}_k$; $p_k=\dfrac{e^{\mathbf{w}_k^\top\mathbf{x}}}{\sum_j e^{\mathbf{w}_j^\top\mathbf{x}}}$; ztráta~= categorical cross-entropy $-\sum_i\log p_{t_i}(\mathbf{x}_i)$.
\end{itemize}

\smallskip
\emph{Rozhodovací stromy.}
\begin{itemize}
  \item \textbf{Stavba:} hladově shora dolů (greedy, rozděl a~panuj); v~každém uzlu vyber řez $(j^*,v^*)=\argmin_{j,v}\bigl[c(\mathcal T_L)+c(\mathcal T_R)\bigr]$; rekurze do zastavovací podmínky (max.\ hloubka, min.\ příkladů). Predikce listu: nejčastější třída (klas.) nebo průměr cílů (reg.).
  \item \textbf{Entropie:} $H(S)=-\sum_k p_k\log_2 p_k$ (bity; $0\log0:=0$; čistý uzel $\Rightarrow H=0$).
  \item \textbf{Informační zisk:} $\mathrm{IG}=H(\mathcal T)-\dfrac{|\mathcal T_L|}{|\mathcal T|}H(\mathcal T_L)-\dfrac{|\mathcal T_R|}{|\mathcal T|}H(\mathcal T_R)$; volíme řez s~maximálním $\mathrm{IG}$.
  \item \textbf{Gini index:} $\mathrm{Gini}(S)=\sum_k p_k(1-p_k)=1-\sum_k p_k^2$; rychlejší výpočet, podobné výsledky jako entropie; pro regresní stromy používáme vážené MSE (rozptyl cílů v~uzlu).
  \item \textbf{Prořezávání} (pruning): neomezený strom overfittuje. \emph{Post-pruning} (cost-complexity): $C_\alpha(T)=\text{chyba}(T)+\alpha\cdot|\text{listy}(T)|$; ladí se $\alpha$ na validaci. \emph{Pre-pruning:} max.\ hloubka, min.\ příkladů v~uzlu.
\end{itemize}

\smallskip
\emph{SVM (metoda podpůrných vektorů).}
\begin{itemize}
  \item \textbf{Hard-margin} (lineárně separabilní): $\min_{\mathbf{w},b}\,\tfrac12\|\mathbf{w}\|^2$ za $t_i(\mathbf{w}^\top\mathbf{x}_i+b)\ge1$; cílové hodnoty $t_i\in\{-1,+1\}$; margin $=\tfrac{1}{\|\mathbf{w}\|}$, šíře pásu $=\tfrac{2}{\|\mathbf{w}\|}$. \textbf{Podpůrné vektory:} trénovací body na hranici pásu ($t_i\cdot y(\mathbf{x}_i)=1$).
  \item \textbf{Soft-margin} (neseparabilní): $\min\,\tfrac12\|\mathbf{w}\|^2+C\sum_i\xi_i$ za $t_i(\mathbf{w}^\top\mathbf{x}_i+b)\ge1-\xi_i$, $\xi_i\ge0$. $C$~velké $\Rightarrow$ menší margin, méně klasifikačních chyb; $C$~malé $\Rightarrow$ větší margin, více chyb.
  \item \textbf{Kernel trick:} nahraď $\mathbf{x}_i^\top\mathbf{x}_j$ jádrem $K(\mathbf{x}_i,\mathbf{x}_j)=\varphi(\mathbf{x}_i)^\top\varphi(\mathbf{x}_j)$ -- výpočet $K$ bez explicitní $\varphi$. \textbf{RBF:} $K(\mathbf{x},\mathbf{x}')=e^{-\gamma\|\mathbf{x}-\mathbf{x}'\|^2}$. \textbf{Polynomiální:} $(1+\mathbf{x}^\top\mathbf{x}')^d$.
  \item \textbf{Multiclass:} one-vs-rest -- $K$ binárních SVM, predikce $=$ třída s~nejvyšším $y_k(\mathbf{x})$; one-vs-one -- $\binom{K}{2}$ SVM pro dvojice tříd, predikce hlasováním (\texttt{libsvm}).
\end{itemize}

\smallskip
\emph{Kombinace více modelů (ensemble).}
\begin{itemize}
  \item \textbf{Bagging} (bootstrap aggregating): $M$ kopií modelu, každá na vlastním bootstrapovém vzorku ($N$ příkladů s~opakováním); predikce průměrem (reg.) nebo hlasováním (klas.). Snižuje \emph{rozptyl}; out-of-bag příklady ($\approx37\,\%$) $\Rightarrow$ zdarma validace.
  \item \textbf{Boosting:} sekvenční -- každý model se zaměří na příklady, které předchozí chybovaly. \textbf{AdaBoost:} hlas slabého klasifikátoru $\alpha_m=\tfrac12\ln\dfrac{1-\varepsilon_m}{\varepsilon_m}$; špatně klasifikovaným příkladům zvýší váhy $w_i$; finální $\hat y=\sign\bigl(\sum_m\alpha_m h_m(\mathbf{x})\bigr)$.
  \item \textbf{Random forest:} bagging rozhodovacích stromů + v~každém uzlu nejlepší řez z~náhodné podmnožiny příznaků ($\approx\!\sqrt{D}$, typicky u~klasifikace); stromy se nechávají přerůst; snižuje korelaci chyb stromů.
  \item Ensemble zlepší pouze při \textbf{nekorelovaných chybách} modelů; jsou-li chyby totožné, kombinace nepomůže.
\end{itemize}

\smallskip
\emph{Statistické testy.}
\begin{itemize}
  \item \textbf{Jednovýběrový t-test:} $H_0\colon\mu=\mu_0$; testová statistika $t=\dfrac{\bar X-\mu_0}{s/\sqrt{n}}$ ($s$ = výběrová směrodatná odchylka); pod $H_0$ platí $t\sim t_{n-1}$. Zamítáme $H_0$, padne-li $|t|$ do kritického oboru (tj.\ $|t|>t_{n-1,1-\alpha/2}$, nebo $p\text{-value}<\alpha$).
  \item \textbf{Dvouvýběrový t-test:} $H_0\colon\mu_1=\mu_2$; při shodných rozptylech $t=\dfrac{\bar X-\bar Y}{s_p\sqrt{1/n_1+1/n_2}}$, df $= n_1+n_2-2$, kde $s_p^2=\frac{(n_1-1)s_1^2+(n_2-1)s_2^2}{n_1+n_2-2}$; \emph{Welchova} varianta (nestejné rozptyly): jiný jmenovatel a~df.
  \item \textbf{$\chi^2$ test dobré shody:} $H_0$: pozorované četnosti odpovídají očekávaným; $\chi^2=\sum_{i=1}^k\dfrac{(O_i-E_i)^2}{E_i}$; pod $H_0$ platí $\chi^2\sim\chi^2_{k-1}$ ($k$ kategorií). Zamítáme, je-li $\chi^2>$ kritická hodnota (nebo $p<\alpha$).
\end{itemize}

\smallskip
\emph{Učení bez učitele: shlukování a~redukce dimenze.}
\begin{itemize}
  \item \textbf{$k$-means} (Lloydův algoritmus): střídej \emph{přiřazení} (každý bod k~nejbližšímu centroidu) a~\emph{přepočet} (centroid~$=$ průměr shluku); kritérium $J=\sum_i\sum_k z_{i,k}\|\mathbf{x}_i-\mu_k\|^2$. Konverguje k~\emph{lokálnímu} minimu ($J$ v~žádném kroku neroste; konečně mnoho přiřazení $\Rightarrow$ zastaví se). Spouštět vícekrát, vzít nejmenší $J$.
  \item \textbf{k-means++:} první střed náhodně; každý další s~pravděpodobností $\propto$ čtverci vzdálenosti od nejbližšího dosud zvoleného -- lepší inicializace, $\mathcal O(\log K)$ aproximace optima v~očekávání.
  \item \textbf{Hierarchické shlukování} (aglomerativní, zdola nahoru): začni se $N$ singletony; opakovaně slučuj dvojici s~nejmenší vzdáleností; výsledkem je \textbf{dendrogram} -- vodorovný řez určuje počet shluků. Linkage: \emph{single} (min vzdálenost; riziko řetězení), \emph{complete} (max vzdálenost; kompaktní shluky), \emph{average}, \emph{Ward} (min nárůst vnitroshlukového rozptylu; výchozí ve \texttt{scikit-learn}).
  \item \textbf{PCA} (analýza hlavních komponent): (1)~centruj data $\tilde{\mathbf{X}}=\mathbf{X}-\bar{\mathbf{x}}$; (2)~kovarianční matice $\mathbf{S}=\tfrac1N\tilde{\mathbf{X}}^\top\tilde{\mathbf{X}}$; (3)~vlastní vektory $\mathbf{v}_1,\dots,\mathbf{v}_D$ seřazené dle $\lambda_1\ge\dots\ge\lambda_D\ge0$ jsou hlavní komponenty; (4)~projekce: $\mathbf{z}_i=(\mathbf{v}_1^\top\tilde{\mathbf{x}}_i,\dots,\mathbf{v}_M^\top\tilde{\mathbf{x}}_i)\in\mathbb{R}^M$. Vlastní číslo $\lambda_k$~= rozptyl dat ve~směru $\mathbf{v}_k$; zachovaný podíl $=\dfrac{\sum_{k=1}^M\lambda_k}{\sum_{k=1}^D\lambda_k}$. PCA je bez učitele -- maximalizuje rozptyl, \emph{ne} oddělitelnost tříd; po projekci může accuracy $k$-NN klesnout.
\end{itemize}
\end{poznamka}


% ============================================================
\section*{Zdroje}
\addcontentsline{toc}{section}{Zdroje}

{\raggedright
\begin{itemize}
  \item M.~Straka, J.~Libovický: slidy přednášek předmětu NPFL129 \emph{Úvod do strojového učení v~Pythonu} (Introduction to Machine Learning with Python), ÚFAL MFF UK, ročník 2024/25 (primární zdroj výkladu; téma SVM z~ročníku 2022/23, kdy bylo součástí osnovy).\\ \url{https://ufal.mff.cuni.cz/courses/npfl129}
  \item M.~Straka a~kol.: zdrojový repozitář kurzu NPFL129 (Markdown zdroj slidů a~generátory obrázků, včetně archivních ročníků).\\ \url{https://github.com/ufal/npfl129}
  \item R.~Šámal: \emph{NMAI059 Pravděpodobnost a statistika~1}, skripta k~přednášce, MFF UK, verze 14.~6.~2025 (rámec testování hypotéz, t-test a~$\chi^2$ test; podrobně v~okruhu M5).\\ \url{https://iuuk.mff.cuni.cz/~samal/vyuka/2425/PSt1/skripta.pdf}
  \item MFF UK: požadavky ke státní závěrečné zkoušce (Informatika, Bc.) a~zadání otázek z~minulých termínů.\\ \url{https://www.mff.cuni.cz/cs/studenti/bakalarske-studium/statni-zaverecne-zkousky/bakalarske-statni-zkousky-studijniho-programu-informatika}
\end{itemize}
\par}

\end{document}
